\chapter{安全水印与隐私}
\label{ch:21}

\paragraph*{本章阅读锚点}
本章的锚点是安全边界。原始 3DGS 默认模型可信可公开；安全研究关注水印、版权、隐私、攻击、防护和模型泄露风险。


随着 3DGS 资产（数字人、建筑模型、城市孪生）的商业化部署，知识产权保护和隐私安全成为不可忽视的现实问题。3DGS 的显式表示特性使其面临独特的安全威胁——场景参数可被直接提取和篡改，训练数据中的隐私信息可能通过渲染结果泄露。与此同时，3DGS 的可微渲染管线也为数字水印、对抗防御等安全技术提供了新的实现途径。

本章系统梳理了~47~篇相关论文，按~4~个子方向展开讨论：数字水印（第~\ref{sec:21_1}~节）、对抗攻击与防御（第~\ref{sec:21_2}~节）、隐写与信息隐藏（第~\ref{sec:21_3}~节）以及其他安全方法（第~\ref{sec:21_4}~节）。

\begin{figure}[htbp]
\centering
\resizebox{\textwidth}{!}{\begin{tikzpicture}[
  route/.style={rectangle, rounded corners=5pt, draw=#1, fill=#1!10,
    text width=3.2cm, minimum height=0.9cm, align=center, font=\small},
  method/.style={rectangle, rounded corners=3pt, draw=#1, fill=#1!6,
    font=\scriptsize, text width=3.0cm, align=center, minimum height=0.6cm},
  arr/.style={-{Stealth[length=4pt]}, techgray!60},
  >=Stealth
]

  \node[route=techgray, text width=4cm, font=\small\bfseries] (root) at (34.8, 5.6) {安全水印与隐私方法分类};
  \node[route=techblue] (b0) at (12.4, 3.1) {数字水印};
  \node[method=techblue] (b0_l0) at (2.0, 0.6) {基础水印框架\\{\tiny 3D-GSW, WATER-GS, GaussianMarker, GuardSplat}};
  \node[method=techblue] (b0_l1) at (6.8, 0.6) {高容量与可泛化\\{\tiny BitC-3DGS, GS-Marker, MarkSplatter}};
  \node[method=techblue] (b0_l2) at (12.4, 0.6) {生成水印与压缩鲁棒\\{\tiny GaussianSeal, CompMarkGS}};
  \node[method=techblue] (b0_l3) at (18.2, 0.6) {统一保护与鲁棒性\\{\tiny GuardMarkGS, RDSplat}};
  \node[route=techgreen] (b1) at (30.2, 3.1) {对抗攻击与防御};
  \node[method=techgreen] (b1_l0) at (23.3, 0.6) {重建质量攻击\\{\tiny CopyRNeRF, GaussianSteal}};
  \node[method=techgreen] (b1_l1) at (28.2, 0.6) {下游检测器攻击\\{\tiny Adv3D, BadGS}};
  \node[method=techgreen] (b1_l2) at (32.5, 0.6) {防御方法\\{\tiny DEGauss, WatermarkDefense}};
  \node[route=techorange] (b2) at (37.4, 3.1) {隐写与信息隐藏};
  \node[method=techorange] (b2_l0) at (37.4, 0.6) {耦合特征与对比学习\\{\tiny X-SG²S, StegaNeRF-GS}};
  \node[method=techorange] (b2_l1) at (42.7, 0.6) {频域加密隐写\\{\tiny FreqStego-GS, DualStream-GS}};
  \node[route=techpurple] (b3) at (47.5, 3.1) {其他安全};
  \node[method=techpurple] (b3_l0) at (47.5, 0.6) {篡改检测与验证\\{\tiny TamperDetect-GS}};
  \node[method=techpurple] (b3_l1) at (52.6, 0.6) {版权与隐私保护\\{\tiny CopyRight-GS, PrivacyGS}};
  \draw[arr] (root) -- (b0);
  \draw[arr] (root) -- (b1);
  \draw[arr] (root) -- (b2);
  \draw[arr] (root) -- (b3);
  \draw[arr] (b0) -- (b0_l0);
  \draw[arr] (b0) -- (b0_l1);
  \draw[arr] (b0) -- (b0_l2);
  \draw[arr] (b0) -- (b0_l3);
  \draw[arr] (b1) -- (b1_l0);
  \draw[arr] (b1) -- (b1_l1);
  \draw[arr] (b1) -- (b1_l2);
  \draw[arr] (b2) -- (b2_l0);
  \draw[arr] (b2) -- (b2_l1);
  \draw[arr] (b3) -- (b3_l0);
  \draw[arr] (b3) -- (b3_l1);
\end{tikzpicture}}
\caption{安全水印与隐私方法分类}
\label{fig:ch21_security_taxonomy}
\end{figure}

%=====================================================================
\section{数字水印}
\label{sec:21_1}

数字水印是~3DGS~知识产权保护的核心技术，旨在将版权信息嵌入~3DGS~资产中，使得从渲染图像或模型文件中可提取水印以证明所有权。3DGS~水印的独特挑战在于：渲染器位于~3D~编码器（高斯参数）和~2D~解码器（消息提取网络）之间，阻断了直接梯度流。水印信号必须先被编码到~3D~高斯参数中，经过~$\alpha$-混合~光栅化渲染为~2D~图像，再从~2D~图像中解码；这个~3D-2D-3D~的信息传递链路比~2D~图像水印的直接编码-解码要复杂得多。此外，高斯的显式结构（点云文件公开透明、每个高斯具有明确的~SH~系数、不透明度等参数）使得水印必须安全地嵌入球谐函数（SH）系数等参数中，同时防止恶意用户通过直接检查或修改模型文件来移除水印。

\subsection{基础水印框架}

本节讨论四条技术路线：（1）高容量与可泛化水印；（2）生成模型水印与压缩鲁棒；（3）统一保护与鲁棒性增强；（4）可解释水印与安全性分析。首先介绍~3DGS~水印的奠基性工作。这些工作建立了从"在~3D~参数中编码信息"到"从~2D~渲染图像中解码信息"这一基本技术链路，为后续的高容量、可泛化等改进奠定了基础。

3D-GSW（2024.09）~\cite{2409_13222} 首次提出~3DGS~鲁棒水印方法。它通过微调（fine-tuning）预训练~3DGS~的参数来嵌入水印，用频率引导致密化（FGD，Frequency-Guided Densification）替代原始~3DGS~的自适应密度控制，同时保持~3DGS~的渲染管线不变。该方法建立了"在~3D~高斯参数中编码二进制消息 → 从任意视角的~2D~渲染图像中解码消息"这一基本技术链路。

\paragraph{Baseline 做法与挑战}
在~3D-GSW~之前，水印方法主要针对隐式~NeRF（Neural Radiance Field，神经辐射场）~表示设计（如~StegaNeRF、WateRF），无法直接迁移到~3DGS~的显式表示。直接将~3DGS~所有高斯用于水印嵌入面临两个核心困难：（1）预训练~3DGS~包含大量冗余高斯，其中许多对渲染质量贡献微乎其微；在这些高斯中嵌入的水印信号在渲染图像中极为微弱，难以被鲁棒解码。（2）大幅修改已优化好的高斯参数会严重破坏渲染质量。此外，3DGS~场景通常包含数十万高斯，对全部参数施加水印微调的计算开销和存储开销都非常高。

\paragraph{整体流程}
3D-GSW~的技术路线分为四个阶段：（1）预训练一个与场景无关的消息解码器~$D_m$（基于~HiDDeN（Hide-and-Seek，一种端到端可训练的图像水印编码-解码架构）~架构），使其能从~2D~图像中提取固定长度的二进制消息；（2）对预训练~3DGS~执行频率引导致密化（FGD），移除冗余高斯并分裂高频区域的高斯；（3）基于~FGD~处理后的高斯构建梯度掩码；（4）在掩码约束下微调~3DGS~，使用小波子带损失保护高频渲染质量。直觉上，前两个阶段是"选好承载水印的容器"，第三个阶段是"控制修改幅度"，最后一个阶段是"实际嵌入信号"。

\paragraph{频率引导致密化（FGD）}
FGD~在微调前对预训练~3DGS~进行结构优化，分为两个阶段执行。

第一阶段：量化并移除低贡献高斯。受基于误差的致密化工作启发，3D-GSW~引入辅助损失函数来度量每个~3D~高斯对渲染质量的贡献。具体地，将所有高斯的颜色参数~$C$~替换为零向量~$C'$（仅在此度量过程中），渲染出~$I^{c'}_\theta$，并计算辅助损失：

\begin{equation}
  L^{aux}_\theta := \sum_{(x,y) \in \text{Pix}} \mathbb{E}_{\theta \sim E_\theta} \| I^{c'}_\theta - I^{gt}_\theta \|_1
\end{equation}

直觉上，这相当于"关闭"每个高斯的颜色，看渲染图像变差多少——变差越多，说明该高斯越重要。通过反向传播，贡献度向量~$V_\theta = \partial L^{aux} / \partial C'$~可以展开为：

\begin{equation}
  V_\theta = \sum_{(x,y) \in \text{Pix}} \sum_{i \in N_G} c_i \alpha_i T_i
\end{equation}

其中~$c_i$、$\alpha_i$~和~$T_i$~分别是第~$i$~个高斯的颜色、不透明度和透射率。$V_\theta \in \mathbb{R}^{N_G \times 3}$~直接反映每个高斯对渲染颜色的贡献权重。实验中，贡献度低于阈值~$V_\theta < 10^{-8}$~的高斯被移除——这些高斯携带水印信号的能力最弱，移除它们不仅不影响渲染质量，反而减少了存储和计算开销。

第二阶段：在高频区域分裂高斯以增强水印容量。将渲染图像划分为~$|P| = 16$~像素大小的~patch，对每个~patch~施加离散傅里叶变换（DFT，将信号从空间域转换到频率域的数学工具）：

\begin{equation}
  F[u,v] = \sum_{m=0}^{M-1} \sum_{n=0}^{N-1} f[m,n] \, e^{-j2\pi(\frac{um}{M} + \frac{vn}{N})}
\end{equation}

然后通过一个强调高频信号的掩码~$Q$~计算每个~patch~的高频强度：

\begin{equation}
  E = \sum_{u,v \in (Q)_{uv}} \|\mathcal{R}(F(P))\|_{U \times V}
\end{equation}

直觉上，这一步的目的是找到图像中纹理丰富的区域。选取高频强度最高的前~$K = 1\%$~的~patch，将其中的高斯分裂为更小的基元。分裂的动机在于：小高斯对渲染质量的贡献更精细，在高频区域分裂后能更好地承载水印信号，同时恢复第一阶段移除高斯可能损失的细节。值得注意的是，FGD~的两个阶段总体上减少了高斯数量（移除量大于分裂增加量），因此在提升水印鲁棒性的同时改善了渲染效率。

\paragraph{3D~高斯梯度掩码}
经过~FGD~处理后的~3DGS~（记为~$G^*_o$）已具备高质量渲染能力，因此微调过程中必须严格控制参数变化幅度，以避免破坏渲染质量。3D-GSW~为每个高斯的每个参数（颜色~$c$、不透明度~$\alpha$、旋转~$r$、尺度~$s$；位置~$\mu$~因梯度本就接近零而无需掩码）计算一个基于参数幅度的指数衰减掩码：

\begin{equation}
  w_i = 1 - e^{-\alpha |p_i|^\beta}, \quad z = \sum_{i=1}^{N_{G^*_o}} w_i
\end{equation}

其中~$p_i$~是第~$i$~个高斯的参数值，$\alpha > 0$~控制梯度操控强度。采用指数函数而非线性掩码，原因在于：3DGS~的某些参数（如不透明度和旋转）包含零值，线性掩码会导致除以零的问题；指数函数则可以平滑地处理零值，同时让参数幅度大的高斯获得更大的掩码系数（允许更大梯度），参数幅度小的高斯获得更小的掩码系数（进一步抑制变化）。在微调过程中，梯度被掩码为~$\nabla L_{total} \odot z$（$\odot$~为逐元素乘积），实验中~$\alpha = 4$~在比特准确率和渲染质量之间取得了最佳平衡。

\paragraph{小波子带损失}
由于~FGD~在分裂高斯时主要针对高频区域，微调过程中~3DGS~倾向于通过增大高斯体积来渲染高频细节，这可能导致高频信息丢失。3D-GSW~利用离散小波变换（DWT，Discrete Wavelet Transform）对~2~层分解后的小波子带施加约束。直觉上，水印信息被嵌入在~LL~（低频）子带中——解码器以~$LL_2$~子带作为输入——因此需要单独保护高频子带的渲染保真度。小波子带损失仅对高频子带（LH、HL、HH）计算~$L_1$~距离：

\begin{equation}
  L_w = \sum_{l} \sum_{S \in \{LH, HL, HH\}} L_1(S_l, S^{gt}_l)
\end{equation}

\paragraph{总损失与训练}
3D-GSW~的总损失为四项的加权和：

\begin{equation}
  L_{total} = \lambda_{rec} L_{rec} + \lambda_{lpips} L_{lpips} + \lambda_w L_w + \lambda_m L_m
\end{equation}

其中~$L_{rec} = L_1(I_w, I_o)$~为重建损失，$L_{lpips}$~为感知损失，$L_w$~为上述小波子带损失，$L_m$~为消息损失（解码消息~$M'$~与目标消息~$M$~之间的二元交叉熵，经~sigmoid~函数约束在~$[0,1]$~范围内）。实验中的权重设置为~$\lambda_{rec} = 1$、$\lambda_{lpips} = 0.2$、$\lambda_w = 0.3$、$\lambda_m = 0.4$。微调仅需~$2$--$10$~个~epoch（每个~epoch~遍历训练视角一次），即可在所有视角的渲染图像中稳定嵌入~$32$--$64$~比特的二进制水印消息。

在~3D-GSW~的基础上，后续工作从不同角度进一步完善了基础水印框架。

\paragraph{以~3DGS~生成本身为编码器的通用水印}

3D-GSW~建立了"微调~3DGS~参数 + 专用致密化 + 梯度掩码"的水印范式。WATER-GS~沿袭了"微调预训练~3DGS~+ 预训练解码器提取"的基本链路，但放弃了~3D-GSW~的三个专用组件——FGD、基于参数幅度的指数衰减梯度掩码和小波子带损失——转而采用更简洁的通用策略。

\paragraph{设计动机}
3D-GSW~的~FGD、梯度掩码和小波子带损失虽然有效，但设计上深度耦合于标准~3DGS~的渲染管线。当目标扩展到~3DGS~压缩框架（点云数量大幅减少）或~2D Gaussian Splatting（参数含义不同）时，这些组件难以直接迁移。WATER-GS~的核心动机是追求通用性（plug-and-play）：不做任何~3DGS~管线修改，仅引入外部解码器和训练时失真增强，使水印嵌入自然适配任意~3DGS~变体。

\paragraph{无编码器设计：3DGS~生成本身即编码器}
与~3D-GSW~类似，WATER-GS~也使用基于~HiDDeN~架构的预训练消息解码器~$\mathcal{D}$，在~COCO~数据集上以~MSE~图像损失和~BCE（Binary Cross-Entropy，二元交叉熵）~消息损失端到端训练编码器~$\mathcal{E}$~与解码器~$\mathcal{D}$，训练完成后丢弃编码器、保留解码器。但关键差异在于水印编码方式：3D-GSW~通过~FGD~和梯度掩码精心控制哪些高斯参数承载水印信号；WATER-GS~则受固定神经网络（FNN，Fixed Neural Networks）启发，直接将原始~3DGS~生成网络~$\mathcal{F}$~本身视为潜在水印编码器——不引入任何额外编码器网络，也不修改~3DGS~的渲染管线。具体地，给定预训练~3DGS~参数~$\Theta = \{\boldsymbol{\mu}, \boldsymbol{s}, \boldsymbol{r}, \alpha, \boldsymbol{h}\}$，在冻结解码器~$\mathcal{D}$~的条件下微调~$\Theta$~至~$\tilde{\Theta}$，使渲染图像~$\tilde{I}$~中携带水印消息~$\boldsymbol{m}$：$\tilde{\boldsymbol{m}} = \mathcal{D}(\tilde{I})$。创作者可通过自定义掩码~$\boldsymbol{M}$~灵活控制参与微调的参数子集（例如仅微调~SH~系数而冻结几何参数），梯度更新为~$\eta \cdot (\partial \mathcal{L}_{tot} / \partial \tilde{\Theta} \odot \boldsymbol{M})$。直觉上，这相当于不改变~3DGS~的"容器结构"，而是利用~3DGS~网络自身强大的特征融合能力来"自然学会"编码水印。

\paragraph{3D~失真层}
3D-GSW~的鲁棒性设计集中在~2D~域（小波子带损失保护高频渲染细节）。WATER-GS~则首次关注~3D~域的鲁棒性。在真实世界中，3DGS~点云文件在传输和处理过程中可能遭受各种扰动（如点云裁剪、高斯噪声污染、随机点丢失），此前的工作均未考虑这些~3D~层面的失真。为此，WATER-GS~在微调阶段引入四类~3D~失真层~$\mathcal{N}$：Identity（保持不变，确保无失真场景下的基本提取能力）、GN（对参数施加高斯噪声模糊，增强对参数微小变化的鲁棒性）、Dropout（以比例~$p$~随机移除部分点，模拟剪枝压缩）、Crop（移除连续区间的点，模拟部分数据丢失）。微调时，3DGS~参数先经失真层扰动为~$\hat{\Theta} = \mathcal{N}(\tilde{\Theta})$，再渲染并解码。这种训练时主动注入~3D~失真的策略使水印在点云遭受各种~3D~域攻击后仍可从渲染图像中可靠提取。实验表明，3D~失真层可使提取准确率提升最高~$20\%$（例如在~LLFF~数据集上，高斯噪声失真下的比特错误率从~$25.62\%$~降至~$4.86\%$）。

\paragraph{训练细节与损失函数}
WATER-GS~的微调损失比~3D-GSW~更简洁：渲染损失~$\mathcal{L}_{rgb} = (1-\beta) \cdot L_1(I_{gt}, I_{wa}) + \beta \cdot \mathcal{L}_{SSIM}(I_{gt}, I_{wa})$~（联合~$L_1$~和~SSIM，不含~LPIPS（Learned Perceptual Image Patch Similarity，学习感知图像块相似度）~感知损失和小波子带损失），消息损失~$\mathcal{L}_m = \text{BCE}(\boldsymbol{m}, \tilde{\boldsymbol{m}})$，总损失~$\mathcal{L}_{tot} = \mathcal{L}_{rgb} + \gamma \cdot \mathcal{L}_m$。微调过程~$10$K$\sim 30$K~次迭代，且保留了~3DGS~原有的自适应密度控制（分裂与合并），确保高斯分布适应水印引起的参数变化。在~LLFF、Mip-NeRF 360~和~Tanks and Temples~数据集上，WATER-GS~在无失真条件下的比特错误率低至~$3.26\%$~（LLFF），同时渲染质量损失可忽略不计。

\paragraph{通用性验证}
WATER-GS~的核心设计目标——不修改~3DGS~管线——使其可直接适配~3DGS~压缩框架（将通用解码器嵌入压缩管线即可应对压缩失真）和~2D Gaussian Splatting（参数含义不同但微调流程一致），无需针对每种变体重新设计专用组件。这与~3D-GSW~的专用~FGD~和梯度掩码形成了对比。

与~3D-GSW~通过~FGD~和梯度掩码来控制水印扰动幅度不同，GaussianMarker（2024.10）~\cite{2410_23718} 从根本性的新角度出发——利用贝叶斯不确定性来指导水印嵌入，而非在微调过程中通过工程化的掩码约束来限制扰动。3D-GSW~的~FGD~移除低贡献高斯并分裂高频区域高斯，然后对剩余高斯施加基于参数幅度的指数衰减掩码，本质上是一种"全局筛选+逐参数加权"的策略。GaussianMarker~则提出了一种更具理论基础的方案：直接估计~3DGS~模型参数的后验分布，量化每个参数对渲染结果的不确定性贡献，然后仅在高不确定性参数中添加扰动来承载水印信号。

\paragraph{不确定性估计的理论动机}
GaussianMarker~的核心洞察来自贝叶斯学习理论：对于已收敛的~3DGS~模型，其参数的后验分布中，方差较大的参数对渲染输出具有更大的不确定性。这些参数的微小扰动对渲染质量的影响更小，因此是承载水印信号的理想位置。形式化地，给定模型参数~$\theta$~和训练数据集~$\mathcal{D}$，预测分布通过对模型后验进行边际化得到：

\begin{equation}
  p(I|V, \mathcal{D}) = \int p(I|V, \theta) \, p(\theta|\mathcal{D}) \, d\theta = \mathbb{E}_{p(\theta|\mathcal{D})}[p(I|V, \theta)]
\end{equation}

其中~$I$~是视角~$V$~下的渲染图像。直觉上，在已收敛模型中，后验方差量化了参数的不确定性。方差越大的参数，其精确值对渲染结果越不关键，因此能容忍更大的扰动而不显著影响渲染质量。

\paragraph{基于~Laplace~近似的不确定性建模}
为实际计算~3DGS~参数的后验分布，GaussianMarker~采用~Laplace~近似——一种经典的后验近似方法，将后验近似为以~MAP（Maximum A Posteriori，最大后验）估计为均值的高斯分布，其协方差矩阵为观测~Fisher~信息的逆矩阵：

\begin{equation}
  p(\theta|\mathcal{D}) \approx \mathcal{N}(\theta^*, \Sigma), \quad \Sigma = -H^{-1}
\end{equation}

其中~$\theta^*$~是已收敛的~3DGS~参数（即~MAP~解），$H$~是~Hessian~矩阵。3DGS~参数的不确定性因此可通过~Hessian~矩阵来估计。假设渲染图像~RGB~通道的协方差为~1（遵循~FisherRF~的设定），Hessian~矩阵可简化为：

\begin{equation}
  H_i = \nabla_\theta f(V_i; \theta^*)^\top \nabla_\theta f(V_i; \theta^*)
\end{equation}

其中~$f(V_i; \theta^*)$~是在视角~$V_i$~下用收敛参数~$\theta^*$~渲染的图像。由于~Fisher~信息具有可加性，模型整体不确定性~$U$~可由所有训练视角的~Hessian~矩阵汇总得到：

\begin{equation}
  U = \sum_{i=1}^{N} H_i
\end{equation}

最终，模型不确定性覆盖所有高斯参数类型（位置~$\mu$、旋转~$R$、缩放~$S$、颜色~$c$、不透明度~$\alpha$）：$H = H_\mu + H_R + H_S + H_c + H_\alpha$。直觉上，这一计算过程的含义是：那些在所有训练视角下梯度变化都较大的参数，对渲染输出非常敏感（低不确定性）；而在多数视角下梯度都较小的参数，对渲染输出不敏感（高不确定性），因而是水印嵌入的安全区域。

\paragraph{基于不确定性的致密化策略}
与~3D-GSW~的~FGD~先移除低贡献高斯再分裂高频高斯不同，GaussianMarker~采用了一种"致密化而非修改"的策略——不修改任何原始高斯参数，而是仅在高不确定性高斯附近创建新的致密化高斯来承载水印扰动。具体地，从原始高斯集合~$G$~中，筛选不确定性高于阈值~$\theta_{unc}$~的高斯：

\begin{equation}
  \Delta = \{g(G_i) \mid G_i \in G, \, U_i > \theta_{unc}\}
\end{equation}

其中致密化函数~$g(\cdot)$~对第~$i$~个高斯~$G_i$~进行随机位置采样~$\mu_i' \sim \mathcal{N}(\mu_i, \sigma_i)$~（在原位置附近按尺度采样新位置），其余参数（颜色、不透明度、旋转等）从原高斯克隆。阈值~$\theta_{unc}$~默认设为所有高斯不确定性的均值。致密化后的高斯集合~$\Delta$~即为~GaussianMarker，与原始高斯~$G$~组合构成含水印的高斯~$\hat{G} = G \cup \Delta$。由于~GaussianMarker~的位置采样分布~$\mu_i' \sim \mathcal{N}(\mu_i, \sigma_i)$~仅产生细微的几何差异，且这些高斯位于高不确定性区域，它们对渲染质量的视觉影响极小，但能有效地将水印信号传递到渲染图像中。

\paragraph{双域消息解码}
与~3D-GSW~仅支持从~2D~渲染图像解码水印不同，GaussianMarker~实现了双域消息提取——版权信息可同时从~3D~高斯参数和~2D~渲染图像中可靠提取。这样，3DGS~资产所有者无论面对何种攻击场景（攻击者仅获取点云文件或仅获取渲染图像），均可验证所有权。

对于~2D~域，GaussianMarker~采用与~3D-GSW~相同的~HiDDeN~架构作为~2D~消息解码器~$D_{2D}$，先在图像上预训练编码器和解码器（训练时施加高斯噪声、随机旋转、裁剪和~JPEG（Joint Photographic Experts Group，联合图像专家组的图像压缩标准）~压缩等~2D~失真以增强鲁棒性），预训练完成后仅保留解码器用于后续优化。

对于~3D~域，GaussianMarker~利用~3DGS~的显式点云结构，采用~PointNet（一种直接消费点云的深度网络架构）~作为~3D~消息解码器~$D_{3D}$——将高斯均值~$\mu$~作为点位置，其他参数（颜色、不透明度、缩放、旋转）作为点特征，从随机采样的含水印高斯子集中解码版权消息。同时训练一个~3D~判别器~$D_{cls}$~（同样基于~PointNet）来区分原始高斯点和含水印高斯点，提供对抗性训练信号。训练时在~3D~域施加多种失真（高斯噪声、随机旋转、随机平移、随机裁剪），确保~3D~消息解码对点云扰动具有鲁棒性。

从~GaussianMarker~过渡到下一个方法，水印研究的焦点从载体选择策略转向解码器训练范式的革新。

\paragraph{两阶段优化}
GaussianMarker~的优化分为两个阶段，均保持原始高斯~$G$~完全不变，仅优化~GaussianMarker~$\Delta$~的参数。

第一阶段：从~2D~到~3D~的水印知识蒸馏。利用预训练的~2D~解码器~$D_{2D}$~作为教师网络，通过知识蒸馏将~2D~图像空间的水印能力迁移到~3D~高斯空间。具体地，冻结原始高斯~$G$，优化~GaussianMarker~$\Delta$，使含水印图像~$I_w$（由~$\hat{G} = G \cup \Delta$~渲染）中解码的消息~$\hat{M}_{2D} = D_{2D}(I_w)$~与目标消息~$M$~匹配。优化目标为~2D~水印损失：

\begin{equation}
  L_{2D} = \lambda_1 L_{msg} + \lambda_2 L_{rec}
\end{equation}

其中消息损失~$L_{msg}$~为原始消息~$M$~与解码消息~$\hat{M}_{2D}$~之间的二元交叉熵，重建损失~$L_{rec} = \|I_w - I_o\|_2^2$~确保含水印图像与原始图像的多视角一致性。

第二阶段：训练~3D~消息解码器。在~GaussianMarker~$\Delta$~优化完成后，训练~3D~消息解码器~$D_{3D}$~从含水印高斯中提取版权消息，同时训练~3D~判别器~$D_{cls}$~进行对抗优化。优化目标为~3D~水印损失：

\begin{equation}
  L_{3D} = \lambda_1' L_{msg}' + \lambda_2' L_{adv}
\end{equation}

其中~$L_{msg}'$~为~3D~域的消息~BCE~损失，$L_{adv}$~为区分原始点和含水印点的对抗损失。实验中~$\lambda_1 = 10.0$、$\lambda_2 = 1.0$、$\lambda_1' = 2.0$、$\lambda_2' = 1.0$。整个优化仅需~$1000$~步（Blender/LLFF）或~$2000$~步（Mip-NeRF 360），单张~V100 GPU~上~$20$~分钟内即可完成。

在~Blender、LLFF~和~Mip-NeRF 360~数据集上的实验表明，GaussianMarker~在消息解码精度和视图合成质量上均达到当时最优，且对~2D~失真（JPEG~压缩、噪声、裁剪、旋转）和~3D~失真（噪声、旋转、平移、裁剪）均表现出良好的鲁棒性。

GuardSplat（2024.11）~\cite{2411_19895} 针对~3DGS~水印在安全性、容量、不可见性和训练效率四方面的平衡挑战提出高效框架。传统方法面临的核心矛盾是：增强水印鲁棒性通常需要更大的参数扰动，但这会降低渲染保真度。与锚点~3D-GSW~直接微调全部高斯参数并用~HiDDeN~架构从零训练消息解码器不同，GuardSplat~在解码器训练范式、水印嵌入策略和鲁棒性增强三个维度上均做出了根本性改进。

\paragraph{CLIP~引导的消息解耦优化}
3D-GSW~使用基于~HiDDeN~架构的消息解码器，需要同时训练编码器和解码器，优化成本较高。GuardSplat~提出消息解耦优化（MDO，Message Decoupling Optimization）——利用~CLIP（Contrastive Language-Image Pre-training，对比语言-图像预训练）~的文本-视觉跨模态对齐能力，仅训练消息解码器而无需训练编码器，从根本上降低了优化开销。具体流程如下：给定二进制消息~$M \in \{0,1\}^L$，首先通过位到令牌映射（bit-to-token mapping）将其转换为~CLIP~文本令牌序列：

\begin{equation}
  T = \{t_S\} \cup \left\{\bigcup_{i=1}^{L}\varPhi(M_i, i)\right\} \cup \{t_E\}
\end{equation}

其中~$t_S = 49406$~和~$t_E = 49407$~分别为~CLIP~文本令牌的起始和终止标记，$\varPhi(\cdot,\cdot)$~将第~$i$~个比特均匀映射到~$[1, 49405]$~范围内的整数。$T$~经零填充至~77~的固定长度后，送入~CLIP~文本编码器~$\mathcal{E}_\mathcal{T}$~提取文本特征~$F_\mathcal{T} \in \mathbb{R}^{512}$，再由消息解码器~$\mathcal{D}_\mathcal{M}$（仅含~3~个全连接层的~MLP（Multilayer Perceptron，多层感知机））从中提取消息~$\hat{M} = \mathcal{D}_\mathcal{M}(\mathcal{E}_\mathcal{T}(T))$。解码器通过二元交叉熵损失~$\mathcal{L}_{\text{msg}}$~优化。直觉上，这一设计的关键洞察在于：CLIP~在~4~亿图文对上预训练获得的丰富表示，使得解码器无需编码器的辅助即可学会精确提取消息。解码器在~CLIP~文本特征上训练完成后，可直接应用于~CLIP~视觉特征进行消息提取（利用~CLIP~的文本-视觉对齐特性），实现了"在文本域训练、在视觉域提取"的解耦范式。训练仅需~5~分钟（单张~RTX 3090），显著快于~3D-GSW~的逐场景解码器训练方式。

\paragraph{SH~感知消息嵌入}
3D-GSW~通过梯度掩码控制对~3DGS~所有属性（颜色、不透明度、旋转、尺度）的微调幅度，但仍会改变原始~3D~结构。GuardSplat~的分析表明：3DGS~的~4~组属性中，位置~$\boldsymbol{\mu}$、协方差~$\boldsymbol{\Sigma}$~和不透明度~$\boldsymbol{\alpha}$~共同定义~3D~结构，而~SH~特征~$\boldsymbol{h}_i$~仅负责视角相关的颜色表示（如光泽和高光），且~SH~参数仅在场景的少数区域存在显著的视角依赖效果。因此，GuardSplat~的策略是冻结所有属性，仅为每个~3D~高斯创建一个可学习的~SH~偏移量~$\boldsymbol{h}^o_i \in \mathbb{R}^{48}$，将偏移量叠加到原始~SH~系数上作为含水印的~SH~特征：

\begin{equation}
  \boldsymbol{h}_i^{\text{wm}} = \boldsymbol{h}_i + \boldsymbol{h}^o_i
\end{equation}

然后送入~3DGS~光栅化器渲染含水印视角。为约束偏移幅度，引入偏移损失：

\begin{equation}
  \mathcal{L}_{\text{off}} = \frac{1}{N}\sum_{i=1}^{N}\|\boldsymbol{h}^o_i\|_2^2
\end{equation}

这一设计带来了双重优势：（1）完全保留原始~3D~结构——位置、协方差和不透明度均被冻结，水印信号仅存在于~SH~偏移量中；实验表明仅修改~SH~偏移量在~32~比特消息长度下~PSNR~达~39.40~dB，而修改全部属性（如~GaussianMarker~的做法）仅为~33.98~dB；（2）安全性——水印信号与~SH~颜色系数深度融合，恶意用户若要从模型文件中移除水印，必须同时破坏颜色表示，使水印不可分离。

\paragraph{反失真消息提取}
在消息提取阶段，利用预训练的~CLIP~视觉编码器~$\mathcal{E}_\mathcal{V}$~和消息解码器~$\mathcal{D}_\mathcal{M}$，从含水印渲染视角~$\hat{C}$~中提取消息：$\hat{M} = \mathcal{D}_\mathcal{M}(\mathcal{E}_\mathcal{V}(\hat{C}))$。由于~CLIP~本身的泛化能力，消息提取天然对高斯模糊和高斯噪声具有鲁棒性。然而对于旋转、JPEG~压缩等失真，GuardSplat~进一步引入可微失真层——在水印嵌入优化过程中，对渲染视角随机施加裁剪、缩放、旋转、亮度抖动和~JPEG~压缩等失真，使含水印~SH~特征通过优化学习对这些失真的抗性。与~3D-GSW~的小波子带损失（专注于保护渲染质量的高频保真度）不同，GuardSplat~的反失真层专注于增强水印信号本身在各种降质条件下的可提取性。

从基础框架到高容量水印，下一阶段的研究开始关注如何突破水印的信息容量瓶颈。

\paragraph{总损失与训练}
GuardSplat~的水印嵌入损失为：

\begin{equation}
  \mathcal{L} = \lambda_{\text{recon}}(\mathcal{L}_{\text{rgb}} + \mathcal{L}_{\text{lpips}}) + \lambda_{\text{msg}}\mathcal{L}_{\text{msg}} + \lambda_{\text{off}}\mathcal{L}_{\text{off}}
\end{equation}

其中~$\mathcal{L}_{\text{rgb}}$~和~$\mathcal{L}_{\text{lpips}}$~分别为~RGB~重建损失和感知损失（确保含水印视角与原始视角的视觉一致性），$\mathcal{L}_{\text{msg}}$~为消息提取的二元交叉熵，$\mathcal{L}_{\text{off}}$~为上述偏移约束。超参数设置为~$\lambda_{\text{recon}} = 1$、$\lambda_{\text{msg}} = 0.03$、$\lambda_{\text{off}} = 10$。训练分为两阶段：先用~CLIP~引导解耦优化训练消息解码器（约~5~分钟），再冻结解码器优化~SH~偏移量（约~10~分钟），总计约~15~分钟即可完成单个~3DGS~资产的水印嵌入。在~Blender~和~LLFF~数据集上，GuardSplat~在~16/32/48~比特消息长度下的比特准确率分别为~99.64\%/99.04\%/98.29\%，同时~PSNR~保持在~38.90 dB~以上，在安全性、容量、不可见性和效率四个维度上均超越~3D-GSW~等同期方法。

\subsubsection{高容量与可泛化水印}

本路线介绍提升水印信息容量和实现跨场景泛化嵌入的方法。随着~3DGS~资产管理的复杂化，高容量水印（嵌入更多信息比特）成为重要需求。与锚点~3D-GSW~在低层信号空间微调高斯参数的路线不同，语义水印方法（如~GuardSplat）将嵌入空间从像素级扰动提升到~CLIP~语义对齐层面，显著提升了鲁棒性。但受限于~CLIP~固定的~77-token~上下文长度——由于超出此长度的~token~没有可用的学习位置嵌入，基于位到令牌（bit-to-token）编码的水印方法最多只能嵌入~77~比特消息。BitC-3DGS（2026.05）~\cite{2605_29583} 针对这一~token~数量瓶颈提出位压缩框架，在不改变~CLIP~编码器结构的前提下将每~token~承载的信息量从~1~比特提升到~$n$~比特。

\paragraph{位压缩令牌化}
BitC-3DGS~的核心设计是位压缩令牌化方案：将长度为~$L$~的二值消息分割为~$C = \lceil L/n \rceil$~个不重叠的~$n$~比特~chunk，每个~chunk~可以取~$2^n$~种离散值中的一个。随后，通过一个位置感知查找表~$\mathcal{M} \in \mathbb{N}^{C \times 2^n}$~将每个~chunk~值确定性映射为~CLIP~词表中的一个~token——具体地，从~CLIP~词表中去除保留~token（如~[CLS]、[SEP]、[PAD]）得到有效词表~$\mathcal{V}_{\mathrm{valid}}$，随机排列后为每个~chunk~位置分配~$2^n$~个不重叠的~token。最终~token~序列由起始~token、各~chunk~的映射~token~和终止~token~拼接后零填充至~77~个~token~输入~CLIP~文本编码器。这种设计使相同数量的~token~能承载更多信息比特——例如~$n=2$~时，$C$~个~token~可编码~$2C$~比特消息，理论上将容量上限从~77~比特提升至~154~比特。

\paragraph{双分支解码器}
位压缩虽然提高了~token~效率，但也增加了解码难度：解码器不仅需要从~CLIP~图像编码器输出的~512~维特征中恢复每个~chunk~的离散状态，还需进一步从~chunk~状态中精确提取各比特。为此，BitC-3DGS~引入双分支解码器~$\mathcal{D}_M$~在两个粒度上联合优化。

第一个分支是块分支（chunk branch）：先通过两层~MLP~将~512~维~CLIP~特征投影为~$C \times d$~维~chunk~嵌入矩阵并叠加可学习位置编码，再经~4~头~transformer encoder~建模~chunk~间依赖关系，最后通过层归一化和共享线性分类器输出每个~chunk~在~$2^n$~个候选值上的~logits~$\bm{s}_i \in \mathbb{R}^{2^n}$，以交叉熵损失~$\mathcal{L}_{\mathrm{chunk}}$~监督~chunk~级分类。为将~chunk~级预测桥接到比特级监督，进一步构造固定二值码本~$\mathbf{B} \in \{0,1\}^{2^n \times n}$~，通过对~chunk~预测概率的~softmax~输出在码本上边缘化（marginalization）得到投影比特概率~$\tilde{\bm{p}} \in [0,1]^L$，以~BCE~损失~$\mathcal{L}_{\mathrm{proj}}$~提供从~chunk~结构推导比特的间接监督。

第二个分支是比特分支（bit branch）：将~512~维特征映射到隐藏表示后，经~$G$~个并行线性头产生~chunk~级~logits~，拼接后加入位置编码并重塑为~$L_g \times G$~矩阵（$L_g = L/G$），再经门控自注意力层（以可学习门控向量~$\bm{\gamma}$~调制注意力输出）精炼比特间依赖，最终输出直接比特~logits~$\hat{M} \in \mathbb{R}^L$，以~BCE~损失~$\mathcal{L}_{\mathrm{bit}}$~提供直接的比特级监督。

解码器的总预训练损失为三项加权和：$\mathcal{L}_{\mathrm{decoder}} = \lambda_s \mathcal{L}_{\mathrm{chunk}} + \lambda_p \mathcal{L}_{\mathrm{proj}} + \lambda_b \mathcal{L}_{\mathrm{bit}}$。

\paragraph{硬消息采样}
高容量设置下的另一个挑战是训练消息的组合覆盖不足：$L$~比特消息空间包含~$2^L$~种可能（如~$L=128$~时超过~$10^{38}$），随机采样仅能覆盖极小子集，且倾向于生成低汉明重量和结构重复的消息，导致解码器对已见消息过拟合并对未见消息泛化能力不足。BitC-3DGS~引入硬消息采样策略动态刷新训练缓冲区。具体地，维护固定大小的消息缓冲区~$\mathcal{D}$~（初始化为~$K$~条随机消息）和准确率记录字典~$\mathcal{S}$。每次训练迭代后，将解码错误的消息（至少一个比特不正确）记录到~$\mathcal{S}$~中，其准确率通过指数平滑更新（系数~0.5）以追踪历史难度。每个~epoch~开始时，缓冲区~$\mathcal{D}$~按线性递增的硬样本比例~$r_{\mathrm{hard}} = \min(1.0, \tau_0 + \alpha \cdot e)$~从~$\mathcal{S}$~中选取准确率最低的消息作为硬子集~$\mathcal{H}$~，剩余位置用未见消息填充以保持多样性，最终洗牌后作为新一轮训练集。达到预定义冻结~epoch~$E_{\mathrm{freeze}}$~后停止引入新消息，缓冲区固定为准确率最低的~$K$~条消息。这种渐进式难度聚焦策略使解码器在有限训练预算下覆盖更广的消息空间。

\paragraph{水印嵌入}
解码器预训练完成后冻结参数，进入水印嵌入阶段。渲染带水印的视角图像后施加随机视觉扰动（模拟各类降质），经~CLIP~图像编码器~$\mathcal{E}_I$~提取特征后由冻结解码器~$\mathcal{D}_M$~预测消息。3DGS~参数通过三项损失联合优化：比特恢复损失~$\mathcal{L}_{\mathrm{bit}}$~、图像保真损失（SSIM + L1 + LPIPS）和嵌入正则化~$\mathcal{L}_{\mathrm{off}}$~（沿用~GuardSplat~设计），$\mathcal{L}_{\mathrm{watermark}} = \lambda_{\mathrm{bit}} \mathcal{L}_{\mathrm{bit}} + \lambda_{\mathrm{image}}(\mathcal{L}_{\mathrm{rgb}} + \mathcal{L}_{\mathrm{lpips}}) + \lambda_{\mathrm{off}} \mathcal{L}_{\mathrm{off}}$。

在~Blender~和~LLFF~数据集上，BitC-3DGS~支持~128~比特消息容量，恢复精度与现有~64~比特方法（如~GaussianMarker、3D-GSW）的~64~比特设置相当，实现了翻倍的容量而不牺牲解码可靠性。

在追求高容量的同时，另一个重要趋势是将水印从"逐场景训练"推进到"一次训练、任意场景可用"的可泛化范式。下面介绍两类代表性工作。

与锚点~BitC-3DGS~（通过位压缩令牌化在逐场景训练框架内扩展比特容量）不同，可泛化水印的核心挑战是范式转换——从"每个场景都需要一次昂贵的微调"推进到"训练一次、无限场景复用"。GS-Marker（2025.03）~\cite{2503_18718} 针对逐场景迭代优化的低效和泛化性不足提出单次前馈水印方法。现有~3DGS~水印方法需要对每个场景单独优化，耗时且无法泛化到新场景。形式化地，给定由~$N$~个高斯基元组成的~3DGS~表示~$X \in \mathbb{R}^{N \times d}$（$d$~为每个基元的属性维度，涵盖位置、不透明度、颜色、缩放和旋转），目标是嵌入二值消息~$M \in \{0,1\}^{l_b}$~使得从任意视角的渲染图像中均可准确提取该消息，同时保持渲染保真度：$\hat{X} = \arg\min_{X} (\mathcal{L}_w(\hat{X}) + \lambda \mathcal{L}_r(\hat{X}))$，其中水印损失~$\mathcal{L}_w$~度量提取消息与原始消息的均方误差，渲染损失~$\mathcal{L}_r$~联合~MSE~和~LPIPS~确保水印前后的视觉一致性。

GS-Marker~的框架包含三个关键模块：（1）基于~PointNet++~的~3D~嵌入网络，将归一化的高斯属性与水印消息拼接后，通过级联的局部区域聚合块预测残差调整量，适配~3DGS~的无序点集结构。具体地，嵌入网络输出残差后先投影至~$[-1, 1]$~范围，再拆分为各高斯属性（位置、缩放、旋转、颜色、不透明度）的独立调整量，分别乘以各属性的缩放因子~$s_*$（位置~$s_M$、缩放~$s_S$、旋转~$s_R$、颜色~$s_C$、不透明度~$s_A$），实现对不同敏感度属性的差异化调整——位置属性需要极小残差而颜色属性可容忍更大修改；（2）失真层，同时模拟~3D~域（dropout、裁剪）和~2D~域（JPEG~压缩、裁剪、旋转）的降质，确保端到端可微训练；（3）基于~HiDDeN~架构的~2D~解码器，从渲染图像中恢复嵌入消息。解码器采用两阶段训练策略：先在输入~3DGS~渲染图像上预训练~HiDDeN~模型以获得良好初始化，再与~3D~编码器端到端联合微调以优化整体性能。

其关键创新是自适应标记控制（Adaptive Marker Control, AMC）机制——解决~3DGS~水印优化中的局部最优困境。设计动机是：预训练~3DGS~已经处于渲染损失的局部最优点（梯度极小），直接叠加水印目标会导致水印梯度主导优化方向，破坏渲染保真度或陷入不稳定收敛。AMC~通过在初始~3DGS~上施加受控扰动，将参数推离渲染损失的平坦区域，使渲染保真度和水印目标从优化起始就提供有意义的梯度。具体地，扰动由可学习的属性幅度~$\sigma_j$（不同属性的敏感度不同——位置属性需要极小扰动而颜色属性可容忍更大调整）和选择掩码~$m_i$（基于预设选择比率~$\rho$~随机确定，$m_i \sim \text{Bernoulli}(\rho)$，仅对部分高斯基元施加扰动以降低训练复杂度）联合控制：$\Delta_{i,j} = \sigma_j \cdot P_{i,j} \cdot m_i$，其中~$P_{i,j} \sim \mathcal{N}(0, 1)$~为随机噪声样本。此外，引入~W-Marker——一种随机生成的高斯基元集合，与输入~3DGS~拼接后为场景注入额外随机纹理，从而增强水印嵌入效果。嵌入网络同时对原始高斯基元和~W-Marker~进行处理（总数~$N + N_r$，其中~$N_r$~为~W-Marker~数量），输出统一的残差调整。训练损失为加权组合~$L = W_{img} L_{img} + W_{msg} L_{msg}$，其中图像损失~$L_{img}$~包含~MSE~和~LPIPS~两个分量，消息损失~$L_{msg}$~为消息均方误差。

GS-Marker~在解码精度和模型保真度上超越逐场景训练方法，且计算时间大幅减少——从逐场景的数十分钟降至单次前馈推理。

至此，高容量与可泛化水印的讨论告一段落。接下来转向另一条重要路线——在生成过程中同步嵌入水印，以及使水印抵抗模型压缩。

除了提高泛化效率，另一条路线是扩展水印的模态覆盖范围。与~GS-Marker~专注于可泛化框架不同，X-SG$^2$S（2025.02）~\cite{2502_10475} 首次统一~1D~到~3D~水印模型，解决现有方法仅支持单一模态水印而无法同时嵌入多模态消息的局限。其核心设计动机源于对~3DGS~点集水印固有困难的深入分析：（1）点集是无序的——无法像图像那样按空间位置直接定位水印嵌入位置，也无法确定点的排列顺序，尤其在嵌入多条水印时难以追踪位置；（2）单个点的容量有限——每个高斯仅能承载极少信息，需要将大信息量拆分后分布式嵌入全部点中，充分利用累积存储空间；（3）多模态数据格式异构——比特流、图像和点云属于不同向量空间，需要统一表示才能在同一框架内协同处理。X-SG$^2$S~的重要设计原则是不改变原始~3DGS~的参数形式或模型结构，水印模块以附加组件形式运作，兼容预训练~3DGS~管线甚至纯~.ply~点云文件。

基于此，X-SG$^2$S~的框架包含注入器和提取器两个核心模块。注入器中，消息首先被分块（patch）并创建冗余以增强鲁棒性：1D~二值消息和~3D~点云数据通过重复冗余编码，使任何子集均可恢复完整数据；2D~图像则通过特征稀疏~DCAE（深度压缩自编码器）编码为~$4 \times 4 \times 128$~紧凑特征，再经~4~层上采样卷积扩展为带冗余上下文的补丁序列并附加固定位置编码。训练期间随机丢弃部分补丁以模拟敌意攻击，恢复时由~MLP~预测位置编码并通过最近邻匹配将保留的补丁重新定位到正确位置，再经下采样卷积和精炼~MLP~重建完整特征并解码水印图像——这使得即使部分高斯点被恶意删除，仍可从剩余补丁恢复完整~2D~水印。

随后，自适应门（Self-Adaption Gate）通过两个互补得分的乘积筛选最适合承载水印的高斯位置。自得分（self-score）评估每个点在整体场景中的信息重要性：利用诱导集注意力（induced set attention）将~$O(N^2)$~的标准注意力降为~$O(N)$——仅用单个锚点向量~$I \in \mathbb{R}^{1 \times d}$~，先让锚点~attending~全部高斯点获得点集的全局摘要~$H$，再让全部点~attending~该摘要得到每个点的全局重要性~$O_s$，最后经~sigmoid~线性层输出标量得分。交叉得分（cross-score）评估每个高斯点与消息补丁的交互匹配度：采用高效~Transformer~近似标准交叉注意力，令高斯点为~$Q$、消息补丁为~$K$~和~$V$，通过行列分别归一化的近似~$O_c = \text{softmax}_\text{row}(Q)(\text{softmax}_\text{col}(K)^T V)$~降低内存开销。两种得分相乘后，通过~top-k~选择最终嵌入位置——得分高的位置意味着对渲染影响小且与消息交互响应强。对于多模态嵌入，后续模态的位置掩码在先前模态已选位置基础上排除重复，确保不同模态消息占用不同的高斯点。

在选定嵌入位置后，三个独立线性层分别将~1D、2D、3D~消息映射到统一向量空间。多头注入（XD-Injection Heads）为每个模态配备独立的~3$\sim$5~层~MLP（基于~PointNet~架构），将消息补丁与目标位置的~SH~系数拼接后生成注入后的~SH~参数——选择第~3、4~阶高阶~SH~系数作为载体，因为这些高阶~SH~对场景渲染贡献最小（仅包含少量高频信息），水印嵌入对视觉质量的影响最低。

提取器中，可学习选择门（4~层~MLP）从嵌入后的完整点集中识别携带消息的高斯位置——训练时将注入器产生的位置掩码作为真值监督。为避免提取器训练干扰注入器的梯度流，X-SG$^2$S~采用关键的梯度截断策略：将含水印的~3DGS~点集视为常数输入（切断梯度回传），使选择门仅学习定位而不影响已优化的注入参数。多头提取头从对应位置的~SH~参数中恢复各模态的隐藏消息。训练损失联合优化多项目标：1D~消息采用二元交叉熵（BCE）损失，2D~和~3D~消息采用~MSE~损失，SH~参数保真损失约束嵌入前后~SH~系数变化幅度，定位掩码交叉熵损失监督选择门的预测精度。此外，当与预训练~3DGS~管线联合训练时，可选地加入~2D~渲染图像~MSE~损失和~LPIPS~感知损失以进一步保护渲染质量。

与~GS-Marker~的前馈路线不同，MarkSplatter（2025.09）~\cite{2509_00757} 针对现有~3DGS~水印方法依赖逐场景微调的高昂计算开销，提出首个基于~Splatter Image~结构的可泛化水印框架。与~GS-Marker~直接在点云结构上操作不同，MarkSplatter~的核心创新~GaussianBridge~首先将非结构化的~3D~高斯点云转化为结构化的~Splatter Image~格式——后者将~3D~信息编码为~2D~图像张量（每个像素对应一个~3D~高斯的~14~维参数：位置、缩放、旋转、不透明度和颜色），使消息嵌入可直接通过标准神经网络操作完成，无需对每个场景单独迭代优化。GaussianBridge~通过轨道式多视角渲染获取输入图像，再利用~LGM~（Large Gaussian Model，一种前馈~3DGS~重建模型）将点云格式的~3DGS~转换为~Splatter Image~格式，并扩展为支持任意数量输入视角而非原始~LGM~的固定视角设定，确保信息完整性。

具体地，GaussianBridge~首先使用~Pl\"ucker~光线嵌入（Pl\"ucker ray embedding）密集编码相机位姿，将多视角图像与光线嵌入拼接为统一的~9~通道特征图~$F = \{C, R_o \times R_d, R_d\}$（其中~$C$~为图像颜色，$R_o, R_d$~分别为光线原点和方向），随后通过~LGM~的非对称~U-Net（编码器-解码器对称卷积网络）~架构（结合残差层和自注意力机制）进行处理。为增强跨视角信息传播，多视角特征在自注意力层前被展平并拼接（flatten and concatenate），使不同视角的信息充分交互。最终输出特征图包含~14~个通道，每个像素即一个完整参数化的~3D~高斯——这种结构化表示直接对齐~Splatter Image~格式，为可泛化水印提供了神经网络友好的操作空间。GaussianBridge~模块具有即插即用特性，也可扩展至~GRM、StereoGS~等更先进的高斯重建模型，从而同时支持~3D~高斯重建任务和~3D~高斯生成任务的水印保护。

水印嵌入采用属性内嵌入与跨属性约束的双策略设计：（1）属性内嵌入——仅在~Splatter Image~的颜色属性~$x_c$~上施加水印扰动，保留几何属性不变以避免破坏~3D~结构。编码器分别通过~ResNet~下采样网络将颜色属性~$x_c$~和位置属性~$x_\mu$~编码为特征表示~$f_c$~和~$f_\mu$，然后通过~AdaIN（自适应实例归一化）层利用几何特征~$f_\mu$~动态调制颜色特征~$f_c$~（实现外观与几何的联合推理），调制后的表示经多视角自注意力机制进一步精炼（显式建模视角间依赖关系并确保符合~3DGS~多视角几何约束）：$f = \text{MultiViewAttn}(\text{AdaIN}(f_c, f_\mu))$。消息嵌入发生在瓶颈层——可学习二值查找表将二值消息映射为潜在扰动并与特征~$f$~拼接。解码器随后通过额外的多视角自注意力层和~ResNet~上采样网络生成扰动图~$\delta = \theta_{emb}(x_c, \hat{m})$，最终含水印颜色属性为~$x_{cw} = x_c + \delta$，以最小感知失真嵌入可恢复信息。

（2）跨属性约束——提出高斯不确定性感知（GUP）热图，利用~Fisher~信息矩阵的简化~Hessian~近似估计每个高斯的不确定性值~$\mathbf{U} = \log(\nabla_\mathcal{G} I_\mathcal{G} \nabla_\mathcal{G} I_\mathcal{G}^T)$（其中~$\nabla_\mathcal{G}$~为~3DGS~模型参数的梯度，$I$~为渲染图像；取对数是为了重新排列指数级大的值），将不确定性渲染为~2D~热图~$\gamma = \mathcal{R}(\mathbf{U}, v)$~后逐元素调制水印扰动幅度，最终嵌入公式为~$x_w = x + \alpha \cdot \gamma \odot \delta$（其中~$\alpha$~为缩放因子）。值得注意的是，旋转参数~$r$~被从不确定性估计中排除，因为旋转属性的微小变化容易在边界区域产生明显的针状伪影（needling artifacts）。GUP~热图的设计动机是：已有基于参考图像的渲染贡献图或显著性图方法依赖真值数据，在~3DGS~生成模型的推理场景中不可用；而基于不确定性的估计无需参考图像，适用于无真值数据的场景。因此~GUP~确保在视觉敏感区域（不确定性高的区域对扰动更敏感）施加更小的扰动。

与~GS-Marker~基于~HiDDeN~的~2D~解码器不同，MarkSplatter~采用基于分割的消息提取机制——受~WAM（Watermark Anything Model）启发，结合~ViT（Vision Transformer，视觉 Transformer）~编码器进行全局特征提取和像素解码器将特征上采样至输入分辨率，同时输出水印区域分割掩码和逐像素消息~$y = \{y^{mask}, y^{msg}\}$~（共~$1+N$~比特，其中~1~位掩码标志该像素是否包含水印，$N$~位为嵌入消息），解决了水印物体仅占渲染图像极小区域时传统全局池化解码器失效的问题，且无需辅助判别器。提取器支持双重验证：既可从~2D~渲染图像中提取水印（$I_w = \mathcal{R}(x_w, V)$；$y^{mask}, y^{msg} = \theta_{ext}(I_w)$），也可直接从含水印的~Splatter Image~颜色属性中提取——无需场景特定调优，确保了~3DGS~模型参数层面的版权可验证性。训练损失为掩码损失和消息损失的线性组合~$\mathcal{L}_{total} = \lambda_1 \cdot \mathcal{L}_{mask} + \lambda_2 \cdot \mathcal{L}_{msg}$，两者均采用逐像素二值交叉熵。增强策略覆盖~2D~域（几何变换：缩放、旋转、平移、透视、裁剪；光度调整：噪声、JPEG~压缩、模糊、亮度/对比度偏移）和~3D~域（几何编辑：裁剪/平移/旋转；噪声注入），确保水印对~2D~渲染失真和~3D~表示篡改的双重鲁棒性。MarkSplatter~在~Objaverse~和~GSO~数据集上，32~比特容量下达到~94.41\%~位精度和~35.12 dB~PSNR，48~比特下达到~93.36\%~位精度和~33.99 dB~PSNR，推理时间仅需~2.5~秒——相比~GaussianMarker~的~117~秒和~WateRF~的~168~秒，速度提升~50$\times$~以上。

\subsubsection{生成模型水印与压缩鲁棒}

本路线介绍在~3DGS~生成过程中同步嵌入水印以及使水印抵抗模型压缩的方法。前两节的水印方法均针对"从多视角图像重建的~3DGS"进行后处理嵌入。然而，随着~AIGC（AI-Generated Content，人工智能生成内容）~技术的快速发展，~3DGS~的生成方式正在多样化——从文本/图像直接生成~3DGS~模型已成为现实。这带来了新的水印需求：水印需要在生成过程中同步嵌入，而非事后追加。同时，3DGS~模型的压缩存储和传输也对水印鲁棒性提出了额外挑战。

与锚点~GS-Marker~（通过前馈网络在已生成的~3DGS~上嵌入水印，实现"一次训练、任意场景可用"的可泛化后处理水印）不同，GaussianSeal（2025.03）~\cite{2503_00531} 将水印范式从"后处理"推进到"生成过程内嵌"——水印不是在~3DGS~资产创建之后追加，而是在~3DGS~生成模型的推理过程中同步嵌入。这一范式转换的核心动机有二：（1）后处理水印方法（包括~GS-Marker~的前馈路线和~3D-GSW~的微调路线）一次只能处理一个~3DGS~对象，对于大规模~3D~资产生成场景效率低下；（2）直接微调生成模型权重（如~AquaLoRA~在~Stable Diffusion~上的做法）会严重影响~3DGS~的生成质量。原因在于~3DGS~生成模型的~UNet~不经过去噪训练，其潜在编码对权重扰动远比扩散模型更脆弱。

GaussianSeal~首次提出~3DGS~生成模型的位水印框架，以~LGM（Large Gaussian Model）~为目标生成模型。LGM~以多视角图像为输入，通过~UNet~架构输出~14~通道的~3DGS~张量（3~维位置、1~维不透明度、3~维缩放、4~维旋转、3~维~RGB~颜色），每个像素对应一个~3D~高斯基元。GaussianSeal~的核心思路是将自适应位调制模块（ABM，Adaptive Bit Modulation Module）嵌入~LGM~的~UNet~的三个关键位置——输入阶段（conv\_in 之后）、中间阶段（encoder 之后）和输出阶段（decoder 之后），使每个阶段的特征图被叠加携带水印信息的调制信号。具体地，给定版权比特串~$m \in \{0,1\}^L$（$L$~为消息长度，默认~16~比特），每个调制模块首先通过~MLP~将~$m$~映射为与当前特征图空间维度匹配的噪声块（block size 为~8），再经卷积层调整通道数后与原始特征图相加：$x' = x + \alpha \cdot \text{Mod}_k(m)$，其中~$k \in \{\text{in, mid, out}\}$~对应三个嵌入位置。关键设计是自适应系数~$\alpha$、$\beta$、$\gamma$（分别对应输入、中间、输出三个阶段）——它们被初始化为~0.1~并通过独立的~AdamW（Adam with decoupled Weight decay，带解耦权重衰减的自适应优化器）~优化器学习更新，而非手动设定固定值。这样设计源于~3DGS~参数对扰动的高度敏感性：3DGS~生成张量的属性值（尤其是位置和缩放）即使微小偏移也会导致高斯点离散和不规则形变。自适应系数通过梯度驱动自动平衡水印解码精度与生成质量。

水印解码端采用预训练的~HiDDeN~解码器（训练时固定参数），从~LGM~生成的~3DGS~渲染图像中提取比特信息。解码器从渲染图像的~DWT~低频~LL~子带中恢复水印比特，这一选择基于低频子带对视觉质量更鲁棒且与高频纹理细节解耦的考虑。训练损失联合优化水印消息损失和~3DGS~保真损失：$\mathcal{L} = \lambda_{\text{gs}} \mathcal{L}_{\text{gs}} + \lambda_{\text{rgb}} \mathcal{L}_{\text{rgb}}$，其中~$\mathcal{L}_{\text{gs}}$~度量含水印~3DGS~张量与原始张量的~MSE~差异，$\mathcal{L}_{\text{rgb}}$~度量两者的渲染图像差异（联合~MSE~和~LPIPS~），$\lambda_{\text{gs}}=1000$~和~$\lambda_{\text{rgb}}=300$~的设定反映了~3DGS~张量空间需要更强的保真约束。训练仅在~LGM~的~Objaverse~数据集上进行（10000~对象训练、100~对象验证），且完全不修改~LGM~的原始权重——仅训练轻量级的调制模块和自适应系数，以最小训练开销实现高精度位解码。

GaussianSeal~在~Objaverse~数据集上的实验表明：（1）位解码精度方面，16~比特和~32~比特设置下均显著优于所有后处理基线方法（3DGS+HiDDeN、3DGS+WateRF、3D-GSW、GaussianMarker），因为水印在生成过程中深度融合于~3DGS~张量而非事后叠加到高斯属性表面；（2）生成质量方面，PSNR/SSIM/LPIPS~指标与无水印的原始~LGM~输出接近，证实自适应系数有效控制了调制扰动对~3DGS~属性的破坏；（3）效率方面，水印嵌入与生成推理同步完成，无需对每个~3DGS~对象执行额外的后处理步骤；（4）鲁棒性方面，对~3DGS~压缩（~C3DGS~25$\times$~压缩）和视角变化保持稳定的解码精度。此外，GaussianSeal~展现出良好的泛化能力——在~Objaverse~上训练的调制模块可直接迁移到~NeRF Synthetic~数据集，且扩展到~100000~对象时性能不退化。安全性分析通过~StegExpose~反隐写检测验证，ROC（Receiver Operating Characteristic，受试者工作特征）~曲线接近随机猜测基线，表明水印信号隐蔽性良好。

在锚点~GaussianSeal~将水印从后处理推进到生成过程嵌入之后，CompMarkGS（2025.03）~\cite{2503_12836} 将~3DGS~水印推向另一个关键场景——模型压缩。3DGS~模型体积庞大，需要压缩存储和传输。但现有水印方法（包括前述~3D-GSW、WATER-GS、GuardSplat~以及~GaussianSeal）均将水印嵌入高斯属性或网络参数中，一旦经历量化压缩——连续的~SH~系数和几何参数被截断为低精度离散值——水印引入的微小扰动即被抹除。实验表明，WateRF~在~HAC~压缩后比特准确率从~91\%~骤降至~54\%，3D-GSW~从~91\%~降至~53\%，GuardSplat~压缩前仅~80\%~、压缩后跌至~53\%，几乎完全丧失版权保护能力。CompMarkGS~首次将压缩鲁棒性作为~3DGS~水印的核心设计目标，使水印在经历量化压缩后仍保持接近无损的比特准确率。

CompMarkGS~选择锚点~3DGS~（如~Scaffold-GS）作为水印载体，将水印嵌入锚点特征~$f \in \mathbb{R}^d$~而非直接操作高斯属性。这一选择具有双重优势。安全性上，锚点特征通过~MLP~间接决定高斯属性，攻击者难以通过分析点云或高斯属性直接检测水印的存在。质量上，锚点特征对渲染质量的影响是间接的（经由~MLP~映射），不像位置、缩放、偏移等属性那样直接控制高斯行为，因此修改锚点特征对渲染质量的扰动更可控。具体地，引入与锚点特征同维的可学习水印嵌入特征~$f' \in \mathbb{R}^d$，通过~$f^w = f + \tanh(f')$~生成含水印的锚点特征。$\tanh$~函数将扰动限制在~$[-1, 1]$~的对称范围内。这一设计源自对锚点属性分布的深入分析：训练良好的锚点属性服从正态分布，无界扰动会增大方差导致梯度不稳定和收敛缓慢；而~sigmoid~的~$[0, 1]$~范围会引入正向偏置；仅~$\tanh$~的对称有界范围能保证稳定收敛和无偏嵌入。水印的解码与渲染解耦：对渲染图像施加~DWT，将低频~LL~子带输入预训练的~HiDDeN~解码器提取~$L$~比特消息，水印消息损失为标准二值交叉熵~$\mathcal{L}_{\mathrm{msg}}$。

核心创新是量化失真层（QDL，Quantization Distortion Layer）——训练时主动注入与实际压缩算法一致的量化噪声：$\tilde{f}^w_i = f^w_i + \mathcal{U}(-\frac{1}{2}, \frac{1}{2}) \cdot q_i$，其中量化步长~$q_i = Q_0 \cdot (1 + \tanh(r_i))$，$r_i = \mathrm{MLP}_q(f^w_i)$~通过小型~MLP~根据含水印特征自适应调整，将步长限制在~$[0, 2Q_0]$~范围内。其物理含义是模拟真实压缩中的取整误差——均匀分布噪声乘以量化步长等价于量化取整的舍入残差。这是首次将量化模拟作为数据增强应用于~3DGS~水印训练，迫使水印编码器学会在量化噪声下仍保持可解码的嵌入策略。实验证明了~QDL~的必要性：在~HAC~压缩下，CompMarkGS~的比特准确率从压缩前~95.95\%~仅微降至~95.92\%（降幅~0.03\%），而所有基线方法均遭遇~30--40\%~的断崖式下跌。

在渲染质量保障方面，CompMarkGS~提出频率感知锚点生长（FAG，Frequency-Aware Anchor Growth）策略来化解水印与质量增强的频率冲突。水印嵌入在渲染图像的低频子带中，而低频子带同时承载全局结构信息——如果质量增强也在低频子带进行，增强算法会将水印信号视为噪声并尝试去除，两者互相干扰。FAG~的策略是将质量增强从低频移至高频：对渲染图像和~GT~图像分别施加~DFT，经高通滤波提取高频分量后用~IDFT~重建高频图像，计算像素级~SSIM~误差并取中位数阈值生成二值掩码，仅在高频区域激活锚点生长，使新增锚点专注于补强高频细节而不干扰低频水印信号。HSV~损失进一步缓解颜色伪影：在~HSV~色相空间而非~RGB~空间计算图像损失，按色调范围构建~R/G/B~三通道的二值掩码定位颜色伪影区域，在掩码内计算~MSE~损失——这更准确地反映人眼对颜色偏移的感知。总损失函数为~$\mathcal{L}_{\mathrm{total}} = \lambda_{\mathrm{img}}(\mathcal{L}_{\mathrm{scaffold}} + \lambda_{\mathrm{hsv}}\mathcal{L}_{\mathrm{hsv}} + \lambda_{\mathrm{freq}}\mathcal{L}_{\mathrm{freq}}) + \lambda_{\mathrm{msg}}\mathcal{L}_{\mathrm{msg}}$，其中~$\mathcal{L}_{\mathrm{scaffold}}$~为~Scaffold-GS~原始重建损失，$\mathcal{L}_{\mathrm{freq}}$~为像素级~SSIM~误差的平均值，四项损失协同优化渲染质量与水印鲁棒性。

\subsubsection{统一保护与鲁棒性增强}

本路线介绍联合优化水印嵌入与对抗防御以实现统一保护，以及针对扩散编辑等新型攻击的鲁棒水印策略。随着水印攻击手段的日益成熟，单纯的水印嵌入已经不足以应对新型威胁——扩散编辑可以直接重新生成高频细节从而隐蔽地擦除水印信号。本节介绍的方法从两个方向应对这一挑战：一是将水印与对抗防御联合优化，实现统一保护；二是针对扩散编辑等特定攻击设计鲁棒水印策略。

本方法修改了~3DGS~的版权保护环节——用统一的双分支优化框架（水印+对抗）替代独立的水印或对抗编辑威慑——3DGS~的场景表示和渲染管线保持不变。GuardMarkGS（2026.05）~\cite{2605_12919} 首次提出统一保护框架，联合优化所有权追踪（水印）和编辑威慑（对抗）。现有保护方法仅覆盖单一维度：水印可追踪未授权使用但无法阻止恶意编辑；对抗编辑威慑可扰乱编辑但无法提供所有权证据。简单叠加两种目标也无法奏效——实验表明，直接将~DEGauss~风格的对抗项加到~3DGSW~风格的水印目标上，会大幅削弱水印恢复精度和渲染保真度，却仅提供有限的编辑威慑。这是因为水印恢复需要全场景优化的广泛支持（确保任意视角可解码），而在同一高斯集合上均匀施加对抗扰动会干扰水印恢复和渲染质量。GuardMarkGS~的核心设计正是解决这一矛盾：在单个优化框架中结合全场景水印目标（覆盖所有高斯的水印嵌入）和选择性对抗目标（仅在掩码选定的高斯子集上施加对抗扰动）。

具体地，GuardMarkGS~的总目标函数为~$\mathcal{L}_{\mathrm{total}} = \mathcal{L}_{\mathrm{wm}} + \lambda_{\mathrm{adv}}\,\mathcal{L}_{\mathrm{adv}}$，其中水印分支~$\mathcal{L}_{\mathrm{wm}} = \lambda_{\mathrm{msg}}\,\mathcal{L}_{\mathrm{msg}} + \lambda_{\mathrm{quality}}\,\mathcal{L}_{\mathrm{quality}}$~覆盖全部高斯，$\mathcal{L}_{\mathrm{msg}}$~为~32~比特所有权消息的二值交叉熵——水印从每个渲染视角的\textbf{低频小波子带}中解码，低频分量对常见失真更具鲁棒性；$\mathcal{L}_{\mathrm{quality}}$~结合~$\ell_1$~和~LPIPS~保真度项，约束受保护渲染与冻结参考渲染之间的视觉质量。对抗分支~$\mathcal{L}_{\mathrm{adv}} = \frac{1}{|V_t|}\sum_{v \in V_t}\!\big(\mathcal{L}_{\mathrm{render}}^{(v)} - \lambda_{\mathrm{lat}}\,\mathcal{S}_{\mathrm{lat}}^{(v)} - \lambda_{\mathrm{traj}}\,\mathcal{S}_{\mathrm{traj}}^{(v)} - \lambda_{\mathrm{xattn}}\,\mathcal{S}_{\mathrm{xattn}}^{(v)}\big)$~在多视角上取平均，其中~$\mathcal{L}_{\mathrm{render}}^{(v)}$~为受保护渲染与冻结参考之间的重建误差，三条分离项各自引导编辑轨迹偏离预期方向，而对抗梯度则通过软掩码权重和参数角色系数进行调制，避免均匀扰动。直觉上，水印分支确保"从任何视角都能解码所有权"，而对抗分支仅对关键区域施加"干扰编辑"信号。两者的空间支持不同，避免了互相干扰。

对抗分支包含三条策略协同转移编辑轨迹：（1）潜在锚点分离~$\mathcal{S}_{\mathrm{lat}}^{(v)} = \|z(\mathcal{R}_v(\mathcal{G}^{\mathrm{prot}})) - z(\mathcal{R}_v(\mathcal{G}^{\mathrm{ref}}))\|_2^2$（$z(\cdot)$~为~VAE~编码器）在潜在空间中将编辑方向推离原始锚点——即使受保护渲染与参考渲染在视觉上几乎相同，潜在表示已被拉远，使扩散编辑器在初始化阶段就偏离正确编辑方向；（2）去噪轨迹转移~$\mathcal{S}_{\mathrm{traj}}^{(v)}$~在扩散去噪过程中引导轨迹偏离预期编辑结果，使用冻结编辑器的去噪轨迹描述符~$\Phi_{\mathrm{traj}}$~度量受保护渲染与参考渲染在相同提示~$\psi$~和噪声~$\epsilon$~下的编辑方向差异；（3）交叉注意力转移~$\mathcal{S}_{\mathrm{xattn}}^{(v)} = \frac{1}{d_{\mathrm{mid}}}\big\|Q_{\mathrm{mid}}^{\mathrm{xattn}}\!\big(\mathcal{R}_v(\mathcal{G}^{\mathrm{prot}});\psi,t,\epsilon\big) - Q_{\mathrm{mid}}^{\mathrm{xattn}}\!\big(\mathcal{R}_v(\mathcal{G}^{\mathrm{ref}});\psi,t,\epsilon\big)\big\|_2^2$~扰乱编辑指令与图像区域在~U-Net~中间块文本条件交叉注意力中的对齐（$d_{\mathrm{mid}}$~为查询维度，$Q_{\mathrm{mid}}^{\mathrm{xattn}}$~为中间块交叉注意力查询 token 的池化描述符），通过扩大受保护和参考渲染之间查询向量的分离，使编辑器无法正确关联文本指令与图像空间位置。三项策略从潜在空间、去噪方向和提示-图像对齐三个层面互补地干扰编辑管线，实验消融表明三者组合才能达到最强威慑。

此外，基于更新显著性的高斯选择策略解决了"对抗信号应该施加在哪些高斯上"的空间分配问题。观察发现，在没有选择策略时，简单场景（如~bear）的对抗更新自然集中在主要物体周围，但复杂场景（如~face）的更新则空间散乱且低效。因此，GuardMarkGS~在优化前通过文本引导的~2D~语义分割（基于~SAM）生成多视角掩码，并将其转换为逐高斯的软系数：$m_i = \min\!\big(1, (s_i / (\tau_{\mathrm{mask}} + \epsilon))^{\gamma_{\mathrm{mask}}}\big)$，其中~$s_i$~为高斯~$i$~的掩码得分（累积其光栅化像素贡献与掩码的加权和），$\tau_{\mathrm{mask}}$~控制饱和阈值，$\gamma_{\mathrm{mask}}$~控制衰减陡度。最终对抗梯度被调制为~$\widetilde{\nabla}_{\theta_{i,k}}\mathcal{L}_{\mathrm{adv}} = m_i\,\rho_k\,\nabla_{\theta_{i,k}}\mathcal{L}_{\mathrm{adv}}$，其中~$\rho_k$~为参数角色系数——不透明度和~DC~颜色（对全局外观影响大）被保守对待，而尺度、旋转和视角相关残差特征则被允许承载更强的对抗信号。直觉上，这一机制相当于为每个高斯的每种参数赋予了不同的"扰动预算"：视觉敏感的参数少扰动，视觉鲁棒的参数多扰动。

从统一保护框架过渡到下一类方法，研究者的关注点转向了更具挑战性的扩散编辑攻击场景。

在~Mip-NeRF 360~和~Instruct-NeRF2NeRF~场景上的大量实验表明，GuardMarkGS~在统一评分~sUCPS（联合评估可追溯性、编辑威慑和保真度三个维度）上比最佳基线方法提升了~10.7\%，位精度达~97.23\%，PSNR~达~30.36 dB。消融实验验证了上述两个分配机制的互补作用：移除高斯选择导致位精度从~97.23\%~骤降至~65.02\%、PSNR~从~30.36~降至~28.74 dB；移除参数角色调制则使~sUCPS~从~0.8622~降至~0.8566。实验使用~NVIDIA RTX A6000 GPU，每场景优化~8~轮，超参数设置为~$\lambda_{\mathrm{msg}} = 0.1$、$\lambda_{\mathrm{quality}} = 1.0$、$\lambda_{\mathrm{lat}} = \lambda_{\mathrm{traj}} = \lambda_{\mathrm{xattn}} = 10^{-4}$。

GuardMarkGS~虽然实现了统一保护，但其对抗分支主要针对指令驱动编辑。而现实中更隐蔽的威胁来自于扩散模型对渲染图像或~3D~资产的直接编辑——这类编辑可以在保持视觉合理性的同时擦除水印。RDSplat（2025.12）~\cite{2512_06774} 首次设计同时抵御~2D~和~3D~扩散编辑的水印框架。现有~3DGS~水印方法仅对经典失真（JPEG~压缩、噪声等）具有鲁棒性，但在扩散编辑面前失效——扩散模型可在~2D~图像层面（对渲染图像做扩散编辑）和~3D~场景层面（直接编辑~3DGS~资产）隐蔽地擦除水印信号，同时保持视觉合理性。近期~2D~方法（如~VINE、EditGuard）虽然通过频率引导和对抗训练实现了扩散鲁棒性，但无法直接迁移至~3DGS：2D~方法仅在渲染图像中嵌入水印，缺乏新视角泛化能力；其解码器在固定视角图像上训练，无法处理新视角渲染的几何可变性。RDSplat~的关键观察建立在对扩散模型的频谱分析之上：Fourier~分析表明，扩散编辑保留~48.53\%~的低频能量但仅保留~0.09\%~的高频能量——扩散模型的去噪过程本质上是低通滤波器，保留低频结构而重新生成高频细节。与~VINE~隐式引导水印向低频迁移的策略不同，RDSplat~直接将水印嵌入原生低频高斯基元，从根本上确保水印信号占据扩散编辑所保留的频谱带。

RDSplat~通过频率感知基元选择（FAPS）将~100~比特水印专门嵌入低频高斯基元，并冻结所有其他基元。FAPS~的设计基于三阶段筛选逻辑，联合捕获基元的空间尺度和观测分辨率以识别低频嵌入目标。第一阶段计算逐基元采样率~$\tilde{\nu}_k = \mathrm{median}_{n: G_k \in V_n}(f_n / z_k^n)$，其中~$f_n$~为相机焦距，$z_k^n$~为深度——表征每个基元在训练相机中的采样密度，高采样率表示信号被良好分辨，低采样率则意味着潜在的混叠风险。第二阶段评估基元在给定空间尺度下是否真正属于低频：Mip~分数~$m_k = s_k^2 \tilde{\nu}_k^2 / \lambda^2$（$s_k = \max_d s_{k,d}$~为主尺度轴，$\lambda$~为~Nyquist~乘子）衡量基元空间范围相对于采样分辨率的比值——$m_k \gg 1$~认证基元空间大且被充分采样，是稳定的低频嵌入目标；$m_k \approx 1$~则表示基元处于混叠边界，应避免。第三阶段引入方向平衡分数~$b_k = (\|\sum_{n: G_k \in V_n} \hat{d}_k^n\|_2 + \epsilon)^{-1}$（$\hat{d}_k^n$~为从基元到相机的单位观测方向），确保被选基元从多视角被一致观测而非仅从单侧可见——当相机均匀分布时方向向量相互抵消，$b_k$~取大值；当相机集中在单侧时方向向量累积，$b_k$~取小值。两个分数经~min-max~归一化后乘法组合~$\xi_k = \bar{m}_k \cdot \bar{b}_k$——乘法形式确保两个条件必须同时满足，任一分数接近零的基元都被排除。最终保留按~$\xi_k$~排序的前~$\lfloor K \cdot r \rfloor$~个且在至少~20\%~训练视角中可见的基元作为水印载体，其余全部冻结。

训练效率通过代理策略实现。直接将扩散编辑器纳入训练循环在计算上不可行，而~Fourier~分析表明~Gaussian blur~与扩散编辑方法表现出近乎相同的频率保留行为（48.53\%~低频保留~vs.~0.09\%~高频），因此~RDSplat~用~Gaussian blur~作为计算高效的扩散编辑代理。增强管线将频率增强~$\mathcal{F}$（Gaussian blur、JPEG~压缩、噪声等）与几何增强~$\mathrm{Gaug}$（旋转、弹性形变、随机裁剪、缩放恢复）组合为统一管线，在单~RTX 4090 GPU~上~3--7~分钟完成微调。消融实验验证了代理策略的必要性：完全移除增强使训练时间减半，但扩散攻击鲁棒性一致下降；而用真实扩散前向传播替代~blur~代理将训练时间增加至~69.3~分钟（~$11\times$~），仅带来微弱的鲁棒性提升。消融还揭示了嵌入位置的关键作用：将水印嵌入高频基元虽然经典鲁棒性略高，但扩散鲁棒性显著下降且训练时间暴涨~$17\times$。这直接证实了扩散编辑摧毁高频信号，使其成为不可靠水印载体的结论。

专用解码器~GeoMark~基于~ViT-S/16~（Vision Transformer，视觉~Transformer~架构）和空间周期秘密嵌入，联合抵抗扩散编辑和新视角渲染的几何变换。现有~2D~解码器各有局限：HiDDeN~处理经典失真但在~VAE（Variational Autoencoder，变分自编码器）~增强下产生可见伪影；VINE~实现扩散鲁棒性但缺乏新视角解码所需的几何不变性；OmniGuard~的~CNN（Convolutional Neural Network，卷积神经网络）~解码器抵抗扩散编辑但在宽基线几何变换下失效。GeoMark~从水印图像提取~[CLS]~token~特征~$f = F_{\mathrm{ViT}}(\phi(I_W))$，通过线性投影恢复~$L$~比特预测（$\hat{w} = W_{\mathrm{out}} \mathrm{LN}(f) + b_{\mathrm{out}}$），其中~ViT~的~patch~级自注意力提供比~CNN~更强的几何变换不变性。空间周期秘密嵌入将~100~比特水印映射为~$32 \times 32$~的空间瓦片并通过复制覆盖全分辨率——任何不小于瓦片尺寸的局部裁剪都包含一份完整的水印副本，确保新视角渲染中裁剪区域的可解码性。GeoMark~在~MIR Flickr~数据集上分两阶段训练：第一阶段施加渐进式几何和光度增强，第二阶段加入~VAE~重建增强。训练完成后解码器被冻结，仅选定的低频基元在嵌入阶段被优化。消融实验表明，用~HiDDeN~替换~GeoMark~在经典攻击下位精度达~0.94~但在扩散攻击下崩溃至~0.59；用~VINE~替换则在旋转攻击下接近随机猜测（0.49），确认现有~2D~解码器无法满足~3DGS~的几何不变性需求。

在~Blender、LLFF、Mip-NeRF 360~和~Instruct-NeRF2NeRF（IN2N）四个基准数据集上的实验覆盖了七种~2D~扩散攻击（随机再生、InstructPix2Pix~全局编辑、DiffEdit~全局编辑、SD1.5~修复、SDXL~修复、SD~超分辨率、Latent~超分辨率）和迭代~3D~编辑（Instruct-GS2GS）。RDSplat~在~100~比特容量下（水印空间从~$2^{48}$~扩展至~$2^{100}$），经典鲁棒性（8~种攻击平均位精度~0.811，TPR@1\%FPR~为~0.903）和扩散鲁棒性（7~种扩散攻击平均位精度~0.701，TPR@1\%FPR~为~0.659）均达到竞争力水平，PSNR~为~27.24、干净渲染位精度~0.952。

在最具实战威胁的~IN2N~场景上（扩散编辑与其训练分布对齐、语义连贯，接近真实攻击者使用编辑工具的方式），RDSplat~对比~GaussianMarker~在扩散攻击上取得~0.82 vs. 0.65~的位精度优势（+0.17），其中~SDXL~修复下~1.00 vs. 0.80、DiffEdit~下~0.88 vs. 0.63。灵敏度分析显示位精度随扩散步数单调下降：10--20~步时保持~$>$0.8，超过~80~步降至~0.6~以下，160~步趋于随机猜测（0.5），确认强再生对当前水印方案构成严重威胁。迭代~3D~编辑测试进一步揭示了一个更根本的局限：Instruct-GS2GS~在~100~次迭代时仍可检测水印，但~300~次迭代后位精度大幅退化至接近随机——因为每次迭代对不同的训练视角施加扩散编辑，累积覆写承载水印的低频结构，即使~FAPS~选定的基元也无法承受无限次的生成式改写。这表明迭代~3D~编辑构成比单次~2D~攻击更严重的威胁，仍是当前方法的开放挑战。

RDSplat~的代价是较低的视觉保真度（PSNR 27.24），反映了从鲁棒性优先角度所做的刻意权衡。未来方向包括原生~3D~编码器-解码器联合训练、扩散感知优化和逐场景自适应嵌入比例。

从静态~3DGS~扩展到动态~4DGS~场景，水印技术面临时间维度的新挑战。

除了针对扩散编辑的鲁棒水印，水印技术的时间维度扩展同样值得关注——当场景从静态~3DGS~演进到包含运动信息的~4DGS~时，水印嵌入必须兼顾时间一致性。与锚点~GuardMarkGS~（保护静态~3DGS~场景免受指令驱动编辑威胁）不同，4D-GSW（2026.05）~\cite{2605_22342} 将水印扩展到~4DGS~动态场景，面临全新的时间维度挑战。4DGS~资产的显式参数化特性使其极易被提取、再分发或恶意篡改运动轨迹，但与隐式表示的水印不同，~4D~水印需要确保版权信息从任意~360°~视角均可恢复——将版权信息"烘焙"进~4D~辐射场中。此前唯一的~4D~隐写方法~Hide-in-Motion~依赖不透明度启发式而非显式运动物理建模，在复杂非刚性变形下~FVD（Fréchet Video Distance，衡量视频生成质量的核心指标）~严重退化。更根本的是，传统逐帧水印注入破坏了运动轨迹的物理连续性，导致严重的时间闪烁和~FVD~崩溃。4D-GSW~的核心目标是嵌入~$L=48$~比特二值水印序列~$B \in \{0,1\}^L$，通过策略性调制外观流形~$S = \{s_i^t\}_{i,t}$（即~SH~系数的零阶~DC~分量）实现，使水印信号随动态场景几何流畅演化，同时严格保持光度保真度。

4D-GSW~的运动感知框架建立在运动学分析之上：定义高斯基元~$p_i$~的时空曲率（STC）指标~$\kappa_i^t = \|v_i^t \times a_i^t\|_2 / (\|v_i^t\|_2^3 + \epsilon)$（其中~$v_i^t = (x_i^t - x_i^{t-1})/\Delta t$~和~$a_i^t = (v_i^t - v_i^{t-1})/\Delta t$~分别为瞬时速度和加速度），量化轨迹的运动不稳定性。该指标源自经典曲线几何与运动感知理论，衡量速度向量在大小和方向上的变化强度。高~STC~区域对应"动态瞬间"（如碰撞或急转弯），其参数流形高度敏感且脆弱，注入水印梯度会扭曲运动语义导致严重视觉伪影。因此~4D-GSW~将~STC~映射为水印嵌入置信权重~$w_i^t = \exp(-\beta \kappa_i^t)$，作为梯度路由的门控信号：稳定区域（$w_i^t \approx 1$）允许高信噪比水印注入，高动态区域（$w_i^t \approx 0$）实施物理隔离以保护运动完整性。由于~$\kappa_i^t$~的动态范围可能很大，训练中对~$w_i^t$~施加~mini-batch normalization~后跟~sigmoid~截断，确保权重在~$[0,1]$~区间内良好分布，防止权重稀疏化。

全局一致性通过联合~HMM-MRF（Hidden Markov Model - Markov Random Field，隐马尔可夫模型-马尔可夫随机场）~能量最小化模型实现。关键挑战是~4DGS~优化过程中基元经历频繁的致密化、剪枝和重采样，导致帧间拓扑突变，基于固定索引的对应关系失效——这种对应断裂在一致性损失中引入系统性偏差，训练中表现为不规则振荡。4D-GSW~将高斯场的时间演化建模为随机过程，利用最优传输（OT，Optimal Transport）理论解决这一隐变量匹配问题——将相邻时间戳的基元集合视为离散测度~$\mu_t$~和~$\mu_{t-1}$，通过几何感知的熵正则化~Wasserstein~距离~$\mathcal{W}(\mu_t, \mu_{t-1}) = \min_{\Pi} \sum_{i,j} \Pi_{ij} \|x_i^t - x_j^{t-1}\|_2^2 - H(\Pi)$~求解最优传输计划（其中~$H(\Pi)$~为熵惩罚项，通过~Sinkhorn~算法高效求解），用重心映射~$\hat{z}_i^{t-1} = \sum_j \Pi_{ij}^* z_j^{t-1} / \sum_k \Pi_{ik}^*$~恢复对齐的先验状态（$z = [x, s]^\top$~同时包含空间坐标和外观属性），确保一致性正则化在真实物理运动残差上操作。为保持训练效率，~OT~对齐仅在子采样的关键帧和高斯基元子集上执行。整个框架被形式化为~MAP~推断问题：能量泛函~$E(Z_t) = L_{data} + \lambda_{con} L_{con}(Z_t, \hat{Z}^{t-1})$~中，$L_{con}$~作为时空马尔可夫先验，运动学权重~$w_i^t$~充当该先验的精度因子——在稳定区域强制相位锁定以锚定水印比特，在动态瞬间自适应放松约束，确保正则化在真实物理运动残差而非拓扑噪声上操作。门控一致性损失~$L_{con} = \sum_i w_i^t (\|z_i^t - \hat{z}_i^{t-1}\|_2^2 + \lambda \sum_{j \in N_i} \omega_{ij} \|z_i^t - z_j^t\|_2^2)$~联合约束时间轨迹和~$K$-近邻空间邻域~$N_i$~内的水印相位一致性（$\omega_{ij}$~为空间邻接权重）。

各向异性梯度路由机制则在反向传播时对水印相关梯度施加运动学权重调制：$\nabla_{s_i^t} L_{total} = \nabla_{s_i^t} L_{rec} + w_i^t \cdot \nabla_{s_i^t} (\lambda_{wm}L_{wm} + \lambda_{wav}L_{wav} + \lambda_{con}L_{con})$，从而仅对外观属性（DC~分量~$s_i^t$）的门控水印梯度确保水印嵌入严格解耦于光度重建保真度和运动优化。频率域水印监督采用~DWT：对时间戳~$t$~视角~$v_j$~的渲染图像执行~DWT~分解为低频子带~$I_{LL}$~和高频子带集~$I_H = \{I_{LH}, I_{HL}, I_{HH}\}$，水印消息从稳定的~$I_{LL}$~子带中经预训练解码器~$D$~解码（~$L=48$~比特，BCE~损失~$L_{wm}$），高频子带则通过小波子带损失~$L_{wav}$~保持渲染保真度。总损失函数~$L_{total} = \lambda_{wm}L_{wm} + \lambda_{wav}L_{wav} + \lambda_{con}L_{con} + L_{rec}$~中，$L_{rec}$~包含了光度重建、~LPIPS~感知约束、几何边界掩码、生成增强和~SDS（Score Distillation Sampling，分数蒸馏采样）~几何一致性等多个子项，全面维护~4D~资产的结构和光度完整性。变分分析进一步揭示，该优化框架在连续极限下收敛到各向异性扩散-反应~PDE（Partial Differential Equation，偏微分方程）：$\lambda_T r(x) - \lambda_S \nabla \cdot (D(x) \nabla r(x)) = G_{wm}$，其中~$D(x) = \nabla^2 w(x)$~为曲率导出的扩散张量，$r(x) = z^t - \hat{z}^{t-1}$~为残差状态场。

这一理论框架解释了运动保真度的保持机制——稳定区域扩散张量的传导率促进信号沿流形传播，高动态区域~$D(x) \to 0$~自然惩罚对脆弱运动轨迹的急剧更新。实验在~Consistent4D~数据集上基于~SC4D~骨干（通过~Linear Blend Skinning~以~$M \leq 512$~个稀疏控制点驱动~$N$~个密集高斯）进行，单~NVIDIA RTX A6000 GPU~上每场景训练约~1~小时。~4D-GSW~在~FVD~指标上仅从~1297.18~微增至~1331.19（而朴素~2D~水印集成导致~FVD~飙升至~1601.50），位精度达~98.05\%（无攻击下达~99.10\%），PSNR~和~SSIM~甚至超过未加水印的~SC4D~基线——论文将此归因于~HMM-MRF~能量模型为辐射场提供了结构化的时空先验正则化。论文也指出局限性：~STC~指标假设~$C^2$~连续的高斯轨迹，在极端非刚性事件（如拓扑断裂）处该连续性被破坏，运动学流形不可微分，导致嵌入权重~$w_i^t$~局部振荡、扩散屏障减弱，在奇异区域水印鲁棒性有所退化。

\subsubsection{可解释水印与安全性分析}

本路线介绍对水印嵌入位置与机理进行可解释分析，以及对水印安全性进行系统评估的方法。前面几节的水印方法大多将"哪些高斯承载水印信号"视为优化过程的黑盒决策，缺少对水印嵌入位置和机理的可解释分析。与此同时，水印攻击技术也在快速发展，使得安全性评估成为不可或缺的研究方向。

从前面关注鲁棒性的方法过渡到这里，研究视角从"如何使水印更鲁棒"转向"如何理解水印为什么工作以及是否真的安全"。

与~GuardMarkGS~将"哪些高斯承受对抗信号"视为基于~2D~语义分割的掩码选择问题不同，Where, What, Why（2026.03）~\cite{2603_08809}~将分析视角从"防御有效性"转向"可解释性"。核心问题是：水印信号究竟嵌入在哪些高斯上（Where），嵌入的是什么内容（What），以及为什么选择这些载体而非其他（Why）。这一视角转变源于对现有方法"载体选择黑盒"的诊断：无论是~GuardMarkGS~的~SAM~掩码投影还是~RDSplat~的频率基元筛选，载体决策的证据链条对用户不可见、不可审计，也无法提供逐高斯的归属分析。

\paragraph{三专家载体评估}
Where, What, Why~的关键设计决策是将载体评估完全锚定在~3DGS~参数空间而非渲染域。它不依赖像素级梯度或~2D~语义分割，而是从原生参数（位置~$x_i \in \mathbb{R}^3$、尺度~$s_i \in \mathbb{R}_+^3$、旋转~$q_i$、不透明度~$\alpha_i$、~SH~系数~$h_i$）直接提取特征，确保视角一致性并降低计算开销。三个专家按语义分组参数，各自独立评分：
\begin{itemize}
\item 几何专家（$C_{\mathrm{geo}} = \{x, s, q\}$）通过尺度各向同性~$Iso_i = \bar{s}_i / \|s_i\|$、邻域旋转一致性~$RotCons_i = (1/k)\sum_{j \in N_k(i)}|q_i \cdot q_j|$~和紧凑足迹~$\sigma_i = (\prod_{d} s_{i,d})^{1/3}$~度量结构稳定性——各向同性高斯在不同视角下渲染贡献更一致，适合承载水印；形状极端各向异性的高斯则对视角敏感，作为载体会引入不可预测的解码波动。
\item 外观专家（$C_{\mathrm{app}} = \{\alpha, h^{(0)}, h^{(1)}\}$）通过~AC~高频能量比~$\sigma_i^{\mathrm{hf}}$、双侧不透明度门控~$g(\alpha)$~和~DC~强度带通~$c_i$~度量跨视角外观一致性——高~AC~能量意味着高斯承载视角相关的高频细节，修改它会引发可见伪影；低不透明度高斯在~$\alpha$-混合~中贡献微弱，其参数修改对渲染几乎无影响但对水印解码也不可靠。
\item 冗余专家（$C_{\mathrm{red}} = \{x, s, q, h^{(0)}\}$）通过重叠加权邻域相似度~$r_{ij}$（综合~DC~颜色和形状相似度，权重~$w_{ij} = \exp(-d_{ij}^2 / (\sigma_o^2(\Delta s_i^2 + \Delta s_j^2)))$~近似投影重叠）度量分布密度和可替代性——在高密度区域中，单个高斯的渲染贡献可被邻域替代，修改它对视觉影响最小。
\end{itemize}
每个专家~$k$~将其特征~$z_k(i)$~映射为证据包~$E_k(i) = [U_k(i), S_k(i)]$，其中质量分数~$S_k(i) = \mathrm{Norm}(z_k(i))$~和不确定性~$U_k(i) = \mathrm{Norm}(\mathrm{Disp}_{N(i)}(z_k) + \mathrm{Penalty}_k(i))$~解耦——邻域离散度（$\mathrm{Disp}$）度量局部一致性，专家特定惩罚项捕获各自关注的风险信号。这一解耦设计使后续门控可以同时考虑质量和置信度，避免低置信度的高分高斯被误选。

\paragraph{安全预算感知门}
SBAG~将载体选择建模为"在自适应预算约束下的多目标排序"问题。排序阶段将三专家证据映射为代理分数~$R_k(i) = \mathrm{clip}(S_k(i) - \beta\, U_k(i), 0, 1)$（$\beta$~为跨专家共享常数），分别解读为几何稳定性~$R_1$、外观安全性~$R_2$~和冗余确定性~$R_3$。为避免假设任何单一专家的主导地位，采用对称点积效用~$u_i = (R_1(i) \cdot R_2(i) \cdot R_3(i))^{1/3}$——乘法形式确保三个维度必须同时达标，任一维度接近零的高斯被自然排除。

预算分配引入场景自适应的有效比特系数~$\mu_{\mathrm{eff}} = \mu_0\,\bar{v}/\bar{v}^*$，其中~$\bar{v}$~为场景平均可见度（通过一次~DC+不透明度渲染获取逐高斯可见度~$v_i$~后取均值），$\mu_0$~为嵌入设计常数，从而将所需载体数量~$B = M / \mu_{\mathrm{eff}}$~与场景复杂度动态关联。可行集通过分位数约束定义：$\mathcal{F} = \{i \mid R_k(i) \geq Q_q(R_k),\; v_i \geq Q_q(v)\}$，排除低于~$q$-分位数阈值的低质量候选。在~$\mathcal{F}$~内按~$u_i$~进行确定性水位选择~$\mathrm{WM}_0 = \mathrm{top}\text{-}B\{u_i : i \in \mathcal{F}\}$。为弥合视角覆盖缺口，SBAG~进一步构建紧凑证据向量~$e_i = \mathrm{Norm}(R_1(i), R_2(i), R_3(i), v_i, h_i, c_i)$，计算~$\mathrm{WM}_0$~的原型~$\mu_e = (1/|\mathrm{WM}_0|)\sum_{i \in \mathrm{WM}_0} e_i$，按余弦相似度招募近邻扩展为~$\mathrm{WM}_{\mathrm{parent}} = \mathrm{WM}_0 \cup \mathrm{WM}_{\mathrm{prox}}$。最终，每个父高斯分裂为~$N_s$~个子高斯，一个子高斯路由至水印分支~$\mathrm{WM}_*$，其余子高斯充当视觉补偿器~$\mathrm{VIS}$——这一分裂-路由设计使得水印载体和视觉补偿器在空间上紧密相邻，补偿器可以精确地中和载体嵌入引入的局部伪影。

\paragraph{通道级组掩码：正交梯度路由。}
通道级组掩码确保水印载体和视觉补偿器在反向传播中接收到完全正交的梯度信号。具体而言，对于高斯~$i$，根据~Trio~证据和单次渲染先验预计算逐组掩码~$m_g^{\mathrm{wm}}(i)$~和~$m_g^{\mathrm{vis}}(i)$，其中~$g \in \{\delta_{\mathrm{dc}}, \delta_{\mathrm{rest}}, \delta_{\mathrm{opa}}, \delta_{\mathrm{rot}}, \delta_{\mathrm{sca}}\}$~覆盖五个参数通道。视觉掩码~$m_g^{\mathrm{vis}} = \max(\mathrm{clip}(\mathrm{mean}(w_g^{\mathrm{vis}}[\mathrm{VIS}]), 0, \mathrm{cap}_g), \mathrm{floor}_g)$~对~VIS~集合取通道均值并钳位到~$[\mathrm{floor}_g, \mathrm{cap}_g]$；水印掩码~$m_g^{\mathrm{wm}} = \mathrm{clip}(\mathrm{median}(w_g^{\mathrm{wm}}[\mathrm{WM}_*]), 0, \mathrm{cap}_g)$~对~$\mathrm{WM}_*$~集合取通道中值。梯度路由规则为：

\begin{equation}
\nabla_{\theta_i^g} \mathcal{L} = \begin{cases} m_g^{\mathrm{wm}}(i)\,\nabla_{\theta_i^g} \mathcal{L}_{\mathrm{wm}}, & i \in \mathrm{WM}_* \\ m_g^{\mathrm{vis}}(i)\,\nabla_{\theta_i^g} \mathcal{L}_{\mathrm{vis}}, & i \in \mathrm{VIS} \end{cases}
\end{equation}

为完全消除梯度耦合，$\mathcal{L}_{\mathrm{vis}}$~和~$\mathcal{L}_{\mathrm{wm}}$~分两次独立反向传播执行——$\mathrm{WM}_*$~和~$\mathrm{VIS}$~的参数更新在数学上正交，消除了联合优化中的目标冲突。

\paragraph{解耦微调：双目标正交优化。}
视觉目标~$\mathcal{L}_{\mathrm{vis}} = \lambda_{\mathrm{rec}}\mathcal{L}_{\mathrm{rec}} + \lambda_{\mathrm{lpips}}\mathcal{L}_{\mathrm{LPIPS}} + \lambda_{\mathrm{wav}}^{\mathrm{high}}\mathcal{L}_{\mathrm{wav}}^{\mathrm{high}}$~组合~L1~重建、~LPIPS~感知约束和多级~DWT~高频子带保真度——仅作用于~VIS~高斯。水印目标~$\mathcal{L}_{\mathrm{wm}} = \lambda_{\mathrm{wm}}^{\mathrm{clean}}\mathcal{L}_{\mathrm{wm}}^{\mathrm{clean}} + \lambda_{\mathrm{wm}}^{\mathrm{eot}}\mathcal{L}_{\mathrm{wm}}^{\mathrm{eot}} + \lambda_{\mathrm{wav}}^{\mathrm{low}}\mathcal{L}_{\mathrm{wav}}^{\mathrm{low}}$~仅作用于~$\mathrm{WM}_*$~高斯，其中水印嵌入在~DWT~低频（$LL$）子带中，同时用~$\mathcal{L}_{\mathrm{wav}}^{\mathrm{low}}$~正则化低频失真。鲁棒性通过期望变换（EOT）实现：对每个载体~$i$~在可见视角~$V_i$~上平均解码 logits（$\hat{\ell}_i^{\mathrm{cl}} = (1/|V_i|)\sum_{v \in V_i} D_{\theta}(\mathcal{R}_v)$，$\hat{\ell}_i(t) = (1/|V_i|)\sum_{v \in V_i} D_{\theta}(t(\mathcal{R}_v))$），变换~$t \sim p(\mathcal{T})$~从标准退化族（模糊、旋转、缩放、裁剪、噪声、JPEG~压缩）中采样。耦合微调的对照实验揭示了优化冲突：在耦合设置下~BitAcc~出现瞬态退化而~PSNR~持续改善，呈现清晰的~trade-off；而解耦优化收敛更快且全程维持更高的~PSNR~和~BitAcc。

\paragraph{实验结果。}
在~Blender、LLFF~和~Mip-NeRF 360~共~25~个场景上的实验表明，Where, What, Why~在~PSNR~上比最优基线提升~+0.83~dB、位精度提升~+1.24\%。在~7~种图像失真攻击（高斯噪声~$\sigma=0.1$、旋转~$\sim\pi/6$、缩放至~75\%、模糊~$\sigma=0.1$、裁剪至~40\%、JPEG~质量~50\%、噪声+裁剪+JPEG~复合攻击）和~3~种模型级扰动（移除~20\%高斯、克隆~20\%高斯、参数噪声~$\sigma=0.1$）下均保持鲁棒。消融实验验证了三个组件的互补性：移除~SBAG~导致载体选择退化，位精度和渲染质量同时下降；移除组掩码使梯度冲突重现，视觉保真度受损；移除解耦微调则暴露~BitAcc~瞬态退化——三者缺一不可。所有实验在单~NVIDIA A800 GPU~上完成，训练~2--10~轮，使用冻结~HiDDeN~解码器（32/48/64~比特），低贡献高斯（$V_i < 10^{-8}$）被剪枝。

可解释性为单个水印方法提供了微观层面的理解，而从宏观层面审视整个~3DGS~水印领域的安全现状同样重要。安全性分析方面，Zhao et al.（2026.02）~\cite{2602_03878} 发布首篇~3DGS~知识产权保护综述，针对该领域研究碎片化的现状——现有工作横跨水印、隐写、篡改定位和编辑防护等多个任务但缺乏统一分类体系，且对底层技术共性缺少系统分析——提出自底向上的分析框架，从三个互补层次系统梳理。

综述的第一层聚焦底层的高斯扰动机制，这是所有~3DGS~知识产权保护方法共享的技术基础。

综述沿三个维度归纳~24~种现有方法的共有技术选择：（1）属性选择机制——决定哪些高斯参数承载扰动，分为~SH~专有扰动（仅修改颜色外观而保持几何完整，视觉保真度高但优化不稳定且统计可检测性高）、混合属性扰动（同时修改~SH~系数和几何相关属性如位置、协方差、不透明度，保护能力更强但几何参数对扰动极为敏感，微小位置偏移即可导致渲染质量显著退化）和辅助属性耦合（在外部附加专用于承载保护扰动的新属性通道而非修改任何原生高斯参数，可完全保持视觉质量且支持高容量嵌入，但由于辅助通道不与语义属性或渲染关键参数耦合，攻击者可在不降质视觉质量的情况下剪枝或清除这些通道，实际鲁棒性有限）三种策略；（2）分布策略——决定扰动如何在高斯基元间分配，包括基于傅里叶频率引导的局部策略（高频分量关联精细外观细节而感知敏感，低频分量编码稳定结构和语义信息而鲁棒性更强，混合频域策略通过协调不同频域的扰动实现比单频域策略更优的不可感知性-鲁棒性权衡）、基于不确定性引导的局部策略（选择对渲染质量贡献较小的高不确定性高斯作为载体，实现不可感知的修改，但高不确定性高斯的数量和空间分布有限，嵌入容量和灵活性受到约束）和全局策略（在所有高斯基元上均匀分布扰动，最大化信息容量并避免对特定空间或频谱区域的依赖，广泛用于隐写场景，但需要精细优化以保持视觉保真度）；（3）扰动注入管线——包括逐资产微调（直接优化高斯参数实现场景特定的鲁棒性-保真度权衡，但泛化性差且时间成本随资产数量线性增长）、可泛化映射（在大型~3DGS~数据集上训练通用编码器，单次前向传播即可向任意新~3DGS~资产注入扰动，但复杂扰动建模引入训练不稳定性且多~GPU~训练开销大）和生成中注入（将扰动直接集成到~3DGS~生成模型中，通过交叉注意力或~U-Net~块输出在创建时即赋予资产保护性扰动，实现大规模内容生产的高效率，但仅保护对应生成模型产出的资产而无法适用于已有内容）。

综述的第二层覆盖中层的被动与主动保护范式。

被动保护实现事后追踪溯源而不主动限制资产功能，涵盖三种范式：（1）3DGS~水印——嵌入低容量二进制标识符实现版权归属验证，包括基于不确定性感知嵌入的~GaussianMarker、通过对抗训练提升鲁棒性的~GuardSplat、针对扩散编辑场景将水印嵌入低频协方差分量的~RDSplat、支持局部侵权检测的~NGS-Marker、强调压缩鲁棒性的~CompMarkGS、采用双网络设计实现跨场景水印的~GS-Marker、将高斯重组为网格化泼溅图像进行~CNN~嵌入的~MarkSplatter、同时支持白盒和黑盒提取的~WATER-GS、在重建过程中从多视角图像迁移水印的~MantleMark，以及将水印集成到生成管线实现"天生安全"资产的~GaussianSeal；（2）3DGS~隐写——以高容量和不可感知性为目标嵌入多模态载荷（如位序列、2D~图像或~3D~资产），包括通过~SH~系数调制实现高容量嵌入的~GS-Hider~和~ConcealGS、通过位移和透明度映射嵌入~3D~内容的~Splats in Splats、引入通用编码器-解码器管线的~GaussianStego、通过~LSB~操作实现超快嵌入的~Hide A Bit、利用密钥保护解码器实现访问控制的~KeySS、在~Scaffold-GS~锚点内采用解耦加密的~SecureGS、通过频域分解在高频域隐藏消息的~StegaGaussian、利用跨模态冗余的~X-SG$^2$、在~4DGS~时空变形场中嵌入载荷的~Hide-in-Motion，以及在生成推理过程中通过交叉注意力层注入消息的~InstantSplamp；（3）3DGS~篡改定位——目前仅有~GS-Checker~一种方法，无需预嵌入扰动即可直接分析被篡改的~3DGS~资产，通过为每个高斯附加~3D~篡改属性并采用~3D~对比机制揭示局部不一致性来检测疑似操纵。

主动保护则预先限制未授权使用，目前主要聚焦编辑防护：AdLift~通过将有界~2D~对抗扰动提升为~3D~高斯防护来阻止指令驱动的编辑，DEGauss~通过视角聚焦梯度融合和双差异优化增强跨视角鲁棒性。综述还提出~3DGS~知识产权保护的统一目标框架：最小化保护容量损失和鲁棒性损失的加权和，同时约束不可感知性预算，涵盖视觉保真度、保护能力和鲁棒性三个基本属性。

综述的第三层审视上层的鲁棒性威胁，从经典失真扩展到~AIGC~时代的新型攻击：（1）传统~2D~和~3D~失真——2D~失真包括几何变换（旋转、裁剪）、光度变化（亮度、对比度调整）和信号退化（高斯模糊、噪声注入、JPEG~压缩）；3D~失真分为几何失真（剪枝、克隆、致密化改变高斯密度分布，空间变换扰乱位置和协方差从而错位几何依赖性编码）和退化失真（高斯噪声同时破坏几何和外观，模型压缩通过量化和参数合并抑制低幅度扰动）；（2）生成式净化——2D~域方法如~DiffPure~利用扩散逆向过程系统消除对抗噪声，3D~域方法如~GSPure~和~GMEA~利用高斯冗余性剪枝对解码具有区分度但对重建贡献最小的元素，这些方法在保持视觉保真度的同时有效破坏主流保护方案，暴露了生成式净化作为关键但研究不足的脆弱性；（3）生成式编辑——扩散驱动的~2D~和~3D~编辑模型作为强力的语义修改器，通常充当激进低通滤波器，过滤掉~3DGS~资产中的高频保护性扰动，多数水印方法的高频嵌入在生成式编辑面前极为脆弱。

综述揭示了三个关键研究空白：（1）缺乏对高斯扰动底层机制的统一分析——不同保护任务中扰动属性选择、分布策略和注入管线的权衡缺乏系统比较；（2）对~AIGC~驱动的新型鲁棒性威胁的系统性表征不足——大多数方法未显式纳入净化感知训练或评估；（3）被动保护占据主流而主动保护严重不足——全生命周期资产保护从初始不可学习性到事后可删除性远未被探索。

基于此综述提出~6~个未来方向，分属三个维度：鲁棒性方面（R1）抵抗生成式净化与编辑——当前防御常以牺牲视觉质量换取对扩散编辑的抗性，同时实现高不可感知性和对生成式失真的抗性仍是重大挑战；（R2）跨表示存活问题——3DGS~资产在分发管线中常被转换为网格、体素或~NeRF~等其他格式，大多数方法无法在这些表示转换中存活，需探索表示无关的扰动技术，将信息嵌入高层几何拓扑或频域；（R3）统一基准测试——缺乏标准化评估框架阻碍方法间的公平比较，亟需涵盖从传统失真到生成式攻击的标准化攻击套件、多样化测试场景和统一的不可感知性与鲁棒性指标。效率方面（E1）可泛化与通用扰动——大多数方法依赖逐资产微调，对大规模资产生管线不实用，需开发在大规模~3D~数据集上预训练、推理时单次前向传播即可注入扰动的通用方法；（E2）重建与生成~3DGS~的内置安全——随着~3D~资产主要来源向前馈重建和生成管线转移，保护机制应直接集成到模型权重或潜空间中，使产出资产天生具有保护能力而无需事后处理。

保护范式方面（P1）从被动验证到主动防护的范式升级——现有研究主要聚焦事后验证，全生命周期保护应覆盖从初始不可学习性、不可迁移性到数据遗忘的完整链条，为~AIGC~时代的~3DGS~资产知识产权治理建立技术可执行性、可追溯性和法律可审计性。

Deng et al.（2026.02）~\cite{2602_02602} 从方法论层面指出当前水印研究的根本问题：缺乏明确的安全目标和现实的威胁模型定义。该位置论文的核心论点是~3DGS~水印的安全假设和攻击面比传统音视频媒体更为复杂。保护同时跨越模型域（3D~表示参数）和渲染域（2D~输出图像），攻击者可能访问其中任何一个；3DGS~模型在分发和处理过程中频繁重参数化和重排序，可破坏基于索引锚定的保护机制；几何与外观紧密耦合且不同参数的感知敏感度高度不均匀；实时渲染接口允许查询任意视角的渲染图像，使自适应攻击（如多视角探测和针对检测器的优化攻击）成为可能。论文系统回顾了传统音视频水印数十年的工业实践经验与教训——从~1990~年代数字音频/图像/视频的大规模复制催生知识产权保护需求，到~SDMI（Secure Digital Music Initiative，安全数字音乐倡议）公开检测预言机被攻破的教训——强调~3DGS~水印研究正在重复传统媒体曾经走过的弯路：许多工作未明确对抗访问能力便声称安全性，将场景简单划分为"模型级"或"渲染输出级"而缺乏细粒度分类，若干系统直接复用~HiDDeN~等公开检测器（其检测流程可被攻击者完整复现），且未引入密钥管理机制，暴露于逆向工程、自适应查询和伪造攻击等标准安全风险之下。

论文的核心贡献是形式化的安全评估模型。它将~3DGS~水印系统定义为三阶段管线——设置（指定密钥~$K$~并经安全通道分发至验证门户）、嵌入（将~$L$~比特消息~$m \in \{0,1\}^L$~注入原始模型~$M$~生成水印模型~$M_w$）和验证（独立的验证门户提供有限查询访问，支持模型级检测器~$D_M$~和图像/视频级检测器~$D_I$~两种检测通道，仅向认证用户授予有限次数的验证查询以防止自适应查询攻击）。基于此，论文定义了七元组访问向量~$\mathcal{A} = [\text{Key-}K, \text{Access-}M, \text{Access-}E, \text{Access-}D, \text{Oracle-}D, \text{Oracle-}R, \text{Access-}R]$~来刻画攻击者的能力——每个条目取值为~1（拥有该访问权限）或~0（没有），默认假设攻击者不知道密钥且不能直接访问原始模型（即~$\text{Key-}K = \text{Access-}M = 0$）。三种经典对抗箱体设置被定义为访问向量的子集：黑箱~$\mathcal{A}_{bb}$（攻击者仅通过预言机查询与系统交互，如向渲染预言机提交输入以获取渲染输出）、白箱~$\mathcal{A}_{wb}$（攻击者完全访问水印模型以及嵌入和检测算法的实现细节，但密钥和原始模型不可用）和灰箱~$\mathcal{A}_{gb}$（介于两者之间的任意中间访问模式，满足~$\mathcal{A}_{bb} \subset \mathcal{A}_{gb} \subset \mathcal{A}_{wb}$）。论文将所有潜在攻击按攻击者操作域分为三个层级：3DGS~级攻击（直接操纵高斯参数、剪枝、重采样）、图像/视频级攻击（仅访问渲染输出，进行压缩、裁剪、重编码、屏幕录制等后处理）和神经网络级攻击（利用公开或可查询的检测器，通过学习优化擦除或伪造水印信号）。

论文倡导场景驱动的研究范式——首先明确应用和安全目标（不同场景的成功标准、可接受成本和权衡差异显著），然后将对抗能力形式化为访问向量，用访问向量的子集或簇定义威胁等级以实现系统评估和定量比较，最后确保算法设计匹配所选威胁模型（特别关注密钥管理和检测器可用性），并按统一协议评估效果与风险。论文以嵌入式指纹取证（embedded fingerprinting，即~Hollywood forensic watermarking）场景为示范，详细分析了不同箱体设置下的具体子场景及其访问向量：黑箱下的云端重流（捕获~3D~资产渲染帧或录屏再广播）、渲染固定视角的被动泄露（最弱攻击者，访问向量全为零）和利用追踪门户的迭代探测；灰箱下的离线取证追踪（追踪合法接收者、供应链泄露、共谋重建）和预言机引导的模型规避（结合模型操纵与反复查询追踪接口进行迭代水印移除）；白箱下攻击者已知检测机制全部细节的极端威胁。在密钥设计方面，论文梳理了传统媒体水印中经实践验证的密钥机制：黑箱设置下的密钥聚合（将逐帧或逐片段的统计量累积为带置信度的嫌疑集合）和挑战-响应协议（由密钥决定渲染什么视角或从服务器指定的视角集中选择什么内容）；灰箱设置下的密钥个体化码字生成和随机嵌入规则（支持共谋追踪，从泄露副本解码出非授权分发的来源）；白箱设置下的密钥子集选择、投影和置换（确保获取权重不等于可靠规避）以及签名或消息认证码（绑定声明以降低诬陷风险）。

论文还构建了参考框架组织现有方法并澄清设计选择与对抗假设的映射关系。该框架从比特载荷（bit string）、嵌入域、嵌入阶段和检测器设计四个维度对~24~种代表性方法进行结构化分析，揭示了若干关键发现：许多方法仅在单一载荷长度上报告结果，掩盖了载荷-鲁棒性-保真度的三重权衡；全局嵌入虽能保持保真度但某些几何敏感参数（如位置和方向）的微小扰动可导致不成比例的几何伪影；后验优化（post-hoc optimization）是最广泛采用的稳定范式但耗时显著，单次前馈（single-forward）方案虽避免迭代优化却依赖跨域变换引入信息损失；多项工作直接复用~HiDDeN~作为~2D~渲染水印的检测器且未引入密钥约束，面临标准安全风险。

论文还讨论了三种代表性替代观点——工程导向的鲁棒性基准足够论（论文指出鲁棒性是安全的必要非充分条件）、公开检测器促进开放评估论（引用~SDMI~教训指出公开检测器将攻击面扩大至梯度攻击和选择性擦除）和密钥可选论（指出这本质上是"通过隐匿实现安全"，违反已知安全原则）——逐一分析其局限性。最后，论文实现了经典扩展频谱嵌入方案作为~3DGS~场景下的可复现基线：将嵌入消息表示为~$B$~比特~$\pm 1$~载荷向量，通过码分复用（CDM）在原生~3DGS~参数空间的变换域载波池上嵌入，三个独立密钥分别控制载波选择、声明-消息绑定和伪随机扩频码生成，采用非盲检测协议（给定原始模型和嫌疑模型形成变换域残差，通过相关解码和符号检验恢复每个比特），在~Blender、LLFF~和~Mip-NeRF 360~三个标准基准上实验验证。

实验表明，即使渲染差异在视觉上不可察觉，扩展频谱方案仍会导致可测量的保真度下降（~PSNR~和~SSIM~随嵌入强度~$\delta$~降低和载荷长度~$B$~增加而单调劣化），证实了~3DGS~水印场景下经典"容量-鲁棒性-保真度"权衡的存在性和可量化性。论文最终呼吁社区采用场景驱动的最低披露标准：每项工作应声明部署场景和安全目标、明确威胁模型和访问向量、说明方案是否采用密钥、按统一协议报告效果与风险，从而推动~3DGS~水印研究从即兴式数据隐藏走向场景扎根、可比较且可验证的安全研究。

与~Zhao et al.~综述的互补关系在于：Zhao et al.~从技术分类角度自底向上梳理已有方法，聚焦于归纳和分类；Deng et al.~则从安全分析角度自顶向下审视根本性问题，聚焦于批判和规范。两者共同构成了~3DGS~水印领域"技术全景"与"安全反思"的双重审视。

从理论分析转向实战攻击，下面介绍两种从攻击者角度评估水印脆弱性的方法。

理论分析揭示了水印的脆弱性，而实战攻击则进一步验证了这一判断。水印攻击方面，GSPure（2025.11）~\cite{2511_22262} 首次提出~3DGS~水印净化框架。与锚点~Where, What, Why~（从防御者角度优化水印载体选择）形成对照，GSPure~从攻击者角度逆向利用了同一事实：通过参数优化嵌入的水印与特定高斯基元在位置、不透明度和渲染贡献方面表现出强相关性。GSPure~首先排除了随机剪枝、特征缩放和噪声注入等朴素方法（这些方法要么无法完全移除水印，要么同时严重降质原始场景），提出两阶段净化管线：（1）视角感知高斯权重累积——为每个高斯基元计算其在~$N$~个渲染视角下的平均~$\alpha$-混合~渲染贡献权重~$\omega_k$，核心观察是水印相关的高斯在不同视角间呈现不一致或视角特异性的渲染行为（渲染贡献低），而原始场景相关的高斯在多视角间贡献一致且较高；（2）几何准确的特征聚类——将标准化后的高斯位置、不透明度和视角累积权重拼接为高维特征向量~$F_k$，采用~HDBSCAN（Hierarchical Density-Based Clustering，层次密度基聚类）自适应捕获不同场景的多样分布，以聚类为单位而非单个高斯进行剪枝——水印相关高斯通常在几何和属性上呈现强相关性，聚类级剪枝可以完整移除整个水印相关簇。自适应剪枝策略通过两个动态阈值（聚类剪枝因子和噪声剪枝因子）基于全局平均高斯贡献自动调整，GSPure~使水印~PSNR~降低最多~16.34~dB~且原始场景保真度损失不到~1~dB。

GMEA（2025.08）~\cite{2508_07263} 首次提出针对~3DGS~水印系统的通用黑盒攻击框架。与~GSPure~（利用渲染贡献差异进行聚类剪枝）不同，GMEA~从优化角度建模水印移除问题。在黑盒设置下（无法访问水印检测器内部），GMEA~将水印移除建模为大规模多目标优化问题——同时最小化视觉失真~$F_1$（L1~距离与~SSIM~的加权组合）和水印残留~$F_2$。$F_2$~的关键设计是间接目标函数：通过最小化卷积网络提取特征的标准差~$D(F_c) = \sqrt{\frac{1}{H'W'}\sum_{h,w}(F_c(h,w) - \mu_c)^2}$~使特征图失去信息量，从而"致盲"水印检测器——这是一种模型无关策略，对所有依赖卷积特征提取的解码器都有效。面对~3DGS~模型庞大的搜索空间（数十万高斯各含~14+维参数），采用~K-Means~空间聚类将高斯按~xyz~坐标分为~$k$~组，每组独立优化：个体的基因表示~$x^{(i)} = [m^{(i)}, c^{(i)}]$~包含二值掩码（决定剪枝哪些高斯）和颜色扰动向量（修改剩余高斯的~DC~颜色分量），通过~NSGA-II（Non-dominated Sorting Genetic Algorithm II，非支配排序遗传算法）~风格的多目标进化算法（非支配排序+拥挤度选择）逼近~Pareto~最优前沿，平衡视觉保真度和水印破坏。最终将~$k$~个优化子模型的并集重构完整对抗模型，可有效移除~1D~和~2D~水印，同时保持高视觉保真度。

MarkCleaner（2026.02）~\cite{2602_01513} 则专门针对语义水印提出了微几何扰动攻击策略。与~GSPure~和~GMEA~（在~3DGS~参数空间操作）不同，MarkCleaner~选择在~2D~图像域操作。其核心发现是语义水印的鲁棒性依赖于精确的相位对齐——空间位移可以打破这种对齐从而移除水印。MarkCleaner~采用~2DGS~解码器显式参数化几何扰动：将输入图像表示为一组~2D~高斯基元（具有可学习的位置、协方差、颜色和不透明度），通过学习高斯中心点的微偏移场（而非强迫网络回归像素级变化）实现空间重排。这样设计的动机是~2DGS~的显式表示解耦了几何（位置、协方差）与外观（颜色、不透明度），水印失效主要源于空间错位而非语义退化——因此只需学习几何位移即可破坏水印同时保持语义内容。训练时采用双域掩码编码：频域掩码（随机低通/高通滤波）破坏频谱水印签名，空间域随机掩码迫使编码器从不完整信息中提取连贯语义特征。关键创新是微几何扰动监督策略——不将模型输出约束为重建原始输入，而是监督其重建经微小旋转、平移、缩放变换后的图像，迫使模型将几何位移内化为生成过程的一部分。推理时这种学习到的行为自然产生语义一致但几何位移的输出，从而破坏水印解码。为保持语义一致性，引入冻结~DINOv2~编码器的自监督特征对齐（余弦相似度约束），以及~AdaLN（Adaptive Layer Normalization，自适应层归一化）~条件化解码器。

\begin{takeaway}
数字水印是 3DGS 知识产权保护的基础：频谱嵌入利用正交变换将水印分散到所有高斯基元，几何不变性设计抵抗仿射变换和剪枝攻击，而净化攻击的出现则推动了鲁棒水印的持续演进。
\end{takeaway}

%=====================================================================

\subsection{近期方法脉络：投毒与取证检测、溯源与知识产权、Steganography for Gaussian scenes}

这一组近期工作延续 3DGS pipeline、3DGS assets、3DGS，主要把改动集中在增量改进、安全化等环节。下面只保留相对既有路线的变化，便于和前文方法直接对照。

\paragraph{Characterizing Detectability in 3DGS Poisoning}
Characterizing Detectability in 3DGS Poisoning（2026.06）~\cite{2606_03499} 延续 3DGS pipeline 在“投毒与取证检测”上的基本设定，主要改变是：Poison-3DGS：stage-wise detectability benchmark，分析输入、训练中间表示和最终模型中的投毒信号。3DGS 的显式高斯表示和可微渲染能力保持为基础。
其设计动机是：将3DGS安全研究从攻击成功率扩展到检测和取证，指出不同攻击阶段会在不同中间表示中留下异质信号。
具体而言，论文围绕3DGS重建管线的多阶段中间产物构建检测任务，覆盖幻觉物体注入、计算成本放大和水印类后门等攻击。
适用边界是：作为基准工作，检测器泛化和真实攻击覆盖范围仍需扩展。

\paragraph{GaussTrace}
GaussTrace（2026.06）~\cite{2606_10612} 延续 3DGS assets 在“溯源与知识产权”上的基本设定，主要改变是：GaussTrace：基于高斯参数统计画像、假设驱动编辑模拟和证据式LLM推理构建有向溯源图。3DGS 的显式高斯表示和可微渲染能力保持为基础。
其设计动机是：提出面向3DGS资产传播和二次修改的provenance analysis框架，把内在参数统计和模拟编辑证据转化为结构化溯源推理。
具体而言，方法提取高斯属性的统计特征，模拟常见编辑操作产生辅助证据，再让LLM执行链式结构化推理以推断模型之间可能的转换路径。
适用边界是：依赖编辑模拟库覆盖范围；LLM推理可解释但需防止幻觉；对强对抗变换仍需验证。

\paragraph{IBRSteG}
IBRSteG（2026.06）~\cite{2606_30024} 可作为“Steganography for Gaussian scenes”中的新锚点候选。相较于 3DGS，主要改变是：前馈 3D Gaussian embedding 与 Gaussian Attributes Steganography。3DGS 的显式高斯表示和可微渲染能力保持为基础。
其设计动机是：将 3D steganography 表述为前馈 3D Gaussian embedding，缓解既有方法与特定场景刚性耦合的问题。
具体而言，通过 GAS 等模块把 secret scene 的高斯属性嵌入 cover/stego Gaussian scene，在保持视觉不可察觉的同时支持跨场景泛化。
这一锚点判断来自当前分类中的机制关系：候选锚点：安全化，且相关性 0.95；后续方法可围绕 利用 Gaussian 属性承载隐写信息和秘密三维场景。 继续比较和扩展。
适用边界是：安全性依赖 threat model；压缩、编辑或重训练可能破坏隐藏信息。
\section{对抗攻击与防御}
\label{sec:21_2}

3DGS~的对抗安全研究涵盖两个方向：（1）攻击——通过注入微小扰动破坏~3DGS~的重建质量或欺骗下游检测器；（2）防御——保护~3DGS~系统免受此类攻击。3DGS~的对抗攻击面比~2D~图像和~NeRF~更广：它不仅在训练数据层面（输入图像、相机位姿）和模型参数层面（高斯属性）存在攻击入口，而且~3DGS~的可微渲染管线使攻击者可以利用梯度信息精确优化对抗扰动。攻击者可以通过~3DGS~的~CUDA~光栅化器反向传播损失函数的梯度，精确定位哪些像素或高斯参数的微小修改会产生最大的攻击效果。

\subsection{对抗攻击与防御方法}
\label{ss:21_2}

本节讨论三条技术路线：（1）重建质量攻击；（2）下游检测器攻击；（3）防御方法。

\subsubsection{重建质量攻击}

本路线介绍以破坏3DGS重建质量为目标的攻击方法，通过投毒训练数据或操纵SfM管线从根源上破坏三维重建。

本节关注以破坏~3DGS~重建质量为目标的攻击方法。它们通过投毒训练数据或操纵~SfM（Structure from Motion，从运动中恢复结构）~来从根源上破坏~3D~重建管线。

Poison-splat（2024.10）~\cite{2410_08190} 修改了~3DGS~的训练数据输入环节——用经~TV~最大化优化的投毒图像替代原始干净图像——3DGS~的自适应密度控制机制保持不变，但被恶意利用为攻击的放大器。该方法揭示了~3DGS~一个此前未被关注的安全盲区：攻击目标不是重建质量，而是计算资源，极端情况下可导致拒绝服务（DoS，Denial of Service）。

\paragraph{威胁模型与问题形式化}

Poison-splat~在数据投毒（data poisoning）框架下运作。受害者是~3DGS~服务提供商，使用多视角图像和相机位姿组成的数据集~$D = \{(V_k, P_k)\}_{k=1}^N$~训练模型。攻击者的能力假设是：可以向数据集注入投毒数据——例如伪装成普通用户提交恶意数据，或隐蔽篡改合法用户提交的数据。攻击者不需要知道受害者的底层硬件设备。在白盒设置下，攻击者知道受害者的具体~3DGS~配置；在黑盒设置下，攻击者仅知道受害者使用某种~3DGS~变体。攻击者从干净数据~$D$~出发，仅修改图像内容（不修改相机位姿配置文件），生成投毒数据集~$D_p = \{(\tilde{V}_k, P_k)\}_{k=1}^N$。这一攻击可以形式化为如下~max-min~双层优化问题：

\begin{equation}
\max_{D_p} \, \mathcal{C}(G^*), \quad \text{s.t.} \quad G^* = \arg\min_{G} \mathcal{L}_G(D_p)
\end{equation}

其中外层优化最大化计算成本~$\mathcal{C}(G)$（如~GPU~显存占用和训练时间），内层优化是受害者的标准训练过程——最小化重建损失~$\mathcal{L}_G$（即~$(1-\lambda)L_1 + \lambda \mathcal{L}_{\text{SSIM}}$~的加权组合）以获得最优模型~$G^*$。直觉上，攻击者在寻找"最糟糕的训练数据"，使受害者在不知情的情况下用这些数据训练出一个计算代价极高的模型。

\paragraph{Baseline：自适应密度控制}

3DGS~的自适应密度控制（ADC，Adaptive Density Control）是其训练过程中周期性执行的内置机制。具体而言，在每个密度控制周期，3DGS~检查每个高斯在视图空间的位姿梯度（view-space positional gradient）$\partial_i^L$：当梯度超过预设阈值~$\tau_g$（默认~$0.0002$）时，触发分裂（split）或克隆（clone）操作以增加该区域的高斯密度；当高斯不透明度~$\alpha_i$~低于阈值~$\tau_\alpha$~时，执行剪枝（prune）移除冗余高斯。直觉上，纹理复杂、边缘锐利的区域产生更大的视图空间梯度，因此被分配更多高斯来恢复细节——这是~3DGS~实现高质量重建的关键机制，但同时也是一把双刃剑，因为高斯数量与~GPU~显存和训练时间呈正相关。

\paragraph{攻击思路：用高斯数量近似计算成本}

攻击者的直接目标是最大化训练的计算成本，但计算成本本身不可微分，无法直接优化。Poison-splat~的核心观察是：由于~3DGS~的计算成本（显存占用、训练时间）与高斯数量~$|G|$~呈正相关（论文~Figure 2~验证了这一关系），因此可以用高斯数量作为计算成本的代理指标。然而，高斯致密化操作本身也不可微分，梯度无法直接反向传播到输入图像。Poison-splat~进一步发现一个关键经验规律：3DGS~倾向于为具有更高总变分（TV，Total Variation）分数的区域分配更多高斯。TV~是衡量图像锐度/非平滑度的经典指标。直觉上，表面越粗糙、纹理越不光滑，3DGS~需要越多高斯来拟合~2D~投影中的所有细节。基于此，高斯数量可以用训练数据的~TV~分数来近似：$|G| \approx S_{\text{TV}}(D)$。

\paragraph{核心公式}

TV~分数的具体定义为：

\begin{equation}
S_{\text{TV}}(\tilde{V}_k) = \sum_{i,j} \sqrt{|\tilde{V}_k^{i+1,j} - \tilde{V}_k^{i,j}|^2 + |\tilde{V}_k^{i,j+1} - \tilde{V}_k^{i,j}|^2}
\end{equation}

即对投毒图像~$\tilde{V}_k$~中每个像素与其右邻和下邻像素的差值取~$L_2$~范数后求和。直观理解，TV~分数越高意味着图像中相邻像素的差异越大，即图像越"锐利"或越"不平滑"。

为了平衡攻击强度与隐蔽性，攻击者可选择施加~$L_\infty$~范数约束：

\begin{equation}
\tilde{V}_k = \text{clamp}(\tilde{V}_k, V_k - \epsilon, V_k + \epsilon)
\end{equation}

即每个像素的扰动幅度不超过~$\epsilon$。当~$\epsilon$~设为最大值（如~255）时，约束退化为无约束优化。

\paragraph{实现：代理模型迭代优化}

直接对每张视角图像独立最大化~TV~分数并不能产生有效攻击——独立优化会在不同视角的投毒图像间引入不一致性，反而削弱攻击效果。Poison-splat~的解决方案是训练一个~3DGS~代理模型（proxy model）~$G_p$~来保证多视角一致性。具体迭代流程为：（1）在干净数据上初始化代理模型~$G_p$；（2）采样一个随机相机视角~$k$，用代理模型渲染图像~$V'_k = R(G_p, P_k)$；（3）以渲染图像为起点，沿~TV~分数梯度的符号方向迭代~$\tilde{V}_k = \tilde{V}_k + \text{sign}(\nabla_{\tilde{V}_k} S_{\text{TV}}(\tilde{V}_k))$，并裁剪到~$\epsilon$-ball~内；（4）用~$(V'_k, \tilde{V}_k)$~对代理模型执行一步优化更新，模拟受害者的训练行为；（5）重复以上步骤直至收敛。这一"外层优化（投毒数据）-内层优化（代理模型更新）"的交替展开近似求解了原始的~max-min~双层优化问题，同时通过~3DGS~渲染函数的多视角一致性天然保证了投毒数据的跨视图一致性。

\paragraph{攻击效果}

在~NeRF-Synthetic、Mip-NeRF 360~和~Tanks and Temples~三个数据集上的实验表明：即使施加~$\epsilon = 8/255$~的小幅扰动约束，Poison-splat~也能将峰值~GPU~显存增至~2~倍以上（总体峰值超过~24~GB，超过~RTX 3090/4090/A5000~等常见~GPU~的显存容量）。在无约束设置下，攻击效果更加极端——例如在~NeRF-Synthetic~的~ship~场景中，峰值~GPU~显存增至~21~倍，接近~80~GB，逼近当前最先进~GPU（H100、A100）的上限。该攻击在黑盒设置下（攻击~3DGS~的扰动迁移到~Scaffold-GS）同样有效，其关键原因是~Poison-splat~的核心直觉——提升输入图像的非平滑度同时保持~3D~一致性——不依赖于特定~3DGS~变体的实现细节，而是对所有具有自适应复杂度的~3D~重建算法都有效。此外，论文还分析了朴素防御策略的局限性：直接限制高斯数量上限虽然可以缓解计算成本的增加，但会显著降低重建质量。因为~3DGS~无法自动区分必要的细节纹理和被投毒的纹理，会导致场景重建失败。

上述攻击从计算成本角度揭示了~3DGS~的脆弱性，而另一类攻击则瞄准重建管线中更上游的结构模块——不攻击~3DGS~的训练过程本身，而是攻击其上游依赖的~SfM~位姿估计。

从计算资源攻击转向重建质量攻击，研究视角进一步扩展到攻击共享基础设施的更具根本性的策略。

\paragraph{设计动机：从计算资源攻击转向重建质量攻击}

Poison-splat~攻击~3DGS~的自适应密度控制机制，将高斯数量膨胀为计算负担，但其破坏范围受限于特定~3DGS~变体的实现细节。PatchPoison（2026.04）~\cite{2604_13153}~则选择了一条更具根本性的攻击路径：瞄准~SfM~这一~3DGS、NeRF~等多种重建系统共享的前端管线。SfM~的输出——相机位姿和稀疏~3D~点云——是下游重建的基础设施：一旦位姿出错，无论下游使用何种重建方法，都无法恢复一致的场景表示。与~Poison-splat~利用~3DGS~的密度控制机制不同，PatchPoison~不需要访问或修改~3DGS~的训练过程，仅通过向输入图像注入不可察觉的扰动即可生效。这使得攻击门槛更低，也更难被检测。

\paragraph{攻击目标：SfM 特征匹配阶段}

现代~SfM~管线（如~COLMAP）的处理流程为：特征提取（SIFT~关键点）、特征匹配、几何验证（RANSAC，Random Sample Consensus）、光束平差（BA，Bundle Adjustment）。PatchPoison~瞄准的是特征匹配阶段——这是整个管线中对重复高频模式最敏感的环节。SIFT~在~DoG（Difference-of-Gaussian）金字塔中检测尺度空间极值作为关键点，而密集高频结构会在原生图像尺度上产生大量高对比度极值，从而生成大量具有强描述子响应的关键点。

\paragraph{棋盘格补丁的设计}

PatchPoison~的核心设计是在每张图像的固定位置（角落区域）注入小型棋盘格补丁。补丁由四个参数化维度完全确定：$(P, b, \alpha, c)$——补丁尺寸~$P$（像素）、块大小~$b$（控制空间频率）、混合强度~$\alpha \in [0,1]$、颜色对比度~$c$（明暗区域的灰度差）。

在空间频率方面，块大小~$b$~决定了棋盘格的频率特性。实验表明，极小块（$b=1,2$）生成的模式类似噪声，被~SIFT~的特征检测器忽略；而适中频率（$b=4$）产生最强的攻击效果，因为该尺度与~SIFT~原生图像尺度对齐，生成稳定、可重复的关键点——这些关键点在所有视角的相同像素坐标处被一致检测到，因此被特征匹配器以极高置信度匹配。

在尺寸方面，仅~$12\times12$~像素的补丁（约占~$800\times800$~图像宽度的~$1.5\%$）即可产生显著攻击效果。更大的补丁反而可能适得其反——当补丁过大时，3DGS~可能开始直接重建棋盘格图案本身，部分恢复重建质量。

\paragraph{可控混合强度~$\alpha$}

补丁通过~$\alpha$-混合融入原始图像，数学定义为：

\begin{equation}
I_i^* = I_i \cdot (1 - M_\alpha) + M_\alpha
\end{equation}

其中~$M_\alpha$~是嵌入在图像区域~$\Omega$~（角落）的补丁模板，在该区域外为零。$\alpha$~控制补丁的可见性：$\alpha=0$~时补丁完全透明（无攻击），$\alpha=1$~时补丁完全不透明。实验表明，极低的混合强度~$\alpha=0.1$（即补丁几乎不可见）已足以显著破坏重建，而继续增大~$\alpha$~对攻击效果的提升有限，却会降低视觉隐蔽性。这一设计确保了扰动在感知上的不可察觉性——毒化图像与原始图像的~SSIM~保持在高水平。

\paragraph{攻击机制：对极几何违反与~BA~不稳定}

攻击的因果链条如下：（1）棋盘格补丁在不同视角的相同像素坐标处产生大量类角点关键点；（2）这些关键点的~SIFT~描述子在跨视角间高度一致，被特征匹配器以高置信度配对；（3）然而，这些匹配对应关系并非来自任何真实~3D~场景结构——它们是纯粹由~2D~图案产生的伪对应；（4）因此，这些对应关系违反了对极几何约束（即同一~3D~点在不同视角的投影必须满足的基础矩阵约束）；（5）BA~尝试联合优化相机位姿和~3D~点位置以最小化重投影误差，但由于被迫调和本质上不相容的约束（来自真实场景结构的有效约束与来自补丁的虚假约束），优化问题变得病态（ill-conditioned），产生不稳定或严重错误的相机位姿估计；（6）下游~3DGS~使用这些错误位姿进行训练，由于几何一致性被破坏，无法恢复连贯的场景表示。

\paragraph{攻击效果}

在~NeRF-Synthetic~数据集上，仅~$12\times12$~像素的补丁即可将~SSIM~从~$0.969$~降至~$0.793$，LPIPS（Learned Perceptual Image Patch Similarity，衡量感知相似度的指标）~增加~6.8$\times$。攻击不需要修改任何~3DGS~代码或参数，且仅毒化~50\%~的训练数据即可维持有效攻击。在低毒化率（如~25\%）下，攻击效果方差较大，表明降质不完全一致。与高斯模糊、高斯噪声、几何变换、JPEG~压缩等~baseline~对比，PatchPoison~是唯一同时实现强重建降质与高视觉隐蔽性的方法——高斯模糊不影响重建，高斯噪声虽可降质但显著降低隐蔽性。局限性方面，PatchPoison~在纹理丰富的实景场景（Mip-NeRF~360）中效果较弱。因为棋盘格补丁容易融入自然背景，降低了对~SIFT~特征检测器的扰动效果。

PatchPoison~通过棋盘格补丁经验性地验证了攻击~SfM~特征匹配阶段的有效性，但其攻击设计缺乏理论解释——为什么棋盘格补丁能破坏匹配、破坏到什么程度会引发级联崩溃，均未给出形式化分析。PoInit-of-View（2026.04）~\cite{2604_16540}~将这一经验直觉发展为精准的模块化攻击理论，不仅提出可解释的攻击机制，还建立了从梯度不一致性到~SfM~全局崩溃的理论预测。

\paragraph{设计动机：从经验攻击到精准模块化理论}

与~PatchPoison~采用手工设计的棋盘格补丁不同，PoInit-of-View~从一个理论问题出发：能否设计一种攻击，使得~SfM~模块的崩溃是可预测、可解释的？该动机源自对~SfM~在~3D~重建管线中角色的深入分析。SfM~的输出——相机内参/外参和稀疏~3D~点云——是几乎所有下游重建方法（3DGS、NeRF、MVS，Multi-View Stereo）的共同基础。形式化地，给定输入视图~$\{I_i\}_{i=1}^N$，SfM~通过特征提取、描述子匹配、几何验证（RANSAC~配合对极约束~$\mathbf{x}'^\top F_{ij} \mathbf{x} = 0$）和光束平差（BA，最小化鲁棒重投影误差）联合优化全局相机位姿~$\{(R_i, \mathbf{t}_i)\}$~和~3D~点~$\{\mathbf{X}_j\}$。一旦这一初始化出错，下游优化器无法收敛到一致的场景表示——无论下游使用~3DGS~的高斯基元还是~NeRF~的体密度场。因此，攻击~SfM~等于攻击"共享基础设施"，攻击效果天然具有跨系统迁移能力。

\paragraph{核心机制：跨视角梯度不一致性}

PoInit-of-View~的攻击原理基于~SfM~的一个隐含假设：同一~3D~点在不同视角下的局部外观应近似一致。具体而言，对于同一~3D~点在两个干净视角~$(I_i, I_j)$~中的投影，其~Sobel~梯度应满足~$\|G(I_i(\mathbf{p}_i)) - G(I_j(\mathbf{p}_j))\|_1 \leq \tau_g$，即跨视角梯度一致性（CVGC，Cross-View Gradient Consistency）。这一性质是描述子可重复性的基础——SIFT~等描述子从局部梯度方向直方图中构造特征向量，梯度结构越相似，描述子越可能匹配成功。

PoInit-of-View~的攻击正是系统性地违反这一假设。通过引入跨视角梯度不一致性损失（CVIL，Cross-View Inconsistency Loss，$L_{\text{CVI}}$），对抗扰动被优化为在同一~3D~点的不同视角投影处产生矛盾的梯度信号：

\begin{equation}
L_{\text{CVI}} = \|G(R_i) - G(A_j)\|_1
\end{equation}

其中~$R_i$~是参考视图（保持干净），$A_j$~是被投毒视角。直觉上，当~$L_{\text{CVI}}$~超过一致性阈值~$\tau_g$~时，局部描述子的偏差将与梯度偏差成正比（由描述子的~Lipschitz~连续性保证）：$\|\psi(G(R_i)) - \psi(G(A_j))\|_2 \geq L_r^{-} \cdot \Delta$，进而使描述子匹配概率指数级下降。

\paragraph{理论预测：SfM 崩溃阈值}

PoInit-of-View~进一步给出了~SfM~全局崩溃的形式化条件。当关键视角对的跨视角不一致性超过阈值~$L_{\text{th}} = \tau_g + \delta_d / L_r^{-}$~时，匹配概率急剧下降，导致~RANSAC~无法为位姿图的关键边收集足够的内点——位姿图变为欠约束，全局~SfM~优化以高概率失败。论文在实验中经验估计了~$(\tau_g, \delta_d, L_r^{-})$~的值，验证了理论预测的崩溃阈值与实际观察一致。

\paragraph{SfM 模块级联失败的具体流程}

攻击引发的因果级联为：（1）\textbf{描述子偏差}——跨视角梯度不一致性迫使描述子偏差超出匹配距离阈值；（2）\textbf{特征匹配失败}——描述子偏差导致~SIFT~最近邻匹配的~Lowe~比率测试拒绝正确对应，匹配概率指数级下降（由描述子分布的亚指数尾部保证）；（3）\textbf{位姿估计错误}——由于内点数量不足，RANSAC~无法恢复可靠的相对位姿（本质矩阵~$E_{ij} = K_j^\top F_{ij} K_i$~估计失败），位姿图的关键边断裂；（4）\textbf{三角化失败}——缺乏有效匹配对意味着无法进行双视角三角化，稀疏~3D~点云严重稀疏或错误；（5）\textbf{下游重建崩溃}——无论使用~3DGS（以~SfM~稀疏点为高斯锚点和致密化起点）还是~NeRF/Instant-NGP（依赖准确位姿进行光度优化），错误的初始化都使下游优化无法收敛到一致的场景表示。

从~SfM~层面的攻击进一步扩展，研究者开始关注一种更隐蔽的威胁类别——后门攻击。

\paragraph{代理模型引导的优化}

由于~COLMAP~等~SfM~管线不可微分，PoInit-of-View~训练一个可微的~3DGS~代理模型~$P$~来近似跨视角几何行为。优化采用投影梯度上升（PGD），总损失函数为：

\begin{equation}
\max_{\tilde{I}} \, L_{\text{CVI}} - \lambda_{\text{SSIM}}(1-\text{SSIM}(\tilde{I}, I)) - \lambda_{\text{TV}} \text{TV}(\tilde{I}), \quad \text{s.t.} \, \|\tilde{I} - I\|_\infty \leq \epsilon
\end{equation}

其中~$L_{\text{CVI}}$~最大化跨视角不一致性，SSIM~和~TV~正则项确保扰动在感知上不可察觉。代理模型周期性刷新（每~$K=10$~次迭代）以跟踪投毒图像演化过程中的视角可见性变化。这一交替优化方案使投毒图像在视觉上几乎不可区分，但系统性地破坏了~SfM~初始化所需的多视角一致性。

\paragraph{攻击效果与跨系统迁移}

在~Tanks and Temples~数据集上，PoInit-of-View~使~SfM~的平均崩溃率超过~70\%（注册图像、三角化关键点和稀疏~3D~点均大幅下降），下游~PSNR~下降约~25\%，SSIM~下降~15--18\%，LPIPS~接近翻倍。攻击效果在~COLMAP、Instant-NGP~和~3DGS~三种下游管线上表现一致，验证了~SfM~层面攻击的下游无关性（downstream-agnostic）。在黑盒迁移设置下（以~3DGS~为代理优化，迁移到~NeRF），PoInit-of-View~相比单视角~baseline~在~PSNR~上超出~25.1\%，在~SSIM~上超出~16.5\%，证明了攻击~SfM~共享基础设施的强跨系统迁移能力。与~PatchPoison~相比，PoInit-of-View~的优势在于：攻击机制有理论支撑（崩溃阈值可预测），攻击范围更广（不仅限于~SIFT~对高频模式的敏感性），且扰动是优化产生的而非手工设计的，因此更难被针对性防御。

与上述在训练数据层面实施攻击不同，后门攻击直接在模型训练过程中植入恶意行为，属于更隐蔽的威胁类别。Poison-splat~投毒训练数据攻击计算资源，PatchPoison~和~PoInit-of-View~攻击~SfM~特征匹配破坏重建质量，但这些方法的攻击效果在训练完成后即固定——它们影响的是模型的被动属性（计算开销或重建精度）。GaussTrap（2025.04）~\cite{2504_20829} 首次系统研究~3DGS~后门威胁，将攻击范式从上述"一次性投毒"转向"按需激活"的后门攻击——在模型训练过程中植入仅在特定视角触发的恶意渲染行为。攻击目标不再是全局的计算资源或重建质量，而是推理阶段特定视角的场景语义。

\paragraph{威胁模型}

GaussTrap~的后门攻击场景假设：攻击者创建并分发植入了后门视角的~3DGS~模型~$G_{\text{poisoned}}$（例如通过域名仿冒攻击或恶意服务提供商），欺骗受害者下载使用。攻击者的双重目标是：（1）后门激活——当触发视角~$T$~作为输入时，被投毒模型渲染的图像~$I_{\text{poisoned}}(T)$~应与攻击者指定的恶意图像~$I_{\text{attack}}(T)$~高度一致，即~$I_{\text{poisoned}}(T) = I_{\text{attack}}(T)$；（2）功能保持——在所有非触发视角~$V \neq T$~下，被投毒模型的行为与干净模型完全一致：$I_{\text{poisoned}}(V) = I_{\text{clean}}(V)$。这种攻击在安全关键领域（自动驾驶、AR/VR、医疗模拟）中构成严重威胁。例如在自动驾驶场景中，攻击者可使模型在特定路口视角产生虚假障碍物或使交通标志消失的幻觉内容。

\paragraph{核心挑战：多视角一致性排斥}

GaussTrap~面临~Poison-splat~等投毒攻击不存在的独特挑战。3DGS~的渲染基于~3D~高斯的~2D~投影与~$\alpha$-混合，具有天然的多视角连续性——相邻视角的渲染结果应平滑过渡。当攻击者仅在单一视角注入恶意内容时，该视角与相邻正常视角之间缺乏有效的连续性学习，导致相邻视角出现明显伪影（artifacts）和大面积黑色云团。这一现象本质上是因为~3DGS~的多视角一致性机制会自然排斥与周围视角不协调的孤立恶意内容。这与~Poison-splat~中~3DGS~的自适应密度控制反而被攻击者利用的情况形成对比——多视角一致性在这里成为了攻击者的障碍而非助力。

\paragraph{三阶段优化管线}

为解决上述挑战，GaussTrap~提出三阶段优化管线，涉及三类视角的协同优化：攻击视角~$P_{\text{atk}}$（攻击者指定的触发视角）、稳定视角~$P_{\text{stab}}$（由~VES~生成的辅助视角）和正常视角~$P_{\text{train}}$（原始训练集中的视角）。三个阶段依次执行：（1）攻击阶段——在目标攻击视角~$P_{\text{atk}}$~注入恶意视角内容，优化高斯参数使渲染结果包含攻击者指定的幻觉内容（如路标消失、虚假障碍物、限速标志等），每轮~$T_a = 25$~次迭代中随机选择相机视角渲染并与目标图像计算损失；（2）稳定阶段——GaussTrap~的关键创新阶段，通过视角集成稳定（VES，View Ensemble Stabilization）缓解单一恶意视角导致的相邻视角伪影和黑色云团，每轮~$T_c = 5$~次迭代；（3）正常训练阶段——继续在正常视角~$P_{\text{train}}$~上优化场景，确保非目标视角的渲染质量不受影响，每轮~$T_t = 5$~次迭代。整个训练共~$E = 2500$~轮（epoch），每轮依次执行三个阶段。

\paragraph{视角集成稳定（VES）}

VES~是~GaussTrap~解决多视角一致性与后门隐蔽性矛盾的核心机制。朴素方案仅在攻击视角嵌入恶意内容并正常训练，会导致相邻视角出现严重伪影——如论文可视化结果所示，直接训练产生的渲染在攻击视角周围出现大面积黑色区域和失真。VES~的思路是：在攻击视角周围构建一组微小角度偏移的辅助视角（稳定视角），用干净~3DGS~模型渲染这些视角的~GT~图像并加入训练集，使攻击视角与正常视角之间存在渐进过渡而非突变。具体构建方式为：设攻击视角相机的外参旋转矩阵和平移向量为~$R_{w2c}$~和~$T_{w2c}$，角度扰动集合为~$A = \{(\delta_p^{(1)}, \delta_y^{(1)}), \ldots, (\delta_p^{(n)}, \delta_y^{(n)})\}$，其中~$\delta_p$~和~$\delta_y$~分别表示俯仰角（pitch）和偏航角（yaw）偏移。对每对偏移量，生成绕~X~轴和~Y~轴的基本旋转矩阵：

\begin{equation}
R_x(\delta_p) = \begin{bmatrix} 1 & 0 & 0 \\ 0 & \cos\delta_p & -\sin\delta_p \\ 0 & \sin\delta_p & \cos\delta_p \end{bmatrix}, \quad R_y(\delta_y) = \begin{bmatrix} \cos\delta_y & 0 & \sin\delta_y \\ 0 & 1 & 0 \\ -\sin\delta_y & 0 & \cos\delta_y \end{bmatrix}
\end{equation}

然后更新攻击视角相机的旋转分量：$R_{w2c}' = R_x(\delta_p) \cdot R_y(\delta_y) \cdot R_{w2c}$。偏移量枚举所有~$(\delta_p, \delta_y)$~组合（包括~$(0,0)$~即攻击视角本身），在合成数据集上设置为~$13^\circ$~和~$15^\circ$，在真实场景中设置为~$5^\circ$。多个微小偏移的辅助视角构成"视角集成"，模拟攻击视角周围的连续视点分布，有效缓解单一恶意视角导致的不稳定性，实现更鲁棒的重建结果。

\paragraph{损失函数}

三个阶段共用相同的损失函数形式，区别仅在于目标图像~$V_{\text{target}}$~的选取（攻击阶段使用恶意图像，稳定和正常阶段使用~GT~图像）：

\begin{equation}
\mathcal{L} = (1 - \lambda) \mathcal{L}_1(V', V_{\text{target}}) + \lambda \mathcal{L}_{\text{D-SSIM}}(V', V_{\text{target}})
\end{equation}

其中~$V' = R(G, P)$~为当前高斯分布在视角~$P$~下的渲染图像，$\mathcal{L}_1$~保证像素级精度，$\mathcal{L}_{\text{D-SSIM}}$~保证结构一致性，$\lambda$~控制两者权衡。这一损失函数的设计与~3DGS~标准训练一致，确保后门注入过程不会引入异常的训练信号模式。

\paragraph{攻击效果与多后门扩展}

在~Blender~合成数据集和~Mip-NeRF 360~真实场景数据集上的综合实验表明，GaussTrap~能够在攻击视角成功注入高相似度的恶意幻觉内容（如虚假限速标志、行人和车辆等），同时在训练视角、测试视角和稳定视角上均保持与干净模型相当的渲染质量。与~baseline~方法~IPA-NeRF+3DGS~相比，GaussTrap~在所有四种评估视角上均取得更优表现。此外，论文验证了多后门攻击的可行性：同时注入~$n = 2$~到~$n = 5$~个攻击视角时，渲染质量仅有轻微下降，表明~GaussTrap~具备多后门注入能力。消融实验进一步验证了~VES~的关键作用——使用~VES~后，攻击视角周围的伪影显著减少，训练和测试视角的图像质量指标均得到提升；增加~3DGS~致密化参数~$D$~可以提升整体指标，因为更高的致密化水平增强了模型表示复杂分布的能力，使其在各视角上更好地对齐~GT~数据的同时保留攻击图像的质量。

后门攻击通过三阶段管线植入恶意视角，而另一种更具隐蔽性的策略则是利用~3DGS~点云本身的稀疏性来隐藏攻击载荷。与~GaussTrap~需要从零开始三阶段训练不同，StealthAttack（2025.10）~\cite{2510_02314} 直接操纵已训练好的~3DGS~点云，无需重新训练。这意味着攻击者可以在模型部署后随时注入恶意内容，威胁时间线从"训练期"延伸至"全生命周期"。StealthAttack~的形式化攻击目标为：给定原始~3DGS~点云~$G$，寻找投毒后的点云~$G^*$，使得投毒视角~$v_p$~的渲染结果~$\mathcal{R}(G^*, v_p)$~包含攻击者指定的幻觉物体~$O_*$（即~$\|\mathcal{R}(G^*, v_p) - I_*\|_2^2$~最小化，其中~$I_*$~为~$G$~在~$v_p$~的原始渲染与~$O_*$~的叠加），同时所有非投毒视角~$v_k \neq v_p$~的渲染与干净模型保持一致~$\|\mathcal{R}(G^*, v_k) - \mathcal{R}(G, v_k)\|_2^2 \approx 0$。

StealthAttack~的设计动机源于对两种朴素方案的失败分析：（1）直接在目标视角的训练图像上叠加幻觉物体并重新训练——3DGS~的多视角一致性机制会将其视为噪声（违反一致性约束）并自动消除；（2）将幻觉物体直接反投影到高斯点云中——难以确定正确深度，可能被现有几何遮挡。由此~StealthAttack~得出关键结论：有效的~3DGS~投毒必须同时利用点云的显式表示特性和稀疏性结构，将幻觉物体放置在正常视角不可见的~3D~区域。

从后门攻击和点云操纵转向下游检测器攻击，攻击者的目标从破坏重建质量扩展到欺骗感知系统。

为此，StealthAttack~提出密度引导的点云攻击策略，包含场景空间分析、连续密度估计和最优位置选择三个步骤。在场景空间分析阶段，首先确定~3DGS~点云~$G$~的轴对齐包围盒（AABB），即找到所有高斯点坐标的最小/最大值，构建完全包围~3D~场景的矩形包围盒，然后将包围盒体素化为均匀网格~$S$。对每个体素~$s$，计算其密度~$\rho(s) = \sum_{g \in s} \alpha(g)$，即该体素内所有高斯点的不透明度之和。在连续密度估计阶段，基于逐体素密度~$\rho(s)$，使用高斯核密度估计（KDE）获得任意~3D~位置的连续密度值：

\begin{equation}
f(x) = \frac{1}{|S|}\sum_{s \in S} K_h(x - c(s))\rho(s), \quad K_h(x) = \frac{1}{(2\pi h^2)^{3/2}} \exp\left(-\frac{\|x\|^2}{2h^2}\right),
\end{equation}

其中~$c(s)$~为体素~$s$~的质心，$K_h$~为带宽~$h$~的高斯核函数。消融实验表明，中等带宽~$h = 7.5$~达到最优平衡——过小的带宽（$h = 0.1$）过度限制点放置范围，过大的带宽（$h = 10.0$）降低攻击精度。在最优位置选择阶段，从投毒视角的虚拟相机中心~$C$~通过幻觉物体覆盖的每个像素投射光线，方向为~$d$，沿光线~$x = C + td$~采样并查询~KDE~密度值，寻找密度最低的~3D~位置：

\begin{equation}
x_* = \arg\min_{x = C + td, \; t \in [t_{\min}, t_{\max}]} f(x),
\end{equation}

其中~$t_{\min} = 0.3$（排除相机附近的浮动点（floaters），这些点在~3DGS~优化中常见但不应作为攻击放置位置），$t_{\max}$~为原始场景在投毒视角每个像素处的深度值。低密度区域意味着从正常视角看不到这些位置——要么位于正常视角的视线盲区，要么被现有几何体遮挡——因此是隐藏幻觉物体的理想位置。在选定位置~$x_*$~处注入新的高斯点，其颜色直接从幻觉物体的对应像素获取，其余属性（协方差、不透明度等）设置为使幻觉在投毒视角清晰可见而从其他视角不可见。

然而，对于训练图像视角重叠度较高的场景，仅靠密度引导的点云攻击仍面临挑战——过多的正常视角覆盖使得低密度区域难以找到。为应对这一情况，StealthAttack~进一步提出视角一致性破坏攻击，通过向正常视角的训练图像注入可控高斯噪声（投毒视角保持干净），有意识地削弱~3DGS~的多视角一致性约束，从而更好地保留注入的幻觉。对训练图像~$I_k$~（视角~$v_k$），噪声注入方式为：

\begin{equation}
I'_k = \mathbb{1}_{v_k = v_p} \cdot I_k + \mathbb{1}_{v_k \neq v_p} \cdot \mathrm{clip}(I_k + \epsilon), \quad \epsilon \sim \mathcal{N}(0, \sigma_t^2),
\end{equation}

其中~$\mathrm{clip}(\cdot)$~确保像素值不超出有效范围，$\sigma_t$~为随训练迭代~$t$~衰减的噪声强度。噪声强度的设计原则是：训练早期使用强噪声以破坏多视角一致性（因为注入各训练图像的噪声相互独立，不同视角间的监督信号不再一致），后期逐步减弱以恢复正常视角的渲染质量。论文探索了三种衰减调度策略：线性衰减~$\sigma_t = \sigma_0(1 - t/T)$~均匀降低噪声；余弦衰减~$\sigma_t = \sigma_0 \cos(t\pi/2T)$~提供平滑的起始和加速的后期衰减；平方根衰减~$\sigma_t = \sigma_0\sqrt{1 - t/T}$~在更长时间内维持较高噪声后快速下降。实验表明，较高的初始噪声强度~$\sigma_0 = 100$~配合较慢的线性衰减策略取得最优平衡——攻击效果显著增强且对正常视角渲染质量的影响可控。

在~Mip-NeRF~360、Tanks and Temples~和~Free~三个数据集上的实验中，StealthAttack~在投毒视角的幻觉渲染质量（V-Illusory~的~PSNR、SSIM、LPIPS）上全面超越~IPA-NeRF（Nerfacto）、IPA-NeRF（Instant-NGP）和~IPA-Splat~三种~baseline，同时非投毒视角（V-Test）的渲染质量保持稳定。攻击成功的判定标准为~V-Illusory~的~PSNR~$> 25$~dB~且~V-Test~的~PSNR~下降~$\leq 3$~dB。论文还提出了基于~KDE~密度排序的攻击难度评估协议：利用相机内参和~10\%~采样半径计算各视角覆盖范围内的场景密度，按密度排序选择最小密度（$v_{\min}$，最易攻击）、中位密度（$v_{\text{mid}}$）和最大密度（$v_{\max}$，最难攻击）三个代表性视角，实验验证了场景覆盖密度与攻击成功率之间的负相关关系。在多视角攻击（同时在~2、3、4~个投毒视角注入幻觉）场景下，StealthAttack~同样优于所有~baseline。消融实验表明三个攻击组件（直接图像投毒、密度引导点云攻击、视角一致性破坏）中，单独的图像投毒几乎无效，密度引导点云攻击显著提升效果，三者组合取得最优表现。

\subsubsection{下游检测器攻击}

本路线介绍通过操纵~3DGS~场景来欺骗目标检测器、深度估计器等下游~AI~模型的攻击方法。重建质量攻击的目标是破坏~3DGS~本身的渲染结果，而另一类更具实际威胁的攻击则将目标锁定在下游感知系统——通过操纵~3DGS~场景来欺骗目标检测器、深度估计器等下游~AI~模型。这类攻击在自动驾驶、监控等安全关键场景中尤其危险。

3DGS~在自动驾驶等安全关键应用中的广泛部署引发了对抗攻击的严重安全关切。CLOAK（2025.06）~\cite{2506_00280}（Concealed Localized Object Attack Kinematics，隐蔽定位对象攻击运动学）首次利用~3DGS~视角相关外观实现了视角特异性伪装——其用视角分割的训练数据集替代均匀纹理的训练数据集——3DGS~标准~SH~训练管线保持不变。这是首个通过操纵训练数据来利用~3DGS~固有~SH~视角依赖性的攻击方法。

\paragraph{Baseline：3DGS 标准训练如何表示颜色。}
在标准~3DGS~训练中，每个高斯不存储固定~RGB~颜色，而是用球谐函数（SH）系数定义外观。SH~是一种用球面基函数逼近视角相关颜色的数学工具——直觉上，可以把~SH~系数想象成一组"滤镜模板"，当观察者移动到不同位置时，不同的模板被激活并混合，产生不同的颜色。训练过程中，来自不同相机视角的图像通过可微渲染反向传播，共同优化同一组~SH~系数，使高斯在所有训练视角下都呈现正确的外观。

\paragraph{改了什么：从均匀纹理到视角分割对抗数据集。}
CLOAK~不修改~3DGS~管线本身，而是操纵训练数据。具体地，攻击者首先选定一组目标相机视角~$C^* \subset C$，然后为这些视角生成对抗图像——将目标物体的纹理替换为对抗纹理（如路面纹理、停止标志等）。形式化地，原始训练数据集~$\mathcal{D} = \{(x_i, c_i)\}_{i=1}^N$~被替换为攻击数据集~$\mathcal{D}' = \{(A(x_i, c_i), c_i)\}_{i=1}^N$，其中视角选择性替换函数定义为：

\begin{equation}
    A(x, c) =
    \begin{cases}
        \tilde{x}, & \text{if } c \in C^*, \\
        x, & \text{otherwise},
    \end{cases}
\end{equation}

即仅对落入目标视角集合~$C^*$~的图像替换为对抗版本~$\tilde{x}$，其余图像保持不变。替换后，3DGS~照常执行标准训练，SH~系数在对抗视角和正常视角的"推拉"优化中自然编码出视角特异性的外观。

\paragraph{设计动机。}
这种训练数据层面攻击的设计动机有三层考量：（1）可行性——攻击者无需访问~3DGS~模型参数或训练过程，只需在数据采集阶段替换少量图像即可实施，攻击门槛极低；（2）隐蔽性——从非目标视角观察时，场景完全正常，连人类检查员也无法发现异常，因为~SH~视角依赖性本身就是~3DGS~的标准特性，不是漏洞特征；（3）鲁棒性——对抗外观被编码进~SH~系数后成为场景表示的一部分，任何使用该场景文件的下游应用都会受到影响，且源图像通常不随场景文件分发，事后追溯困难。

从基于训练数据操纵的~CLOAK~到基于参数空间直接操作的~DAGGER，下游检测器攻击的路线开始分化。

\paragraph{实验效果。}
在实验中，CLOAK~使用~Blender Cycles~渲染器从~210~个相机角度渲染一辆汽车，覆盖半球形区域实现完整~360°~可视化，嵌入三组外观：110~个正常角度显示真实汽车外观，80~个俯视角度显示"路面"纹理，20~个正后方角度显示"停止标志"纹理。训练完成后，对抗纹理在目标视角（俯视和正后方）清晰呈现，但对角线视角会模糊这些对抗修改，使人类观察者和检测器都难以从非目标视角发现外观不一致。在黑盒评估中，相机视角从正常地面角度平滑过渡到对抗俯视和后方角度，YOLOv8~在~110~个正常视角中成功检测到汽车~80~次，但在~80~个对抗俯视视角中仅有~2~次检测成功（78/80~次漏检，失败率~97.5\%），且在对抗纹理完全可见时出现完全漏检，证明了视角特异性伪装对下游检测器的严重威胁。

CLOAK~利用~3DGS~固有的~SH~视角依赖性实现伪装，而~DAGGER（2025.06）~\cite{2506_00280}~则走了一条不同的路线——直接在参数空间中操作。DAGGER~采用定向对抗攻击策略，在无需访问训练数据的前提下直接对~3D~高斯参数施加扰动——通过投影梯度下降精确定位并修改~3D~高斯的位置、颜色或不透明度属性，使渲染图像中的对抗扰动欺骗~Faster R-CNN~等多阶段目标检测器的区域提议和分类阶段。与~CLOAK~利用~SH~视角依赖性不同，DAGGER~通过对高斯参数的微小修改在所有渲染视角中产生一致的对抗效果，揭示了~3DGS~参数空间直接操纵对下游检测器的广泛威胁。

上述白盒攻击方法需要访问模型参数，在实际场景中限制较大。ComplicitSplat（2025.08）~\cite{2508_11854} 首次实现~3DGS~对下游检测器的黑盒攻击。与~CLOAK~（同样操纵训练数据利用~3DGS~固有~SH~视角依赖性）的相似之处在于同样操纵训练数据，但~ComplicitSplat~的核心区别在于完全的黑盒设定。

\paragraph{改了什么：从白盒到黑盒。}
CLOAK~假设攻击者了解~3DGS~训练流程但无需修改管线，而~ComplicitSplat~进一步收紧威胁模型——攻击者仅知道~3DGS~使用高斯基元和~SH~系数表示场景（这是公开知识），但不访问内部场景参数、中间表示、最终训练结果或任何下游模型。攻击者的唯一能力是操纵训练图像：替换或注入训练集中的图像，但不修改~3DGS~算法本身、损失函数或优化过程。这使得攻击门槛极低——只需在数据采集阶段替换图像即可。

从纯黑盒的~ComplicitSplat~过渡到物理世界攻击，研究者的关注点从数字域扩展到物理可实现性。

\paragraph{技术实现：多区域视角分割对抗数据集。}
与~CLOAK~选定一组目标相机视角~$C^* \subset C$~不同，ComplicitSplat~定义~$n$~个目标视角区域（每个区域对应一种隐蔽外观）。形式化地，第~$i$~个目标视角区域由参考视角~$c_{ref,i}$~和角度阈值~$\theta_i$~定义：

\begin{equation}
    R_i^* = \{c \in C : \angle(c, c_{ref,i}) < \theta_i\}, \quad i=1,\ldots,n,
\end{equation}

其中~$\angle(c, c_{ref,i})$~表示相机光轴方向之间的角度距离。对落入每个区域的图像，攻击者将目标物体的外观替换为对抗伪装纹理（如路面、草地），生成对抗图像~$\tilde{x}^{(j)}$，非目标区域图像保持不变。最终形成攻击数据集~$\mathcal{D}' = \{(A(x_i, c_i), c_i)\}_{i=1}^N$，其中视角选择性替换函数为~$A(x,c) = \tilde{x}^{(j)}$~当~$c \in R_j^*$，否则~$A(x,c) = x$。该算法对~$N$~张图像仅遍历一次，时间和空间复杂度均为~$O(N)$。

\paragraph{SH 推拉优化机制。}
3DGS~标准训练使用~$\ell=2$~阶~SH（每颜色通道~9~个系数，RGB~共~27~个系数），这些系数参数化一个连续的方向颜色函数，在训练中通过多视角图像的可微渲染反向传播共同优化。ComplicitSplat~的关键洞察是：当训练数据同时包含正常视角图像和对抗视角图像时，SH~系数在不同视角图像间的"推拉"优化过程自然将多视角对抗外观编码进场景表示——对抗视角的梯度推动~SH~系数朝对抗外观方向优化，而正常视角的梯度则保持原有外观，两者的竞争最终使同一组高斯在不同观察角度呈现截然不同的外观。例如，使车辆从地面视角完全正常，而从俯视角度呈现路面纹理从而对空中检测器隐形。这种对抗外观的嵌入仅受~SH~表示容量的限制，理论上可以编码任意多种隐蔽外观。

\paragraph{对训练收敛速度的影响。}
由于对抗视角和正常视角的外观竞争导致~SH~优化中的~SSIM~等图像相似度指标被反复"推拉"，训练收敛速度约为正常训练的~2$\times$~慢（在~NVIDIA RTX 4090~上，0.5~分钟对比~15~秒达到~SSIM~$>0.95$）。然而在实际中，3DGS~场景通常训练至~30K~次迭代，ComplicitSplat~的场景在此时仍能达到~SSIM~$0.93$~的高渲染质量，视觉效果高度逼真。

\paragraph{设计动机：静默篡改的攻击面。}
设计动机基于~3DGS~生态的现实特征：3DGS~场景广泛公开共享（如~SceneSplat-7K~等数据集提供大规模场景下载），但源图像通常不随场景文件分发。这意味着攻击者在数据采集阶段替换少量图像后，最终~3DGS~场景文件中不会保留任何篡改痕迹——事后检查场景文件无法发现异常，而源图像又不可获取以进行交叉验证。这种"静默篡改"使~3DGS~场景成为不知情的共谋者。此外，由于攻击不依赖任何特定模型参数，具有天然的跨检测器泛化能力。

\paragraph{跨检测器泛化实验验证。}
实验中，ComplicitSplat~成功攻击覆盖三种检测架构的~6~个检测器：单阶段检测器~YOLO（v3/v5/v8/v11）、多阶段检测器~Faster R-CNN~和~Transformer~检测器~DETR。攻击成功率（ASR）从~50\%（DETR）到~91\%（YOLOv8），证明黑盒攻击无需访问下游模型即可广泛生效。值得注意的是，较早的~YOLOv3/v5~反而比~YOLOv8~更具鲁棒性，部分先前评估也表明~YOLOv5~在车辆检测任务上可能优于~YOLOv8。在~AP@0.5~指标上，所有检测器均出现显著下降——草地伪装比路面伪装造成更大的性能退化（如~YOLOv11~在草地伪装下~AP@0.5~下降~0.41，路面伪装下下降~0.32），停止标志上的伪装效果更为突出（时钟纹理使~YOLOv11~的~AP@0.5~下降高达~0.52）。

\paragraph{真实世界物理捕获场景。}
为验证攻击在真实~3DGS~工作流中的有效性，作者使用低成本工具将合成数据方法扩展到物理世界：准备两个相同的物理车模（一个正常蓝色、一个涂有路面伪装），分别用~PolyCam~手机应用（基于摄影测量技术）扫描生成~3D~网格，导入~Blender~后从完全相同的相机位姿分别渲染——地面视角渲染正常车模、俯视视角渲染伪装车模，合并后的图像集训练~3DGS~场景。结果中，Faster R-CNN~的~ASR~高达~98.51\%，证明多阶段检测器在物理部署中可能特别脆弱；YOLOv8~和~DETR~分别保持~58.82\%~和~68.75\%~的中等成功率；YOLOv3~和~YOLOv11~的迁移性较低（18.60\%~和~37.86\%）。这些发现表明~3DGS~攻击可以超越数字渲染泛化到真实物理捕获，但模型架构在真实世界鲁棒性中起重要作用。

前面的攻击方法在数字域中操纵~3DGS~场景，而更具挑战性的是将对抗攻击迁移到物理世界——生成可以被实际制造和部署的~3D~对抗物体。与~CLOAK~在数字域操纵训练数据实现伪装不同，3DGAA（2025.07）~\cite{2507_09993} 的核心目标是生成物理可实现的对抗物体——如路边的一个完整~3D~假路障或涂装了对抗纹理的微型车辆，从任意视角对检测器造成干扰。

从3DGAA的全参数优化到PGA的纯外观迷彩，再到BiTAA的双任务联合攻击，物理世界攻击的路线逐步精细化。

\paragraph{改了什么：从数字域训练数据操纵到物理世界全参数优化。}
CLOAK~保持~3DGS~标准训练管线不变，仅通过视角分割训练数据利用~SH~视角依赖性实现数字域伪装。3DGAA~则直接在~3DGS~的全~14~维参数空间中操作——每个高斯基元~$g_j \in \mathbb{R}^{14}$~由位置~$x \in \mathbb{R}^3$、旋转~$q \in \mathbb{R}^4$（四元数）、缩放~$s \in \mathbb{R}^3$、颜色~$c \in \mathbb{R}^3$~和不透明度~$\alpha \in \mathbb{R}$~参数化，联合优化几何和外观属性。这比传统~2D~对抗补丁（只能贴在平面上）和基于~mesh~的~3D~攻击（受限于简化几何）提供了更高自由度的扰动空间。对抗目标形式化为
\begin{equation}
    \min_{G} \; L_{adv} = \mathbb{E}_{P \sim \mathcal{V}} \left[ f_D\bigl(R(G, P)\bigr) \right],
\end{equation}

其中~$R(G, P)$~是从相机位姿~$P$~渲染的图像，$f_D(\cdot)$~是预训练检测器~$D$~对目标类别的置信度分数，$\mathcal{V}$~是视角分布。直觉上，优化目标是找到一组高斯参数，使得从任何视角渲染出的图像都让检测器"看不到"目标。

\paragraph{四阶段管线设计。}
3DGAA~采用四阶段管线生成物理鲁棒的对抗物体：（1）\textbf{3DGS~生成}——从~$n$~张多视角~RGB~图像和标定相机位姿~$\{P_i\}_{i=1}^n$~出发，经预训练~3DGS~网络生成初始高斯物体~$G = \{g_j\}_{j=1}^N$；（2）\textbf{物理过滤}——通过拓扑剪枝和结构去噪消除伪影，确保多视角渲染的几何一致性（详见下文）；（3）\textbf{对抗优化}——在对抗损失和形状一致性损失之间动态平衡权重，迭代更新高斯参数；（4）\textbf{物理增强}——四步可微变换链模拟部署时的环境变化，提升从仿真到现实的迁移能力（详见下文）。

\paragraph{物理过滤：拓扑剪枝与结构去噪。}
直接在整个~3DGS~参数空间上优化往往产生孤立浮动物体和空间抖动等伪影，这些缺陷在多视角渲染下尤为突出，且物理不可实现。物理过滤模块分两阶段消除这些问题。

\textbf{拓扑剪枝（TP）}：对每个高斯~$g_j$，通过~$k$~近邻估计其在~3D~空间中的局部密度~$\rho(g_j)$，并将密度低于~$p$~百分位数阈值~$\tau_d = \text{Percentile}(\{\rho(g_j)\}_{j=1}^N, p)$~的基元视为离群点移除。这种基于百分位数的方法自动适应不同物体密度，无需手动调节参数。

\textbf{结构去噪（SD）}：在剪枝后的高斯集合上，通过相机感知的各向异性滤波平滑尺度和不透明度的突变。具体地，滤波强度~$\omega_j$~由每个高斯到最近相机平面的投影距离调制
\begin{equation}
    \omega_j \propto \min_i \|R_i x_j + T_i\|_2 / f_i,
\end{equation}

其中~$R_i, T_i, f_i$~分别表示第~$i$~个相机的旋转矩阵、平移向量和焦距。这种空间自适应机制确保远距离区域得到更强的平滑，而近距离细节得到保留。滤波后的尺度更新为~$s_j' = \sqrt{\omega_j^2 + s_j^2 \cdot I}$。两阶段过滤共同消除了~83.1\%~的伪影伪结构，为后续对抗优化提供几何可信的基础。

\paragraph{对抗损失与形状一致性的动态平衡。}
对抗优化的核心挑战是同时追求强攻击效果和物理真实感。为此，3DGAA~引入形状一致性损失，约束优化过程中的几何偏移
\begin{equation}
    L_{shape} = \frac{1}{N} \sum_{j=1}^N \left( \|p_j - p_{j,0}\|_2^2 + \|s_j - s_{j,0}\|_2^2 + \|q_j \cdot q_{j,0}^{-1}\|_2^2 \right),
\end{equation}

其中~$p_{j,0}, s_{j,0}, q_{j,0}$~是优化前的初始位置、缩放和旋转，$\cdot$~表示四元数乘法。总损失为~$L_{total} = \lambda_{adv} L_{adv} + \lambda_{shape} L_{shape}$，其中权重通过动态机制自适应调整：先将~$L_{adv}$~除以缩放因子~$\gamma$（实验中设为~10.0）归一化，然后计算~$w_{adv} = 1/(\hat{L}_{adv} + \epsilon)$，$w_{shape} = L_{shape}$，归一化后得到~$\lambda_{adv}$~和~$\lambda_{shape}$，并钳位~$\lambda_{shape} \geq 0.4$~防止形状约束坍塌。这一设计的直觉是：当对抗损失已经很低（攻击已足够强）时，$w_{adv}$~变大，优化自动转向保持几何真实感；反之则加大对攻击的追求。该动态方案消除了手动调节固定损失权重的需要。

\paragraph{四步可微物理增强变换链。}
为提升对抗物体从仿真到现实的迁移能力，3DGAA~在优化过程中对渲染图像施加四步可微变换链~$T = T_{occl} \circ T_{shadow} \circ T_{photo} \circ T_{noise}$（按从右到左的顺序依次应用）：（1）\textbf{成像退化}~$T_{noise}$——模拟传感器噪声，噪声水平随物体距离增加而增大，$T_{noise}(I) = I + \mathcal{N}(0, \sigma^2(d))$，其中~$\sigma(d) = \sigma_0 + \beta \cdot d$，$d$~为像素级深度值，建模了低质量硬件或大气散射导致的退化；（2）\textbf{光度变化}~$T_{photo}$——通过通道级仿射变换模拟光照不一致、曝光漂移和白平衡变化，$T_{photo}(I_c) = I_c \cdot \gamma_c + \beta_c$，其中~$\gamma_c \sim \mathcal{U}(0.9, 1.1)$，$\beta_c \sim \mathcal{U}(-0.05, 0.05)$；（3）\textbf{阴影投射}~$T_{shadow}$——利用深度引导的~sigmoid~掩码模拟自然光照下的软阴影，$T_{shadow}(x,y) = 1 / (1 + e^{-\kappa(d(x,y) - d_{th})})$，其中~$\kappa$~控制阴影柔和度，$d_{th}$~为阴影投射的深度阈值；（4）\textbf{自适应遮挡}~$T_{occl}$——引入随机透明度的矩形遮挡器模拟现实中的前景杂物和部分遮挡。所有变换均可微分，梯度通过增强管线回传至高斯参数，使对抗优化自然产出对物理环境变化鲁棒的扰动。

\paragraph{三种选择性维度优化模式。}
实际部署场景对物理可变性有不同约束。例如刚性车身不允许形变，而~3D~打印物体则可以自由改变几何。为此，3DGAA~将~14~维参数分为外观子集~$\Theta_{app} = \{\text{rgb}, \alpha\}$（4~维）和几何子集~$\Theta_{geo} = \{\text{pos}, s, q\}$（10~维），提供三种优化模式：（1）\textbf{外观模式}~$3DGAA^{app}$——冻结~$\Theta_{geo}$，仅优化~$\Theta_{app}$。生成的对抗纹理可直接作为可打印的涂装或薄膜施加于车辆表面，无需任何几何修改。在所有攻击方法中达到最佳感知保真度（最低~LPIPS~和最高~PSNR）；（2）\textbf{几何模式}~$3DGAA^{geo}$——冻结~$\Theta_{app}$，仅优化~$\Theta_{geo}$（受~$L_{shape}$~约束）。消融实验表明几何维度（位置~$x$~和缩放~$s$）在~30~个~epoch~内即可快速收敛至高~LCR，揭示了几何扰动对检测器输出的高敏感性；（3）\textbf{全维度模式}~3DGAA——联合优化~$\Theta_{app} \cup \Theta_{geo}$，提供最强攻击效果（最高~LCR、最低~mAP），代价是相较外观模式有轻微的真实感损失。维度敏感性分析进一步揭示：位置~$x$~和缩放~$s$~是最敏感的参数，贡献了早期攻击进展的大部分；不透明度~$\alpha$~和颜色~$c$~呈中等敏感性；旋转~$q$~由于在~3DGS~表示中的非局部和冗余性质，影响最小。

\paragraph{实验效果。}
3DGAA~在~CARLA~模拟器的~20~种车型、5~种天气条件下评估，覆盖全方位~24~个方位角/仰角和三种距离（3m、5m、10m）。在六种主流检测器（Faster R-CNN、Mask R-CNN、SSD、YOLOv3/v5/v8）上，3DGAA~将检测~mAP@0.5~从~87.21\%~降至~7.38\%，LCR（Log Confidence Reduction，对数置信度缩减，衡量检测器对正确标签的置信度下降比例）达~3.5628，显著超过最强基线~RAUCA（LCR~为~1.4705）。外观模式~$3DGAA^{app}$~和几何模式~$3DGAA^{geo}$~也分别达到~3.4410/8.03\%~和~2.6766/13.64\%~的~LCR/mAP，均超越所有基线方法。跨检测器迁移实验表明，在~Faster R-CNN~上优化的扰动迁移至~YOLOv3~时~LCR~达~6.180，证明~3DGAA~通过多维优化捕获了跨架构的可迁移漏洞。在物理验证中，对抗纹理经~UV~展开后打印在哑光贴膜上，贴附于同一刚性微型车辆外壳（无车身修改），在室内/室外、日间/夜间四种光照条件下，3DGAA~将~AP@0.5~从~84.17\%--96.25\%~降至~17.92\%--35.83\%。包含物理增强的训练方案在所有光照条件下均稳定诱导漏检或误分类，验证了仿真到现实迁移的有效性。

3DGAA~展示了~3DGS~在物理对抗攻击中的潜力，而~PGA~和~R-PGA~则进一步解决了物理对抗迷彩在真实环境中的实用性问题。与~3DGAA~（联合优化~14D~参数的全维度攻击）不同，PGA（2025.07）~\cite{2507_01367} 仅优化~SH~系数（表示表面颜色）而冻结几何参数~$\{\mu_g, s_g, q_g, \alpha_g\}$，确保物理部署时无需改变物体形状。对抗迷彩可直接作为可打印的涂装或薄膜施加于车辆表面。

\paragraph{改了什么：从全参数优化到纯外观迷彩生成。}
3DGAA~在~14~维参数空间中联合优化几何和外观属性以追求最强攻击效果。PGA~则采取了更务实的路线：锁定所有几何参数，仅优化~SH~系数~$k_g$~的零阶项~$k_{g,0}$~（即与视角无关的基础颜色分量）。这一选择虽然牺牲了部分攻击强度，但换来了关键的实际部署优势——无需修改物体形状即可将对抗迷彩作为可打印涂装部署。此外，3DGAA~依赖~CARLA~模拟器构建场景，而~PGA~利用~3DGS~的快速重建能力直接从真实照片建模，消除了仿真器与真实世界之间的~domain gap（域差异）。

\paragraph{三模块攻击框架。}
PGA~的核心框架包含三个模块，形成从场景重建到对抗优化的完整管线：（1）\textbf{重建模块}——从多视角图像~$I = \{I_1, I_2, \ldots\}$~出发，通过~3DGS~训练框架重建场景高斯~$G = \{g_1, g_2, \ldots, g_N\}$，无需手动构建~3D~模型或依赖模拟器；（2）\textbf{渲染模块}——在目标物体周围选择多视角相机~$\Theta = \{\theta_{c1}, \theta_{c2}, \ldots\}$（变化距离、俯仰角和方位角确保全面覆盖），通过~3DGS~光栅化器渲染图像~$I_r = R(\theta_c, G)$，然后用~SAM~从渲染图像中提取目标物体掩码~$M = S(I_r, P)$（其中~$P$~为目标物体的提示词）。掩码将扰动约束在目标区域内：保留渲染图像副本~$I_{ori}$，最终送入检测器的图像为~$I_{det} = (I_r \odot M) + (I_{ori} \odot (1 - M))$，其中~$\odot$~为逐元素乘法；（3）\textbf{攻击模块}——将~$I_{det}$~输入受害检测器~$F$~获得检测结果~$B = F(I_{det}; \theta_f)$，计算检测损失，并通过反向传播迭代更新~SH~系数。检测损失定义为
\begin{equation}
    L_{det}(I_{det}) = -\sum_{I} \max_{m^*} \left( c_{m^*}^{(I)} \right), \quad m^* = \arg\max_m \mathrm{IoU}(b_m^{(I)}, b_{gt}),
\end{equation}

即在每个输入~$I$~中找到与真实标注~$b_{gt}$~的~IoU~最大的检测框~$m^*$，最小化其正确类别置信度。SH~系数通过~I-FGSM~迭代更新~$k^{t+1} = k^t + \alpha \cdot \nabla_{k} L_{det}(I_{det})$，攻击完成后通过~SuGaR~从优化后的高斯~$G'$~提取迷彩网格~$\mathcal{C}$~用于物理部署。

\paragraph{跨视角一致性：防止高斯间的互相遮挡与自遮挡。}
PGA~需要解决一个~3DGAA~未遇到的特殊挑战：由于~PGA~冻结了所有几何参数，高斯之间的空间关系无法通过优化调整，必须通过重建阶段的预处理来保证多视角一致性。具体存在两类遮挡问题：（1）\textbf{互相遮挡}——~3DGS~重建时并非所有高斯都精确贴合物体表面，部分高斯位于物体内部或表面之上方。当视角变化时，这些位置不准的高斯会遮挡表面高斯的颜色。PGA~引入~SuGaR~的正则化项，将高斯对齐到物体表面并鼓励降低不透明度，防止高斯被优化到物体内部，确保视角变化时表面颜色不被其他高斯遮挡；（2）\textbf{自遮挡}——高阶~SH~系数赋予高斯强大的视角依赖颜色表示能力，使得同一高斯的不同部分在不同视角下呈现截然不同的颜色。在多视角联合优化中，优化器倾向于精细调整每个视角下高斯的可见部分，导致严重的局部颜色不一致。PGA~仅优化~SH~零阶项~$k_{g,0}$~（视角无关的基础颜色），确保每个高斯表面的颜色变化均匀——同一高斯在所有视角下呈现相同的对抗颜色，从根本上消除了自遮挡问题。

PGA~解决了静态光照条件下的问题，但真实世界的光照和天气是动态变化的。

\paragraph{Min-max 优化：抗背景变化的鲁棒迷彩。}
PGA~的另一个核心挑战是：训练图像的背景多样性有限，传统优化容易陷入局部最优——生成的迷彩只在训练背景附近有效，视角稍有变化即失效。PGA~将此问题建模为~min-max~优化
\begin{equation}
    G' = \min_{G} \max_{\delta} L_{det}\bigl(I_{det}(\theta_c, G) + \delta \odot (1 - M)\bigr), \quad \text{s.t.} \; \|\delta\|_\infty \leq \epsilon,
\end{equation}

其中~$\delta$~为背景像素上的逐像素噪声，$\epsilon$~为噪声预算。具体执行过程为：在每次迷彩优化迭代之前，先对背景添加~I-FGSM~优化的对抗噪声以最大化检测损失（即通过对抗攻击强化背景中的检测线索），然后优化迷彩以最小化该增强后的检测损失。内层~max~使对抗特征对背景变化具有鲁棒性，外层~min~则在此增强条件下追求最强攻击效果。当检测器恢复正确检测或达到迭代上限时停止背景扰动，避免过度干扰使迷彩优化过于困难。

\paragraph{辅助正则化与物理增强。}
除上述核心技术外，PGA~还采用多项辅助技术提升物理域效果：（1）\textbf{EoT~变换}——在优化过程中随机化尺度、对比度、亮度并添加噪声，增强物理鲁棒性；（2）\textbf{不可打印分数（NPS）}——$\mathrm{NPS} = \sum_{\{p\} \in C(I_{det})} \min_{p' \in P} |\{p\} - p'|$，其中~$P$~为可打印颜色集，$C(I_{det})$~为迷彩使用的~RGB~三元组，最小化打印失真；（3）\textbf{主颜色正则化}——从训练集背景像素中提取~K-means~聚类的主颜色，约束迷彩颜色接近背景主色调~$L_{clr} = \frac{1}{|\Omega|} \sum_{(x,y) \in \Omega} \min_i \|I_{det}(x,y) - \mu_i\|_2$（其中~$\Omega$~为掩码非零像素位置集），提升视觉隐蔽性；（4）\textbf{SH~系数扰动约束}——限制攻击前后~SH~零阶系数的~$L_2$~距离。总损失整合为~$L_{total} = L_{det}(T(\cdot)) + \lambda(\mathrm{NPS} + L_{clr} + \|k_{g,0} - k_{g,0}^*\|_2)$，其中~$T$~为~EoT~变换集，$k_{g,0}^*$~为攻击前的初始值。迭代更新公式相应修正为~$k_{g,0}^{t+1} = k_{g,0}^t + \alpha \cdot \nabla_{k_{g,0}} L_{total}$。

\paragraph{实验效果。}
PGA~在数字域和物理域均进行了系统评估。数字域实验使用~CARLA~模拟器（用于与基线公平对比），覆盖两种天气（晴天/阴天）、四种距离（5m/10m/15m/20m）、五种俯仰角（$20^\circ$--$60^\circ$）和~$360^\circ$~全方位（$10^\circ$~间隔），共~1440~张测试图像。以~Faster R-CNN~为白盒目标、在部分覆盖（partial~coverage，更具挑战性且更实用）设定下，PGA~在所有设定下均达到最优攻击效果，显著超越~DAS、FCA、DTA、ACTIVE、TAS~和~RAUCA~六种基线方法。跨检测器迁移实验中，迷彩从~Faster R-CNN~迁移至~YOLOv5、Mask R-CNN~和~Deformable-DETR~均保持强攻击性能。消融实验验证了两个核心技术的必要性：跨视角一致性技术和~min-max~优化框架联合使用时达到最佳攻击效果。物理域实验在两个尺度上验证：（1）\textbf{1:24~缩比模型}——在~Audi Q5~缩比车上部署迷彩，50cm~和~100cm~距离多视角拍摄，多个检测器评估均显示~PGA~显著优于其他方法；（2）\textbf{1:1~真实车辆}——在~GOLF Sportsvan~真车上部署，使用无人机环绕拍摄~282~张图像输入~PGA~框架生成迷彩（部署时仅需可打印标定贴辅助~SAM~分割），以贴纸形式部署，Faster R-CNN~的~AP@0.5~从~88.48\%~降至~25.67\%。这表明仅需简单工具（相机和打印贴纸），PGA~即可对真实车辆实施有效攻击。

PGA~的对抗迷彩在优化时固定环境光照，无法适应部署时从正午到黄昏、从晴天到雾霾的连续光照变化；同时，标准~3DGS~将光照烘焙到视角依赖的~SH~系数中，导致跨视角纹理不一致——特定视角的迷彩呈现不同外观，阻碍收敛到统一的物理纹理。R-PGA（2026.03）~\cite{2603_26067} 作为~PGA~的期刊扩展版本，从场景仿真和优化目标两个维度系统解决了这些局限。

在场景仿真层面，R-PGA~首先对~3DGS~的每个高斯基元~$g_j$~引入显式的物理渲染（PBR）属性：反照率~$a \in [0,1]^3$（本质表面颜色，独立于光照）、金属度~$m \in [0,1]$、法线~$n$~和粗糙度~$r \in [0,1]$，建立物理解耦表示。渲染时，标准~3DGS~的颜色渲染被替换为基于渲染方程的物理着色模型：出射辐射度~$L_o(\omega_o, x) = \int_\Omega f_r(\omega_i, \omega_o, x)\, L_i(\omega_i, x)\, (\omega_i \cdot n)\, d\omega_i$，其中~$f_r$~采用~Disney~微面元~BRDF~模型，入射光~$L_i$~分解为直接光照（基于图像的照明~IBL，用高分辨率环境图表示）和间接光照（球谐函数捕获多次反弹，配合可见性项处理遮挡），法线~$n$~从高斯协方差矩阵的最短轴推导。关键在于：优化反照率~$a$~而非~SH~系数，使对抗模式反映表面的本质颜色而非瞬时照明效果，从根本上消除了~PGA~中~SH~烘焙光照导致的跨视角纹理不一致问题。

然而，对整个场景进行材质分解在开放环境中是不适定的——周边区域视角监督稀疏，导致几何误差、材质/光照分解失败和悬浮伪影相互加剧。R-PGA~因此设计了混合渲染管线：（1）\textbf{前景渲染}——用~SAM~从多视角图像中提取目标物体掩码~$M_{obj}$，仅用物体图像~$I_{gt}^{obj}$~重建目标物体的可重照明~3DGS~$G_{obj}$，渲染公式为~$I_{fg} = R(G_{obj}, c) \cdot M_{tar} + I_{gt}^{obj} \cdot (M_{obj} - M_{tar})$，其中~$M_{tar}$~为迷彩区域掩码；（2）\textbf{背景渲染}——用预训练的流匹配图像翻译模型~LBM，以目标环境图~$E_t$、原始图像~$I_{gt}$~和前景渲染的光度线索~$I_{fg}$~为条件，单步推理生成与前景光照一致的背景~$I_{bg} = H(I_{fg}, I_{gt}, E_t)$；（3）\textbf{场景合成}——通过掩码混合~$I_{det} = M_{obj} \cdot I_{fg} + (1 - M_{obj}) \cdot I_{bg}$~得到完整渲染图像。这种生成式方法绕过了开放场景背景材质分解的不适定问题，实现了无缝的全场景重照明。图像翻译模型在~CARLA~模拟器上用~18~种车型、5~种天气条件下的~21,600~对图像训练，通过流匹配目标学习从源背景分布到目标背景分布的速度场。

在优化目标层面，标准~EoT~框架对所有物理配置均匀采样、最小化期望损失~$\mathbb{E}_{c \in C}[\mathcal{L}]$，忽略了损失方差。R-PGA~观察到一种空间一致性现象：无论攻击模式或初始化如何，高损失区域（攻击失败）始终聚集在特定物理配置附近（如极端俯仰角或特定光照条件），称为固有配置硬度。这些困难配置的优化被来自容易视角的梯度在早期干扰，导致过早收敛到局部最优，在物理参数空间中形成锯齿状损失地形和损失峰值。为此，R-PGA~提出硬物理配置挖掘（HPCM）策略，将连续物理配置空间~$C$~离散化为~$q$~个区间，维护全局难度表~$S = \{s_1, s_2, \ldots, s_q\}$~跟踪每个配置的历史鲁棒性（初始化为高常数以鼓励早期全面探索）。每步迭代用动量更新机制平滑瞬时损失~$s_i^{(t)} = \mu \cdot s_i^{(t-1)} + (1-\mu) \cdot \mathcal{L}_{curr}$（$\mu \in [0,1)$），确保高难度值反映优化全程始终难以攻击的配置。采样概率~$P(c_i) = \frac{\exp(s_i / \tau)}{\sum_{j=1}^M \exp(s_j / \tau)}$~按难度分数的~Softmax~分布，$\tau$~为温度参数（实验中~$\mu = 0.5$，$\tau = 1$）。R-PGA~从理论上证明~HPCM~在数学上等价于对~Log-Sum-Exp（LSE）目标~$J_{LSE}(\mathbf{a}) = \tau \ln \sum_{i=1}^q \exp(\mathcal{L}_i(\mathbf{a}) / \tau)$~执行随机梯度下降——期望梯度~$\mathbb{E}_{c \sim P}[\nabla_{\mathbf{a}} \mathcal{L}(\mathbf{a}, c)]$~严格等于~$\nabla_{\mathbf{a}} J_{LSE}$。由于~LSE~是最大损失的光滑上界（$\max_i \mathcal{L}_i \leq J_{LSE} \leq \max_i \mathcal{L}_i + \tau \ln M$），最小化~LSE~等价于抑制最坏情况损失峰值，有效扁平化损失地形，保证在所有物理配置下的均匀鲁棒性。这比标准~Min-Max~优化的内层显式搜索高效得多——避免了在每步迭代中对复杂渲染管线执行昂贵的梯度搜索。

R-PGA~的总优化目标为~$\min_{\mathbf{a}} \mathbb{E}_{c \sim D_{HPCM}}[\mathcal{L}_{det}(I_{det}; F)]$，其中~$\mathcal{L}_{det}$~最小化与真值框最大~IoU~匹配的检测框的置信度，$\mathbf{a}$~为~3DGS~的视角无关反照率。在~CARLA~模拟器的~6~种天气、4~种距离（5m--20m）、10~种俯仰角（$0^\circ$--$90^\circ$）和~$360^\circ$~全方位拍摄下，每种方法构建~4,320~张测试图像。在~6~种主流检测器（YOLOv3、YOLOX、Faster R-CNN、Mask R-CNN、Deformable-DETR、PVT）上，R-PGA~在白盒~YOLOv3~上将~AP@0.5~降至~0.0227，平均跨检测器~AP~为~0.0677，全面超越包括~PGA~和~RAUCA-E2E~在内的~8~种基线方法。跨天气、距离、俯仰角和方位角四维物理配置的对比中，R-PGA~一致建立新的最优性能。消融实验验证了物理解耦重建、混合渲染和~HPCM~三个模块的协同贡献。对视觉基础模型~GLIP~和~DINO~的迁移实验进一步证实了~R-PGA~的跨架构攻击泛化能力。在物理域实验中，对抗迷彩经打印贴附于~1:24~比例模型车，在强光/正常光/弱光三种条件下，R-PGA~将~AP@0.5~从~0.9131~降至~0.2840，相比最强竞争基线~RAUCA-E2E（0.3452）额外降解~0.0612，验证了高保真重建、混合渲染管线和~HPCM~优化策略在物理部署中的有效性。

前面的攻击方法聚焦于目标检测这一单一任务，但在自动驾驶等应用中，感知系统通常包含多个协同工作的子任务——检测器判断"有没有物体"而深度估计器判断"物体有多远"，两者的协同错误比单一错误更具威胁。BiTAA（2025.09）~\cite{2509_19793} 正是在此背景下提出，实现同时降解检测和偏置深度的双任务对抗攻击。

从物理世界攻击回到数字域，前馈~3DGS~模型带来了全新的攻击面——攻击者不再需要操纵场景文件本身，只需扰动输入图像。

\paragraph{改了什么：从单任务检测攻击到检测+深度的双任务联合攻击。}
3DGAA~和~PGA~均仅针对目标检测设计攻击目标——3DGAA~最小化检测器对目标类别的置信度，PGA~进一步引入~min-max~优化增强跨背景鲁棒性，但两者的攻击损失仅涉及检测分支。BiTAA~将攻击面从单一检测任务扩展到检测和单目深度估计的双任务联合优化：对检测分支施加置信度抑制（与~3DGAA/PGA~的检测目标形式一致），同时新增有符号、幅度可控的对数深度偏置项，在感兴趣区域（ROI）内实现可控的近距离或远距离误判。与~3DGAA~类似，BiTAA~同样在~3DGS~的全~14~维参数空间中操作——启用位置、旋转、缩放、颜色和不透明度的全部自由度进行端到端优化，而非~PGA~仅优化~SH~零阶项的纯外观路线。

\paragraph{前向管线：冻结的检测器与深度估计器并行。}
BiTAA~将~3DGS~生成网络视为冻结骨干，提供初始高斯集合~$G_0$。对每个标定视角~$v \in V = \{v_1, \ldots, v_m\}$，渲染当前高斯~$G$~的图像~$I_v = R(G, v)$~和干净基准图像~$I_v^0 = R(G_0, v)$。渲染结果同时送入两个冻结的任务网络：冻结检测器~$D$~提供类别过滤后的检测框和置信度分数；冻结单目深度估计器~$M$~分别输出扰动图像的深度图~$d_v = M(I_v)$~和干净图像的深度图~$d_v^0 = M(I_v^0)$，后者每视角缓存一次以提供稳定的残差基准。两个任务网络均不参与训练，梯度从图像空间的复合损失通过可微渲染器回传至高斯参数~$G$。

\paragraph{每视角 ROI 的定义和约束。}
深度操纵被限制在每个视角中语义相关的区域内，而非全图深度。具体地，BiTAA~从冻结检测器的输出中提取目标类别匹配的检测框，经置信度过滤后，对每个框施加向内各向同性收缩（收缩比~$\gamma \in (0,1)$）以抑制边界噪声，得到每视角的~ROI~区域~$\Omega_v$。等价地，用二值掩码~$\mathcal{M}_v$~表示~$\Omega_v = \{x : \mathcal{M}_v(x) = 1\}$。值得注意的是，ROI~仅用于定位深度偏置项——检测损失仍在完整图像上计算。对于没有有效检测框的视角（$|\Omega_v| = 0$），深度损失项被自动跳过。

\paragraph{复合损失函数：检测抑制 + 有符号对数深度偏置 + 正则化。}
BiTAA~的复合损失由四个分量加权组合
\begin{equation}
    L_{total}(G) = \lambda_{det} L_{det} + \lambda_{dep} L_{dep} + \lambda_{shape} L_{shape} + \lambda_{print} L_{print},
\end{equation}

其中默认权重为~$(\lambda_{det}, \lambda_{dep}, \lambda_{shape}, \lambda_{print}) = (1.0, 1.0, 0.2, 0.1)$。

\textbf{检测抑制项}~$L_{det}$~在所有视角和~EOT~增强样本上最小化目标类别的最大置信度
\begin{equation}
    L_{det} = \frac{1}{|V|} \sum_{v \in V} \mathbb{E}_t \left[ -\log(p_v(t) + \epsilon) \right], \quad p_v(t) = \max_{k \in K_v} s_{v,k}^{(t)},
\end{equation}

其中~$p_v(t)$~是~EOT~增强图像~$\tau_t(I_v)$~上所有目标类别检测的最大置信度，$\epsilon > 0$~防止数值问题。该形式将最强检测的置信度推向零，与~3DGAA~的对抗目标~$\min_G \mathbb{E}_{P}[f_D(R(G,P))]$~在数学精神上一致——都是最小化检测器对目标类别的响应。

\textbf{有符号对数深度偏置项}~$L_{dep}$~是~BiTAA~区别于所有前述方法的核心创新。对每个视角，定义逐像素的对数深度残差~$\delta_v(x) = \log(d_v(x) + \epsilon) - \log(d_v^0(x) + \epsilon)$，在~ROI~内取平均并与有符号目标~$(s, \delta^*)$~对齐
\begin{equation}
    L_{dep} = \frac{1}{|V|} \sum_{v \in V} \frac{1}{|\Omega_v|} \sum_{x \in \Omega_v} \bigl(\delta_v(x) - s \cdot \delta^*\bigr)^2,
\end{equation}

其中~$s \in \{+1, -1\}$~控制偏置方向（$s = +1$~为推远，$s = -1$~为拉近），$\delta^* > 0$~设定目标偏置幅度。这一设计的关键性质是\textbf{剂量可控性}：目标~$s \cdot \delta^*$~在~ROI~内对应乘法深度变化~$\exp(s \cdot \delta^*)$，使得攻击者可以精确控制深度误判的方向和幅度。例如，$\delta^* = 0.05$~对应约~$5\%$~的对数深度偏移。当~$\delta^* = 0$~时，深度项退化为零，BiTAA~退化为纯检测攻击——这为后续的跨任务迁移实验提供了消融条件。

\textbf{几何稳定性正则}~$L_{shape}$~惩罚位置、缩放、旋转偏离骨干输出~$G_0$~的程度，同时软约束四元数的单位范数
\begin{equation}
    L_{shape} = \sum_{i=1}^{N} \bigl( w_\mu \|\Delta\mu_i\|_2^2 + w_s \|\Delta s_i\|_2^2 + w_q \|\Delta q_i\|_2^2 \bigr) + \gamma \sum_{i=1}^{N} (\|q_i\|_2^2 - 1)^2,
\end{equation}

公式中~$\Delta\mu_i, \Delta s_i, \Delta q_i$~分别表示位置、缩放、旋转相对于骨干输出~$G_0$~的偏移量，而~$w_{(\cdot)}$~为各参数的权重。该正则项在对抗强度和几何真实感之间取得了平衡——消融实验表明，移除~$L_{shape}$~使~TNR$_{det}$~从~69.43\%~升至~71.80\%，但参数漂移指标（Side$^\Delta$）从~$1.772 \times 10^{-3}$~显著增至~$2.553 \times 10^{-3}$。

\textbf{可打印性正则}~$L_{print}$~约束颜色扰动保持在可打印范围内并抑制图像空间的高频伪影，由三个子项组成：颜色预算约束~$L_\Delta$~限制每个高斯的颜色偏移~$|\Delta c_i^{(ch)}| \leq \Delta_c$；全变分~$L_{TV}$~惩罚渲染残差图像的空间梯度；高频抑制~$L_{HF}$~通过离散傅里叶变换的环形掩码惩罚高频分量。三者组合为~$L_{print} = L_\Delta + \lambda_{TV} L_{TV} + \lambda_{HF} L_{HF}$。移除~$L_{print}$~使~LPIPS~从~0.5396~升至~0.5594，感知质量明显下降。

\paragraph{EOT 物理增强与优化过程。}
在检测分支上，BiTAA~应用~EOT（Expectation Over Transformation）——在优化过程中对渲染图像随机施加光度变换和几何变换（与~PGA~的~EoT~策略类似），以增强从仿真到物理世界的迁移能力。实验表明~EOT~不仅提升物理鲁棒性，还降低了跨变换的深度偏置方差（Var$_{EOT}(\Delta)$~从~$1.5256 \times 10^{-5}$~降至~$1.4324 \times 10^{-5}$），同时改善感知质量（LPIPS~从~0.5837~单调降至~0.5396）。优化采用一阶方法更新高斯参数~$G \leftarrow G - \eta \nabla_G L_{total}(G)$，所有任务网络和渲染器保持冻结。

\paragraph{剂量可控性验证。}
通过在~$\delta^* \in [0, 0.10]$~范围内扫描目标幅度、在两个偏置方向~$s \in \{+1, -1\}$~下评估，BiTAA~展现出清晰的近线性剂量响应曲线：随着~$\delta^*$~增加，实际~ROI~平均对数深度偏移~$\Delta$~单调增长，斜率接近~1~且截距可忽略，$s = +1$~（推远）和~$s = -1$~（拉近）两条分支关于原点近似镜像对称。95\%~置信带较窄，表明多视角聚合和~ROI~掩码下的优化方差低、跨视角和跨随机种子一致性好。该剂量响应在~DPT-Large、Depth Anything V1、Monodepth2~三种代表性深度估计器上均成立，模型差异仅反映在斜率紧度而非趋势本身。这意味着攻击者可以通过设定~$(s, \delta^*)$~精确命令"近"或"远"的误判强度。

\paragraph{基线对比：双任务攻击的联合降解优势。}
BiTAA~与两大族基线在统一协议下对比（相同冻结模型、视角集和训练预算）。

\textbf{检测导向攻击}（CAMOU、UPC、DAS、FCA、TT3D）平均将~mAP@0.5~从~74.68\%~降至~46.77\%（均值下降~27.91~点），但对深度估计的影响微弱且不一致——平均~AbsRel~仅从~0.1878~升至~0.18832（+0.00052），DAS~甚至使~RMSE~略降（$-0.0258$）。

\textbf{深度导向攻击}（APA、SAAM、3D2FOOL）使~AbsRel~升至~0.1912/0.1926/0.1952（平均~+0.0052），RMSE~升至~1.3092/1.3137/1.3172（平均~+0.0228），但对检测的影响中等——平均~mAP~为~59.19\%（下降~15.49~点）。相比之下，BiTAA~实现最强的联合降解：mAP~降至~6.05\%（绝对下降~68.63~点，比最佳检测基线~FCA~的~28.24\%~进一步降低~22.19~点）；同时~AbsRel~达到~0.1989（+0.0111，约为深度基线平均增益的~2.1~倍），RMSE~达到~1.3229（+0.0323，约为深度基线平均增益的~1.4~倍）。单一~3D~扰动同时超越检测专用攻击在检测指标上的最佳表现和深度专用攻击在深度指标上的最佳表现，证明了双任务框架的协同增益。

\paragraph{跨任务迁移的不对称性。}
BiTAA~设计了系统的跨任务迁移协议：分别用检测专用损失（$\lambda_{dep} = 0$）和深度专用损失（$\lambda_{det} = 0$）训练，然后在双任务上评估，揭示了一个关键的不对称性——检测$\to$深度的迁移远强于深度$\to$检测的迁移。

\textbf{Det$\to$Depth~方向}：宏平均深度降解~$\bar{\Delta}_{depth}$~达~17.1\%。按检测代理模型分析，两阶段模型引发的跨任务深度效应最大（Mask R-CNN~行均值~28.08\%，Faster R-CNN~13.85\%，SSD~9.40\%）；按深度目标模型分析，基于~Transformer~的~DPT-L~最易受影响（列均值~21.78\%），Depth Anything V1~次之（17.97\%），Monodepth2~最鲁棒（11.58\%），显著单元如~Mask R-CNN$\to$Depth Anything~达~33.74\%。这表明检测器驱动的扰动改变了物体轮廓、阴影和局部光度一致性等物体中心线索，而深度估计器（特别是全局~Transformer~骨干）大量复用这些线索，从而放大了跨任务降解。

\textbf{Depth$\to$Det~方向}：宏平均检测降解~$\bar{\Delta}_{det}$~仅~5.6\%，约为反方向的~$1/3$。DPT-L~和~Depth Anything V1~产生相当的行均值（7.51\%~和~7.41\%），Monodepth2~最弱（1.87\%）；作为目标，单阶段~SSD~明显更脆弱（列均值~10.60\%，峰值~DPT-L$\to$SSD~达~19.00\%），而~Faster R-CNN~和~Mask R-CNN~相对稳健（3.24\%~和~2.96\%）。这归因于单阶段检测器更依赖密集纹理和边界证据——深度导向优化间接扰动了这些特征。总体而言，两个方向的迁移比约为~$3:1$，验证了~BiTAA~双任务框架的设计动机：检测和深度之间的交互既可测量又不对称，单一任务攻击会忽略另一任务上的安全隐患。

\paragraph{物理世界验证。}
BiTAA~在~1:30~比例的~Audi A2~模型上验证物理部署。使用红米~K60 Pro~手机在比例一致的~16.7cm~距离（模拟真实~5m~测试范围）多视角拍摄，对抗纹理通过普通办公打印机输出后简单裁剪粘贴于车模表面（刻意不追求理想制造精度）。在四个视角的~$4 \times 4$~面板对比中，干净图像以~0.98~的置信度被正确检测为车辆，对抗图像则被误分类为雨伞、自行车、水槽等类别。深度热力图显示对抗扰动在车体纹理丰富的面板区域产生有符号的局部对数深度偏移，背景区域基本不受影响——这与~ROI~约束的训练行为一致。值得注意的是，尽管使用粗糙的打印和粘贴工艺（存在轻微色带、边界伪影和表面反光），检测抑制和深度偏置仍然一致有效，表明~EOT~训练对真实世界噪声具有足够的鲁棒性。

以上攻击均针对需要逐场景优化的经典~3DGS~管线——攻击者操纵训练数据或场景参数来嵌入对抗内容。而随着前馈~3DGS~模型的兴起，攻击面发生了根本性转移：攻击目标从场景本身变为输入图像。AdvSplat（2026.03）~\cite{2603_23686} 首次系统研究前馈~3DGS~模型的对抗攻击。与~CLOAK~等攻击逐场景优化~3DGS（操纵训练数据利用~SH~视角依赖性）不同，AdvSplat~的核心区别在于攻击范式的根本转换。前馈~3DGS~模型（如~pixelSplat、MVSplat、DepthSplat）通过大规模预训练实现从少量输入视角快速重建，具有更强的商业部署潜力——但这也意味着攻击者无需操纵场景文件本身，只需向输入图像注入不可感知的扰动即可破坏下游的重建质量。

\paragraph{改了什么：从攻击场景到攻击输入图像。}
传统~3DGS~的对抗攻击将恶意内容嵌入~3D~场景参数或训练数据中，攻击对象是场景本身。前馈~3DGS~的引入彻底改变了这一攻击面——前馈模型将场景重建过程封装为一个黑盒神经网络，输入图像成为唯一的外部接口。攻击者只需构造经过精心扰动的输入图像，即可通过模型的前馈推理管线自动将扰动放大到整个重建结果中。这一范式转换带来两个关键变化：（1）攻击门槛降低——不需要访问训练过程或场景文件，只需篡改上传到互联网的图像即可实施攻击，受害用户可能在不知不觉中下载并使用这些被污染的图像；（2）攻击影响更广——一次生成的对抗图像可被任何使用该前馈模型的用户加载，而非仅影响使用特定场景文件的用户。

从前馈模型的输入图像攻击进一步前移到生成管线的源头，攻击者开始瞄准~3D~生成过程的输入端。

\paragraph{白盒~PGD~攻击：揭示前馈~3DGS~的基础脆弱性。}
作为基础性实验，AdvSplat~首先通过白盒~PGD（Projected Gradient Descent）攻击揭示前馈~3DGS~模型族的固有脆弱性。与分类模型使用分类损失不同，前馈~3DGS~的重建任务没有离散标签，因此攻击目标改为重建质量本身——以~MSE（均方误差）和~LPIPS（学习感知图像块相似度）的加权组合作为攻击损失：

\begin{equation}
    L = L_{\text{MSE}}(\hat{I}_t, \hat{I}_t^r) + \lambda \cdot L_{\text{LPIPS}}(\hat{I}_t, \hat{I}_t^r),
\end{equation}

其中~$\hat{I}_t$~为第~$t$~步的对抗输入图像，$\hat{I}_t^r$~为对应视角的渲染图像，$\lambda$~为平衡权重。这一选择基于两个考量：MSE~和~LPIPS~几乎被所有前馈~3DGS~模型的训练损失所采纳，且通常贡献最大的梯度信号，因此攻击目标与模型优化方向天然对齐。PGD~从干净图像出发，在~$L_\infty$~约束球内沿损失函数关于输入图像的梯度方向迭代更新，精确定位对重建质量影响最大的像素区域。实验表明，白盒攻击可导致重建质量完全崩溃——渲染图像中出现大量黑色伪影，甚至完全坍塌为近乎全白的图像。

\paragraph{跨模型迁移攻击的局限性。}
一个自然的问题是：能否在某个模型上生成对抗样本后迁移到其他模型？AdvSplat~的实验给出了否定的答案。对三个前馈~3DGS~模型（DepthSplat、NoPoSplat、AnySplat）的穷举交叉评估显示，尽管迁移攻击确实能轻微降低重建质量，但降级幅度远小于查询式攻击。其根本原因在于：前馈~3DGS~模型通常基于不同的高容量~Transformer~骨干网络，使用不同的预训练数据、架构偏置和表示几何，导致特征空间差异显著——这使迁移攻击对前馈~3DGS~几乎不可行。

\paragraph{频域参数化：降低查询维度的核心策略。}
由于迁移攻击效果有限，AdvSplat~转向更实际的查询式黑盒攻击。然而查询式攻击面临一个核心挑战：前馈~3DGS~模型的输入图像通常是高分辨率的，搜索维度~$H \times W \times 3$~极其庞大，需要海量查询才能在如此高维的空间中找到有效的对抗方向。为解决这一问题，AdvSplat~提出基于~DCT（离散余弦变换）的频域参数化策略，将扰动约束在低频子空间中。具体地，图像首先被划分为~$J$~个大小为~$n \times n$~的等尺寸图像块，对每个块应用~DCT-II~正交变换将其转换到频域。在频域系数矩阵中，仅扰动左上角~$s \times s$~的低频子块（对应图像的主要结构信息），保持其余高频系数不变，然后通过逆~DCT~将扰动后的系数映射回像素空间。这一设计将优化空间的维度从像素级的~$H \times W \times 3$~大幅降低到~$J \times s \times s \times 3$（实验中~$n=8$，$s=3$，即每个~$8 \times 8$~的~64~维空间仅优化~$3 \times 3=9$~个低频系数），用远少于像素级攻击的查询次数即可逼近有效的对抗方向。

\paragraph{梯度估计黑盒算法：NES~有限差分逼近。}
在频域参数化的基础上，AdvSplat~开发了两种查询高效的实用黑盒算法。第一种是基于~NES（Natural Evolution Strategies）的梯度估计算法。由于黑盒设定下无法直接计算梯度，NES~通过有限差分策略逼近梯度：在频域低频子空间中采样~$M$~对高斯噪声向量，分别构造正向和反向扰动，将它们通过逆~DCT~映射到像素空间并输入模型获取渲染结果，然后用损失差值加权平均来估计梯度方向：

\begin{equation}
    \nabla_{\hat{I}} L \approx \frac{1}{2M\sigma} \sum_{m=1}^{M} \left[L(\hat{I}_m'^{+}, I_{r,(m)}'^{+}) - L(\hat{I}_m'^{-}, I_{r,(m)}'^{-})\right] u_m,
\end{equation}

其中~$u_m$~为标准高斯噪声，$\sigma$~为搜索方差。估计出梯度后，即可在像素空间使用符号梯度更新并投影回~$L_\infty$~约束球内，实现与白盒~PGD~类似但完全无需模型内部的攻击流程。

\paragraph{无梯度黑盒算法：CMA-ES~进化策略。}
第二种是无需任何梯度信息的~CMA-ES（Covariance Matrix Adaptation Evolution Strategy）算法。当模型对噪声不敏感或损失地形不够平滑时，梯度估计可能失效，此时~CMA-ES~提供了一种互补的搜索策略。CMA-ES~在频域低频子空间中迭代采样候选扰动：每轮迭代从以当前均值~$\mu^{(t)}$~和协方差~$\Sigma^{(t)}$~参数化的多元高斯分布中采样~$B$~个候选频域扰动，通过逆~DCT~映射到像素空间并输入模型计算损失。然后按损失降序排列，选择前~$B/2$~个表现最优的候选来更新均值和协方差：

\begin{equation}
    \mu^{(t+1)} = \sum_{\ell=1}^{B/2} w_\ell \, z_{\ell:B}^{(t)}, \quad
    \Sigma^{(t+1)} = (1-c_\Sigma)\,\Sigma^{(t)} + c_\Sigma \sum_{\ell=1}^{B/2} w_\ell \, (z_{\ell:B}^{(t)} - \mu^{(t)})(z_{\ell:B}^{(t)} - \mu^{(t)})^\top,
\end{equation}

其中重组合权重~$w_\ell$~按对数递减分配（排名越靠前权重越大），使搜索方向自适应地偏向损失地形中更有利的区域。完整的~CMA-ES~还包括共轭进化路径、全局步长自适应调整和协方差秩一更新等机制，确保在高维频域空间中的高效搜索。最终，从收敛的分布~$\mathcal{N}(\mu^{(T)}, \Sigma^{(T)})$~中采样的噪声即为所求的频域扰动。

\paragraph{实验效果。}
在~Re10K~和~DL3DV~两个基准数据集上，AdvSplat~对三个代表性前馈~3DGS~模型（DepthSplat、NoPoSplat、AnySplat）进行了全面评估。对抗扰动的~$L_\infty$~范数被严格约束为~8/255——即使在最坏情况下每个像素最多扰动~8/255，对抗图像与干净图像之间的~PSNR~仍高于~30\,dB，扰动对人眼几乎不可感知。然而如此微小的扰动却导致渲染质量急剧恶化：定量指标上，PSNR、SSIM、CLIP~相似度和~DINO~相似度均显著下降，LPIPS~大幅上升（指标含义反转——越低越好变为越高越好）；定性上，渲染图像出现大量黑色伪影甚至完全坍塌，3D~高斯点云的颜色和不透明度发生剧烈变化。在梯度估计和~CMA-ES~两种黑盒变体的比较中，当梯度估计足够准确时，梯度方法优于无梯度方法；但~NoPoSplat~是个例外——由于其使用混合数据集训练导致模型容量更大、梯度地形更复杂，CMA-ES~反而优于梯度估计。频域参数化的效率优势也得到了验证：在相同查询次数下，使用~DCT~低频约束的攻击在收敛时达到更高的损失值，证明降维策略确实用更少的查询找到了更有效的对抗方向。

\paragraph{设计动机。}
AdvSplat~的研究动机源于前馈~3DGS~模型商业化部署的安全风险。一个高度可能的现实攻击场景是：攻击者反复查询前馈~3DGS~模型的~API~生成对抗扰动图像，随后将这些图像上传到互联网。普通用户可能在毫无察觉的情况下下载并使用这些被污染的图像进行~3D~重建，导致输出质量严重下降——这不仅损害公司的商业声誉，更可能动摇用户对~3D~重建技术的信任基础。随着前馈~3DGS~朝着基础模型方向发展，确保其安全性变得与提升其性能同等重要。

前面的攻击方法分别从训练数据操纵（CLOAK）、参数空间扰动（DAGGER）、物理对抗物体（3DGAA/PGA/R-PGA）和前馈模型输入（AdvSplat）等角度攻击~3DGS~管线，而另一条攻击路线瞄准的是~3DGS~的更上游——从源图像到~3D~模型的完整生成管线。与上述方法（攻击~3DGS~场景本身、其训练数据或前馈模型输入）不同，M-IFGSM（2024.12）~\cite{2412_02803} 攻击的是~3D~生成管线的输入端——在源图像上生成对抗噪声，使基于~CLIP（Contrastive Language-Image Pre-training）~视觉语言模型的~3D~生成系统产出可逃避检测的~3D~物体。

攻击方法至此已经覆盖了从源图像到最终输出的完整攻击链。接下来转向防御方法的研究。

\paragraph{改了什么：从攻击3DGS场景到攻击生成管线输入。}
从攻击阶段看，CLOAK~操纵~3DGS~训练数据、DAGGER~扰动~3DGS~参数、AdvSplat~扰动前馈模型输入——这些方法的攻击对象已经是~3D~表示本身（或其直接输入）。M-IFGSM~则将攻击点进一步前移到生成管线的源头：攻击者不触及~3DGS~重建算法，而是在图像采集阶段即对源图像注入对抗噪声，后续~3DGS~重建照常执行，但生成的~3D~模型已携带对抗特征。从攻击目标看，CLOAK~等方法的攻击对象是下游~2D~目标检测器（如~YOLO、Faster R-CNN），而~M-IFGSM~的攻击目标是~CLIP~视觉语言模型的零样本分类能力——使~CLIP~无法正确识别~3D~物体渲染结果所属类别。

\paragraph{掩码迭代快速梯度符号法。}
传统~FGSM（Fast Gradient Sign Method）对整张图像均匀施加扰动——$X^{\text{adv}} = X + \epsilon \cdot \text{sign}(\nabla_X J(X, y_{\text{true}}))$。M-IFGSM~在此基础上引入两个关键改进。首先，将单步更新改为迭代更新（即~I-FGSM），通过多轮梯度累积产生更强的对抗效果。其次，也是核心创新——用分割掩码将扰动约束在物体区域：利用~SAM（Segment Anything Model）提取输入图像中目标物体的分割掩码~$M$，然后仅对掩码覆盖区域施加扰动。形式化地，非定向攻击的迭代更新为
\begin{equation}
    X_{N+1}^{\text{adv}} = \text{clip}_{0,255} \left\{ X_{\text{inv}} + M \odot \left( X_N^{\text{adv}} + \epsilon \cdot \text{sign}\left(\nabla_{X_N} J(X_N^{\text{adv}}, y_{\text{true}})\right) \right) \right\},
\end{equation}

其中~$X_{\text{inv}} = X \odot (1 - M)$~为背景区域（保持不变），$M$~为~SAM~生成的二值分割掩码，$\odot$~为逐元素乘法，$\epsilon$~控制扰动幅度，$\text{clip}_{0,255}$~将像素值裁剪到合法范围。迭代过程中使用早停机制——当真实标签的概率降至零且损失超过阈值（$\theta = 20.00$）时终止迭代，避免过度扰动。

\paragraph{掩码约束的设计动机。}
将扰动集中在物体掩码区域而非整张图像，基于两个分析。（1）扰动效率——背景区域对~CLIP~判断物体类别的贡献很小，在背景上施加扰动不仅浪费对抗预算，还会在~3DGS~重建过程中引入不必要的全局伪影，降低重建质量。将所有扰动能量集中于物体区域，可以在更小的~$\ell_\infty$~范数下达到更强的攻击效果。（2）检测规避——掩码约束使扰动在物体边界处自然截断，避免了全图扰动在物体-背景交界处的明显不连续性，使对抗噪声在人类视觉检查中更难被发现。

\paragraph{攻击管线与实验设置。}
M-IFGSM~的完整攻击管线包含三个阶段：（1）掩码生成——用~SAM~（sam-vit-h~检查点）从输入图像提取目标物体掩码~$M$；（2）对抗扰动生成——用~M-IFGSM~迭代生成掩码约束的对抗图像；（3）3D~模型重建——用~3DGS~从对抗图像重建~3D~模型（35~张训练图像，6~张测试视角渲染）。实验使用~CO3D~数据集的~8~个物体类别（每类~41~张图像，缩放至~$224 \times 224$~像素以匹配~CLIP ViT-B/16~输入尺寸），攻击目标为~CLIP ViT-B/16~模型。该模型将输入图像分割为~$16 \times 16$~的~patch~序列，通过~Transformer~自注意力层处理——M-IFGSM~的梯度扰动正是通过反向传播穿过这些~patch~嵌入和注意力层来精确定位对分类决策影响最大的像素区域。

\paragraph{实验效果。}
在对抗图像上，M-IFGSM~将~CLIP~的~Top-1~准确率从~94.9\%~降至~2.1\%，Top-5~准确率从~99.6\%~降至~6.4\%，接近完全失效。更重要的是，对抗噪声成功迁移到~3DGS~重建后的渲染视角——从训练相机位置渲染的~3D~模型图像，CLIP~Top-1~准确率从~95.4\%~降至~12.5\%；从测试（新）相机位置渲染时，准确率从~91.2\%~降至~35.4\%。测试视角攻击效果的衰减表明，对抗噪声在~3DGS~重建过程中被部分平滑，但仍然显著破坏了~CLIP~在新视角下的分类能力。

\subsubsection{防御方法}

本路线介绍从频率分析防御、检测-净化防御和编辑防护等角度保护3DGS系统安全性的防御方法。

防御方面，DefenseSplat（2026.02）~\cite{2602_19323} 修改了~3DGS~训练管线的图像预处理环节——用~DWT~的高频滤波替代直接使用原始输入图像——3DGS~的渲染质量损失和优化流程保持不变。DWT~是一种同时捕捉空间位置和频率信息的信号分解工具，可将图像分解为四个子带：$LL$（低水平+低垂直，保留整体结构与平滑内容）、$LH$（低水平+高垂直，水平边缘特征）、$HL$（高水平+低垂直，垂直边缘特征）和~$HH$（高水平+高垂直，对角细节）。

\paragraph{设计动机：对抗扰动集中在高频。}
DefenseSplat~的设计起点是一个频率域观察。作者用~DWT~将对抗图像分解后，通过多视角特征匹配（SuperPoint+LightGlue）比较干净图像与对抗图像在各子带上的跨视角一致性。结果显示，高频子带（$LH$、$HL$）的匹配率下降远比低频子带（$LL$）严重——在~Mip-NeRF 360~数据集上，高频匹配率下降~19.56\%，而低频仅下降~5.28\%；在~Tanks and Temples~上，分别为~15.25\%~和~3.45\%。这说明~3DGS~对抗攻击主要将随机噪声注入高频分量，而对低频结构的影响相对较小。$HH$~子带因能量占比极低（在几乎所有数据集上~$<$0.08\%）而被忽略。

\paragraph{核心机制：高频置零滤波。}
基于上述分析，DefenseSplat~的核心操作非常简洁——将~DWT~分解后的三个高频子带直接置零，仅保留~$LL$~低频子带进行逆变换重建：

\begin{equation}
I_f = \mathrm{iDWT}\bigl(\mathrm{DWT}(I')_{LL},\; \mathbf{0},\; \mathbf{0},\; \mathbf{0}\bigr),
\end{equation}

其中~$I_f$~是滤波后的图像，$I'$~是对抗输入图像，$\mathrm{iDWT}$~是逆离散小波变换。直觉上，这相当于在频域切除所有高频噪声，只让低频的场景结构信息进入~3DGS~训练流程。

\paragraph{3DGS~固有模糊效应的协同利用。}
仅靠高频置零并不充分——$LL$~子带仍然可能包含具有低跨视角一致性的人造纹理。DefenseSplat~巧妙地利用了~3DGS~的一个固有性质：当高斯拟合跨视角一致性低的区域时，为了最小化~$L_1$~损失，高斯会自动对多个视角的颜色取平均，产生模糊重建。这种模糊效应在正常场景中并不理想，但在防御场景中恰好进一步抑制了缺乏跨视角一致支持的对抗伪影。

\paragraph{尺度正则化：抑制对抗纹理诱导的拉长高斯。}
另一类对抗伪影则具有高跨视角一致性——攻击注入的人造纹理在多个视角间相互一致，此时~3DGS~会尝试用拉长的高斯（elongated Gaussians）来拟合这些模式，导致高斯数量膨胀和内存消耗增加。为此，DefenseSplat~设计了基于~ReLU~的尺度正则化损失，约束每个高斯沿三个主轴尺度的归一化方差：

\begin{equation}
\mathcal{L}_{\mathrm{scale}} = \mathrm{ReLU}\bigl(\sigma_s - \tau\bigr),
\end{equation}

其中~$\sigma_s$~是归一化方差，$\tau$~是预设阈值。直觉上，该损失只惩罚形状过于拉长的高斯（归一化方差大的），而不影响两种有益的高斯：小球形高斯（重建精细细节、缓解过度重建）和大扁平高斯（重建平滑区域、缓解欠重建）。实验中尺度损失权重设为~$\lambda_{\mathrm{scale}} = 1 \times 10^{-5}$，攻击强度设为~$\varepsilon = 16/255$。

\paragraph{实验效果。}
在~Mip-NeRF 360、Tanks and Temples~和~LLFF~三个基准上，DefenseSplat~在不损害干净数据训练质量的前提下显著增强~3DGS~鲁棒性。由于此前尚无专门针对~3DGS~对抗攻击防御的方法，作者选取了方向最相关的两个方法作为对比基线：Difix3D+（基于扩散模型的图像增强）和~CompactGS（精简冗余高斯）。DefenseSplat~在所有数据集上均位列前两名，在高斯数量控制和渲染速度方面表现尤为突出——验证了"频域滤波+尺度正则化"这条简洁路线的有效性。

Spectral Defense（2026.03）~\cite{2603_12796} 在高斯和图像域联合提出频谱防御。与锚点~DefenseSplat~（通过~DWT~高频置零+尺度正则化防御）采用互补但不同的技术路线。DefenseSplat~在~2D DWT~域操作且置零全部高频子带，是一种全局性的频谱切除；而~Spectral Defense~同时在~3D~参数空间和~2D~频谱域操作，保留自然的高频内容，仅选择性移除攻击诱导的异常——前者是"一刀切"的频域过滤，后者是"精准定位"的频谱剪枝。

\paragraph{3D~频率滤波：傅里叶频谱引导的参数空间剪枝。}
在~3D~层面，Spectral Defense~将剪枝从效率优化工具重新定位为安全原语（security primitive），通过~3D~高斯的傅里叶频谱分析识别并移除对抗性冗余高斯。其理论基础在于：每个~3D~高斯~$G$~的傅里叶响应~$F(t)$~完全由其协方差矩阵~$\Sigma$~决定——空间变量~$x$~和频率变量~$t$~构成~Fourier~对，频谱可分解为幅度衰减~$\exp(-2\pi^2 t^\top \Sigma t)$~和线性相位项~$-2\pi t^\top \mu$，其中相位项仅取决于位置~$\mu$~且不影响频谱分布，因此高斯的频率特性完全由~$\Sigma$~的尺度编码。具体地，最小特征值~$\sigma_{min}$~越小的方向空间扩展越窄，高频衰减越弱、高频响应越强。基于此，Spectral Defense~在固定截止频率~$t$~下定义两个核心指标：高频衰减得分~$S(G) = \exp(-2\pi^2 t^2 \sigma_{min}^2)$，衡量每个高斯在统一频谱参考下的预期高频质量；高频重要性权重~$W(G) = (1-S(G))^\gamma$（$\gamma > 0$~控制非线性放大），将衰减得分映射为重要性——高频响应过强的高斯被分配低权重，贡献适度且结构相关高频内容的高斯保留较高权重。实验中~$t=8$、$\gamma=2$~为最优设置，且在~$t \in [4,10]$、$\gamma \in \{1,2,3,4\}$~范围内结果均稳定。

为整合空间可见性信息，Spectral Defense~引入光线命中指示器~$\phi(G)$——将高斯投影到图像平面后统计其被训练光线交点的次数（沿用~LightGaussian~的定义），反映每个高斯在场景中的空间支撑程度。综合剪枝得分定义为~$\text{score}(G) = W(G) \cdot \phi(G)$，同时反映频谱贡献和空间支撑。得分低的高斯（很少被观测到却具有强高频行为）被判定为攻击诱导的异常。剪枝采用周期性策略：每~$T_{\text{prune}}=100$~步训练迭代执行一次，随机采样~$K'=48$~个代表性相机视角（平衡覆盖度与计算开销），按~$\text{score}(G)$~排序后剪除最低的~$\delta\%$~个高斯。剪枝比例~$\delta$~根据场景复杂度自适应调整：小型场景（NeRF-Synthetic）~$\delta=3.0\%$，大型场景（Tanks and Temples）~$\delta=4.5\%$，复杂场景（Mip-NeRF 360）~$\delta=5.0\%$。

\paragraph{2D~频谱正则化：方向分布引导的各向异性抑制。}
在~2D~层面，Spectral Defense~对渲染图像进行~2D DFT~分解以补偿~3D~剪枝的不足——即使参数空间已过滤异常高斯，受害模型仍在投毒视角上优化，渲染结果中可能仍残留对抗伪影。具体流程为：对渲染图像~$\hat{V}^p$~施加~2D DFT~获得频率响应~$\mathcal{F}(u,v)$，提取由频率阈值~$r_{\text{min}}=0.3$~和~$r_{\text{max}}=0.9$~界定的高频带~$E(u,v)$——该带通选择聚焦于具有显著能量的高频区域，避开低频结构和极高频噪声。随后将角度域~$[-\pi, \pi)$~离散化为~$B=36$~个均匀间隔的角度扇区，第~$b$~个扇区~$\Omega_b$~收集频率向量落在该角度范围内的频谱系数，聚合得到角度能量密度~$E_b$~和方向分布~$P_b = E_b / \sum_j E_j$。$B=36$~的选择经过消融验证——更大的~$B$~使角度直方图更精细但结果差异很小，方法对~$B$~的具体取值具有弱依赖性。

Spectral Defense~用归一化熵~$norm(H) = H / \log B$~（取值~$[0,1]$）量化方向集中度——干净图像的高频能量近似各向同性分布（$P_b$~趋于平坦，熵高，$norm(H) \to 1$），攻击诱导的方向性伪影（条纹状、带状噪声）使高频能量集中在狭窄方向（$P_b$~在特定扇区急剧峰值，熵低，$norm(H) \to 0$）。各向异性损失~$L_{ani} = 1 - norm(H)$~惩罚低熵的方向集中频谱，同时保留各向同性的自然纹理。该损失逐视角计算后在~$K$~个渲染视角上平均，得到频域正则化项~$L_{\text{freq}} = \frac{1}{K}\sum_{k=1}^{K} L_{ani}(\hat{V}_k^p)$。此外，Spectral Defense~还引入全变分损失~$L_{tv}$~促进局部平滑以抑制像素级噪声。

\paragraph{整体训练目标与实验效果。}
Spectral Defense~的总体训练目标将重建损失与频域先验统一为
\begin{equation}
\min \; \bigl( L(\hat{V}^p, V^p) + \lambda \bigl( L_{\text{freq}}(\hat{V}^p) + L_{tv}(\hat{V}^p) \bigr) \bigr),
\end{equation}

其中~$L(\hat{V}^p, V^p)$~是标准~3DGS~重建损失（光度+结构相似性），$\lambda$~控制频域正则化强度。损失权重~$\lambda$~同样根据场景复杂度调整：小型场景~$\lambda=4$，复杂场景~$\lambda=5$——增大~$\lambda$~持续抑制投毒诱导的过度生长并略提升~PSNR，但在小型场景中过大的~$\lambda$~会过度抑制自然纹理导致~PSNR~微降。每个训练迭代通过最小化该目标更新高斯参数，每~$T_{\text{prune}}$~步执行一次频率感知剪枝，两者联合稳定渲染频谱分布并抑制~3D~参数空间中的对抗性致密化。在~Tanks and Temples、NeRF-Synthetic~和~Mip-NeRF 360~三个基准上，Spectral Defense~将过度生长抑制最多~5.92$\times$（高斯数量接近恢复到干净数据训练水平），内存降低~3.66$\times$，渲染速度显著提升（如~TT-Barn~从~120 FPS~提升至~236 FPS），同时~PSNR~在多场景上保持或超过投毒基线（如~NS-chair~达~25.00、MIP-bonsai~达~29.07）。在干净输入下~Spectral Defense~同样保持有竞争力的渲染质量，对~Scaffold-GS~等黑盒受害系统也展现出鲁棒的泛化能力。

频率域防御从信号处理角度过滤对抗扰动，而另一条防御路线是从数据层面检测并净化投毒样本。RemedyGS（2025.11）~\cite{2511_22147} 提出首个有效的黑盒防御框架，针对~Poison-splat~等计算成本攻击。三者的定位差异可以概括为：锚点~DefenseSplat~在~3DGS~训练管线内部修改图像预处理环节（DWT~高频置零+尺度正则化），Spectral Defense~同时在~3D~高斯参数空间和~2D~渲染图像频谱域联合操作，而~RemedyGS~完全不修改~3DGS~训练流程——在训练流程上游独立部署检测和净化模块，对~3DGS~系统完全透明。这样，RemedyGS~可以即插即用地部署在已有~3DGS-as-a-service~系统中，无需重新训练或调整高斯模型。

\paragraph{投毒图像检测器。}
RemedyGS~首先需要一个精确的检测器将投毒图像从正常输入中分离出来。对~DefenseSplat~而言所有图像统一经过~DWT~滤波，不存在区分步骤；而~RemedyGS~的选择性处理策略意味着检测精度直接影响整体防御效果——干净图像被误判为投毒会遭受不必要的净化降质，投毒图像被漏检则直接威胁系统安全。检测器基于~4~层~2D~卷积网络~$f_{det}$~实现，末端连接线性分类头。卷积核天然擅长捕捉局部纹理特征，恰好对齐~Poison-splat~攻击的指纹特征——攻击通过最大化~TV~分数在投毒图像中引入不自然的高频噪声和更显著的边缘结构，这些异常纹理模式成为可靠的检测签名。

检测器的训练数据构建是一个关键设计。作者基于~DL3DV~大规模多视角数据集（$960 \times 540$~分辨率）采样~320~个场景，对每个场景应用~Poison-splat~计算成本攻击，生成约~100~万对（干净图像~$V_{cln}$，投毒图像~$V_{poi}$）的配对训练集~$D_{det} = \{(V_{poi}, y_{poi}=1), (V_{cln}, y_{cln}=0)\}$。检测器以学习率~$1 \times 10^{-4}$~训练~50~个~epoch，批大小~16，使用交叉熵损失：

\begin{equation}
L_{det} = \frac{1}{|D_{det}|} \sum_{i=1}^{|D_{det}|} \mathrm{CE}_{\mathrm{loss}}(y_i, f_{det}(V_i; \theta)),
\end{equation}

其中~$\theta$~为检测器参数。评估时，来自~Mip-NeRF 360~数据集的~$1600 \times 1066$~分辨率图像被随机裁剪为训练分辨率——由于投毒纹理是局部注入的，随机裁剪仍能保留足够的检测线索。实验显示检测器在跨数据集评估中实现接近~100\%的检测准确率。

\paragraph{净化器：信息论驱动的图像恢复。}
被检测器标记为投毒的图像进入净化器。净化器采用编码器-解码器架构的非线性变换网络：编码器~$f_\phi$（5~层~2D~卷积）识别并消除攻击者施加的有毒纹理，逐步提取原始图像特征；解码器~$g_\phi$（对称转置卷积）从提取的特征中重建净化图像~$V_{rec} = g_\phi(f_\phi(V_{poi}))$。编码器与解码器之间通过加法式跳跃连接（additive skip connections）传递细粒度特征，保留原始图像的高频细节。消融实验表明，加法式跳跃连接显著优于通道拼接（channel concatenation）——后者会将攻击者注入的噪声连同有用特征一起传递到解码器，而加法式连接通过逐元素相加使有用信号增强、噪声信号被抑制。

净化器的训练目标从信息论角度建立。设原始干净图像服从分布~$p_{cln}(V_{cln})$，攻击者通过条件分布~$p_A(V_{poi}|V_{cln})$~生成投毒图像，净化器的目标是使恢复图像~$V_{rec}$~尽可能保留关于原始干净图像~$V_{cln}$~的信息，即最大化互信息~$I(V_{cln}; V_{rec})$。直接优化互信息是不可处理的，作者引入~Barber-Agakov~变分下界将其转化为可计算的替代目标。该下界指出，对任意随机变量~$x$~和~$y$，互信息满足：

\begin{equation}
I(x; y) \geq H(x) + \mathbb{E}_{p(x,y)}[\log q(x|y)],
\end{equation}

其中~$q(x|y)$~为任意变分分布。应用到净化问题中：

\begin{equation}
I(V_{cln}; V_{rec}) \geq H(V_{cln}) + \mathbb{E}_{p(V_{cln}, V_{rec})}[\log q(V_{cln}|V_{rec})].
\end{equation}

进一步假设变分分布~$q(V_{cln}|V_{rec})$~服从独立多元高斯分布~$\mathcal{N}(V_{rec}, \sigma^2 I)$（$\sigma^2$~为方差，$I$~为单位矩阵），则：

\begin{equation}
I(V_{cln}; V_{rec}) \geq H(V_{cln}) + \frac{d}{2}\log(2\pi\sigma^2) - \frac{1}{2\sigma^2}\mathbb{E}_{p(V_{cln}, V_{rec})}[\|V_{cln} - V_{rec}\|_2^2],
\end{equation}

其中~$d$~为~$V_{cln}$~的维度。由于~$\sigma$、$H(V_{cln})$~和维度相关项为常数，最大化互信息下界等价于最小化~$\mathbb{E}[\|V_{cln} - V_{rec}\|_2^2]$，即~MSE~损失。最终净化器的训练目标为：

\begin{equation}
L_{pur} = \min_\phi\; \mathbb{E}_{p_{cln}(V_{cln})} \mathbb{E}_{p_A(V_{poi}|V_{cln})} [\|V_{cln} - g_\phi(f_\phi(V_{poi}))\|_2^2].
\end{equation}

直觉上，Barber-Agakov~下界为~MSE~最小化这一常用做法提供了严格的信息论基础——最小化像素级重建误差实际上是在最大化恢复图像与原始干净图像之间的互信息。

\paragraph{对抗训练：解决~MSE~过平滑问题。}
仅用~MSE~训练的净化器输出存在过度模糊问题，这是卷积神经网络在~MSE~监督下的已知局限——网络倾向于对所有可能的重建取平均，牺牲高频细节以换取低~MSE~值。为此，RemedyGS~引入对抗训练机制，强制恢复图像的分布与自然干净图像的分布对齐。

具体而言，条件判别器网络~$F(V|c)$~负责区分来自~$p_{cln}(V_{cln})$~的干净图像与来自~$p_{rec}(V_{rec})$~的恢复图像之间的分布差异，其中条件信息~$c$~是净化器编码器所提取的潜在表示。该判别器由~4~层~2D~卷积堆叠而成；条件潜在表示先通过~2D~卷积上采样至与输入图像匹配的空间维度，再沿通道维度与输入图像拼接后送入判别器，预测拼接图像来自~$p_{cln}$~还是~$p_{rec}$。判别器的训练目标为：

\begin{equation}
L_F = \max_{\theta_F}\; -\mathbb{E}_{p_{cln}}[\log F(V|c)] - \mathbb{E}_{p_{rec}}[\log(1 - F(V|c))].
\end{equation}

给定固定的净化器，最优判别器为~$F^*(V|c) = p_{cln}(V) / (p_{cln}(V) + p_{rec}(V))$。净化器的对抗训练目标则利用该最优判别器反向传播梯度：

\begin{equation}
L_G = \max_\phi\; \{-\mathbb{E}_{p_{cln}}[\log F^*(V|c)] - \mathbb{E}_{p_{rec}}[\log(1 - F^*(V|c))]\}.
\end{equation}

数学上可证明~$L_G$~与~Jensen-Shannon~散度直接相关：

\begin{equation}
L_G = \max_\phi\; \{2\log 2 - \mathrm{JSD}(p_{cln}(V) \| p_{rec}(V))\}.
\end{equation}

由于~JSD~非负，$L_G$~在且仅在~$p_{rec} = p_{cln}$~时达到最大值——理论上保证对抗训练驱动恢复分布收敛到干净分布。

\paragraph{训练流程与损失权重。}
训练采用三阶段策略：（1）~MSE~热身——先用纯~MSE~损失训练净化器~50~个~epoch（学习率~$1 \times 10^{-4}$，批大小~16），稳定基本恢复能力；（2）判别器预热——冻结热身净化器，单独训练判别器~50~个~epoch（学习率~$5 \times 10^{-4}$）；（3）交替对抗训练——以热身模型初始化，净化器和判别器交替更新~40~个~epoch（学习率均为~$2.5 \times 10^{-4}$）。最终的净化器总损失为：

\begin{equation}
L_{pur}^* = \lambda_1 L_{MSE} + \lambda_2 L_{LPIPS} + \lambda_3 L_G,
\end{equation}

其中~$\lambda_1 = 0.23$、$\lambda_2 = 100$、$\lambda_3 = 1$。LPIPS~损失的高权重（$\lambda_2 = 100$）反映了感知质量恢复的优先级——净化器需要逐像素重建准确（MSE），更需要人类视觉感知上的高保真（LPIPS），同时通过对抗损失~$L_G$~确保统计分布一致性。

\paragraph{多场景防御效果。}
在~NeRF-Synthetic、Mip-NeRF 360~和~Tanks and Temples~三个基准上，RemedyGS~在白盒攻击下将~GPU~内存使用恢复到接近正常水平（相比无防御时降低超过~$2\times$），同时~PSNR~相比朴素基线（图像平滑和高斯数量限制）提升最高~4~dB，SSIM~提升~0.24。图像平滑虽然能抑制高频噪声但导致过度防御——对所有图像（包括干净的）统一滤波，造成高达~9~dB~的重建质量降质；高斯数量上限则迫使高斯从非锐化区域重新分配到锐化区域，导致局部细节丢失。相比之下，RemedyGS~的检测器将干净图像直接旁路到~3DGS~训练，不引入任何质量损失。

在黑盒攻击场景中（攻击者使用~Poison-splat~但受害者采用~Scaffold-GS~而非~vanilla 3DGS），RemedyGS~防御效果保持一致，比朴素基线的效用提升超过~1~dB——验证了数据驱动方法对未知受害模型的泛化能力。

在自适应攻击场景中——攻击者知晓防御策略，在攻击目标中额外加入互信息正则项~$I(V_{cln}; V_{poi})$~使投毒图像分布更接近干净分布以绕过检测——RemedyGS~依然保持有效防御。作者设置了三个自适应强度~$\beta \in \{1, 5, 20\}$~控制互信息约束权重，结果表明即使~$\beta = 20$~时（投毒图像与干净图像的互信息极强），防御性能仅有轻微退化。此外，即使~10\%的投毒图像漏检未被净化，系统计算成本仍与正常情况几乎相同——这归因于~Poison-splat~攻击需要大量投毒图像协同作用才能显著推高计算成本，少量漏检的影响可忽略。

上述防御方法针对的是训练数据投毒和重建质量攻击，而随着扩散模型编辑能力的增强，如何保护已训练好的~3DGS~场景不被恶意编辑也成为重要的防御需求。与锚点~GuardMarkGS~（在统一框架中联合优化水印和对抗）和~DefenseSplat~（修改训练管线防御投毒攻击）不同，AdLift（2025.12）~\cite{2512_07247} 的防御目标是已训练好的~3DGS~场景——将~2D~图像域已证明有效的对抗扰动（阻止扩散模型编辑）提升到~3D~高斯表示中，使其在任意新视角渲染时都携带保护信号。

从水印到对抗攻击与防御，再到隐写与信息隐藏，~3DGS~安全研究的版图逐步完整。

\paragraph{问题设定：2D~对抗扰动向~3DGS~的直接扩展为何失败。}
2D~图像的对抗扰动已证明可有效阻止扩散模型编辑。对于单张~2D~图像~$x^t$，保护后的图像~$x^{prot}$~通过优化对抗损失同时约束扰动幅度获得：$x^{prot} = \arg\min_x L_{adv}(x, y)$，$\text{s.t.}~\|x - x^t\|_p \leq \varepsilon$，其中~$L_{adv}$~是针对编辑模型的对抗损失，$\varepsilon$~控制扰动预算。将此扩展到~3DGS~的一个直观尝试是用~2D~方法分别保护各训练视角的图像，然后用保护后的图像~$\{v^{t,prot}\}$~重新拟合~3DGS。然而，各视角独立优化的扰动在视角间不一致，导致跨视角冲突——拟合出的~3DGS~即使在训练视角上也出现欠拟合，且对未见视角的泛化严重退化。现有方法（如~GaussianMarker、GuardSplat）通过在高斯参数空间引入软约束来表示扰动，但软约束无法在不可见性与对抗攻击强度之间取得令人满意的平衡。

\paragraph{渲染空间硬约束：在图像域而非参数域约束扰动。}
AdLift~的核心洞察是：与其在高斯参数空间施加难以校准的扰动约束，不如直接在渲染图像空间施加严格的~$\ell_p$~范数上界。形式化地，学习~3DGS~对抗扰动的目标为：

\begin{equation}
G^{prot} = \arg\min_G \sum_{v^t \in V^t} L_{adv}(R(G, v^t), y) \quad \text{s.t.}~\|R(G, v^t) - R(G^{raw}, v^t)\|_p \leq \varepsilon,~\forall v^t \in V^t,
\end{equation}

其中~$R(G, v^t)$~是~3DGS~模型~$G$~在视角~$v^t$~下的渲染图像，$G^{raw}$~是原始未保护模型。直觉上，在图像域施加约束绕过了软正则化的不稳定性和直接在异构高斯参数（位置、协方差、球谐系数、不透明度）上设定硬约束的困难，确保了跨视角一致且感知上可控的不可见性。

\paragraph{提升~PGD（Lifted PGD）：梯度截断与图像到高斯拟合的交替优化。}
求解上述约束优化问题的难点在于：优化变量是~3D~高斯参数，而约束定义在多视角的~2D~渲染空间——两个空间之间的映射关系由可微渲染器隐式定义。AdLift~设计了提升~PGD（L-PGD）来解耦这两个空间。考虑第~$k$~次迭代，对于训练视角~$v^t$，标准投影梯度下降的更新规则为~$G^{k+1} = \Pi_C(G^k - \alpha_G \nabla_G L_{adv}(R(G^k, v^t), y))$，其中~$\Pi_C$~是投影到可行集~$C = \{G : \|R(G, v^t) - R(G^{raw}, v^t)\|_p \leq \varepsilon\}$~的算子。由于~$C$~在~3D~参数空间中难以显式刻画，L-PGD~将更新分解为两个步骤。

第一步是梯度截断（gradient truncation）：虽然可行集~$C$~在参数空间隐式定义，但在渲染图像空间有显式刻画。设~$x^k = R(G^k, v^t)$~为当前渲染图像，$x^{raw} = R(G^{raw}, v^t)$~为原始渲染，L-PGD~在图像域截断梯度并投影：

\begin{equation}
x^{k+1} = \Pi_{B_p(x^{raw}, \varepsilon)}\!\bigl[x^k - \alpha \cdot \mathrm{sign}(\nabla_x L_{adv}(x^k))\bigr],
\end{equation}

其中~$B_p(x^{raw}, \varepsilon)$~是以~$x^{raw}$~为中心、半径为~$\varepsilon$~的~$\ell_p$~球，$\Pi_{B_p}$~将更新后的图像投影回该球内。这保证了截断后的渲染图像~$x^{k+1}$~严格满足不可见性硬约束，同时具有更小的对抗损失。

第二步是图像到高斯拟合（image-to-Gaussian fitting）：将截断后的图像~$x^{k+1}$~视为监督信号，通过最小化重建损失将~2D~约束传播回~3D~参数空间：

\begin{equation}
G^{k+1} = G^k - \beta \cdot \nabla_G L_{rec}(R(G^k, v^t), x^{k+1}),
\end{equation}

其中~$\beta$~是拟合步的学习率，$L_{rec}$~是光度重建损失（如~SSIM）。通过这一步，高斯参数~$G^{k+1}$~朝向假设的下一步高斯模型~$\tilde{G}^{k+1}$（满足~$R(\tilde{G}^{k+1}, v^t) = x^{k+1}$）优化，从而在保持可行集内约束的同时减小对抗损失。梯度截断与图像到高斯拟合交替执行~$K_p$~和~$K_l$~步，跨所有训练视角~$V^t$~联合渐进优化保护性高斯~$G^{prot}$。这种交替策略避免了在高斯参数上手动分配异构扰动预算的需要——每个中间更新都被约束为在视觉上与原始渲染不可区分，同时图像到高斯拟合将这些约束传播回~3D~空间。

\paragraph{多种对抗损失目标：抵抗~2D~和~3D~编辑。}
L-PGD~与多种对抗损失兼容。AdLift~设计了三种目标分别攻击编辑管线的不同组件：（1）无目标~VAE~攻击~$L_{VU} = -\|E(R(G, v^t)) - E(R(G^{raw}, v^t))\|_2^2$——鼓励渲染图像的~VAE~隐编码偏离原始表示，破坏编辑管线的潜在空间输入；（2）有目标~VAE~攻击（纹理损失）~$L_{VT} = \|E(R(G, v^t)) - E(x^{target})\|_2^2$——将隐编码对齐到无关目标纹理，注入干扰信号误导编辑过程；（3）SAM~攻击——针对~3DGS~局部编辑中广泛使用的~SAM~分割模型，强制预测的分割掩码偏离目标区域~$m^{target}$，从而误导局部编辑的区域定位，组合~BCE~和~Dice~损失。

\paragraph{实现细节与实验效果。}
AdLift~的初始化策略是将原始高斯~$G^{raw}$~完整复制一份作为可训练的保护性高斯（原始高斯固定不动）。在四个前向和~360°~场景上的实验表明，AdLift~在训练视角和未见视角上均能有效抵抗指令驱动的~2D~图像编辑（IP2P）和~3DGS~编辑（Instruct-GS2GS~全局编辑、DGE~局部编辑）。与基于软约束的~GuardSplat~对比，在匹配~PSNR~水平下~AdLift~的编辑抵抗更强；在匹配保护强度下~AdLift~的感知质量更高。此外，AdLift~继承了~2D~对抗扰动的迁移性——使用~SDv1.5-IP2P~作为代理编辑模型训练后，对未见过的~MagicBrush、SDEdit~和~SDXL-IP2P~编辑模型仍能降解编辑质量。在鲁棒性方面，AdLift~对~JPEG~压缩净化保持有效，对基于扩散的净化（DiffPure）虽保护强度有所降低但仍优于未保护基线。

\begin{takeaway}
对抗攻击与防御揭示了 3DGS 独特的安全攻防面：可微渲染管线使梯度攻击成为可能，对抗扰动可从 2D 图像域迁移到 3D 高斯参数空间，而 L-PGD 交替优化策略在保持视觉质量的同时实现编辑抵抗。
\end{takeaway}

%=====================================================================

\subsection{近期方法脉络：可靠性与硬件故障}

这一组近期工作延续 3DGS rendering，主要把改动集中在增量改进等环节。下面只保留相对既有路线的变化，便于和前文方法直接对照。

\paragraph{Single-Event Upsets in 3D Gaussian Splatting Rendering}
Single-Event Upsets in 3D Gaussian Splatting Rendering（2026.06）~\cite{2606_21791} 延续 3DGS rendering 在“可靠性与硬件故障”上的基本设定，主要改变是：GPU并行故障注入和Parallel Support Guard，分析不同字段、位宽和空间范围的single-event upset影响。3DGS 的显式高斯表示和可微渲染能力保持为基础。
其设计动机是：把3DGS安全边界扩展到硬件瞬态故障和静默数据损坏，给出bit-level criticality和空间影响分析。
具体而言，GPU-resident fault injection engine在多场景、多字段、fp32/fp16/bf16位位置上注入数百万次单比特翻转，并评估渲染退化。
适用边界是：故障模型聚焦单比特翻转；真实设备故障相关性和恢复策略仍需更多验证。
\section{隐写与信息隐藏}
\label{sec:21_3}

3DGS~隐写术旨在将秘密~3D~场景或其他信息隐藏在封面~3DGS~资产中，同时保持封面场景的渲染质量不受影响。与水印（嵌入版权标识）不同，隐写关注的是秘密信息的完整传输和不可检测性。3DGS~隐写的独特之处在于：由于~3DGS~的实时渲染能力，训练好的点云文件可能直接在线共享。因此，隐写不仅需要从渲染的~2D~视角中隐藏/提取信息，更需要在~3DGS~点云文件的内在结构中实现安全的隐藏和提取。

GS-Hider（2024.05）~\cite{2405_15118} 首次提出~3DGS~隐写框架，其核心创新是"耦合安全特征属性"的设计。本方法修改了~3DGS~的颜色表示与渲染管线——用高维耦合特征~$f_i$~替代原始~SH~系数~$h_i$~作为可学习颜色属性，同时将~RGB~光栅化替换为耦合特征光栅化——高斯的几何属性（位置~$\mu_i$、旋转~$q_i$、缩放~$s_i$、不透明度~$\alpha_i$）保持不变。

从基础隐写框架~GS-Hider~开始，后续工作从生成式隐写、知识蒸馏到频域加密和密钥安全等多个方向进行了拓展。

\subsection{隐写与信息隐藏方法}
\label{ss:21_3}

本节讨论三条技术路线：（1）耦合特征与对比学习隐写；（2）频域加密与双流隐写；（3）密钥安全与混合解耦隐写。

\subsubsection{耦合特征与对比学习隐写}

本路线介绍通过高维耦合特征替代 SH 系数来同时编码原始场景与隐藏场景的隐写方法。

\paragraph{动机与朴素方案排除}该方法首先排除了两种朴素方案：（1）简单地添加额外的~SH~系数~$h^*_i$~来拟合隐藏场景——不安全，因为攻击者可以从公开的点云文件中轻易检测并移除新增的属性维度；（2）在~3DGS~上附加解码器强制~2D~渲染图像揭示隐藏场景——每个高斯点具有特定的物理含义，联合优化解码器会干扰~3DGS~本身的优化，损害原始~3D~场景的保真度。理想的方案应该让隐藏信息与原始信息"你中有我、我中有你"，而非独立呈现。受经典图像隐写"编码器+解码器"范式的启发，GS-Hider~将~3DGS~的训练与渲染过程本身视为融合两个场景的"编码器"，另外设计轻量解码器网络来分离原始场景和隐藏场景。

\paragraph{耦合安全特征属性}GS-Hider~的核心设计是用高维耦合特征~$f_i \in \mathbb{R}^M$~（维度~$M$~可调）替代原始~SH~系数~$h_i \in \mathbb{R}^{48}$，使同一组参数同时编码原始场景和隐藏场景的颜色与纹理。与~$h_i$~具有固定物理含义不同，$f_i$~是灵活的高维表示，有两个关键优势：（1）信息融合是自适应学习的——两个场景的颜色和纹理在训练过程中自然交织，无需引入独立的编码器或额外参数；（2）安全性——窃取者无法区分特征向量中哪部分对应原始场景、哪部分对应隐藏消息。

\paragraph{耦合特征渲染管线}渲染阶段也做了相应调整。与原始~3DGS~将~SH~系数转换为~RGB~颜色再光栅化不同，GS-Hider~设计了耦合特征高斯光栅化器，直接将~$f_i$~渲染为高维特征图~$F_{coup} \in \mathbb{R}^{H \times W \times M}$（而非~RGB~图像）。特征图仍采用标准的~$\alpha$-混合公式：

\begin{equation}
    F_{coup}[p] = \sum_{i=1}^{N} f_i \, \alpha_i \prod_{j=1}^{i-1}(1-\alpha_j),
\end{equation}

其中~$N$~为覆盖像素~$p=(u,v)$~的高斯数量，$\alpha_i$~为第~$i$~个高斯的不透明度。与~3DGS~渲染~$c_i \in \mathbb{R}^3$~的~3~维颜色不同，特征图~$F_{coup}$~在更高的~$M$~维空间中渲染，提供了更大的信息容量。

\paragraph{双解码器分离两个场景}获得耦合特征图~$F_{coup}$~后，两个并行的轻量卷积解码器分别恢复两个场景：场景解码器~$D_s$~生成原始~RGB~场景~$I_{pred}$，消息解码器~$D_m$~提取隐藏消息场景~$M_{pred}$：

\begin{equation}
    I_{pred} = D_s(F_{coup}), \quad M_{pred} = D_m(F_{coup}).
\end{equation}

为保证实时渲染性能，两个解码器均仅由五层堆叠卷积加~ReLU~激活函数构成。部署时，场景解码器~$D_s$~与点云文件公开共享，而消息解码器~$D_m$~仅授权用户持有——作为特殊的隐私协议。

\paragraph{训练}训练需要原始场景~$\{I^{(n)}_{gt}\}_{n=1}^T$~和隐藏场景~$\{M^{(n)}_{gt}\}_{n=1}^T$~的一一对应视角（$T$~为训练视角数）。可学习参数包括高斯属性~$\{\mu_i, q_i, s_i, \alpha_i, f_i\}_{i=1}^K$~以及两个解码器~$D_s$~和~$D_m$。训练目标包含两个场景各自的重建损失：

\begin{equation}
    \mathcal{L}_{rgb} = (1-\lambda)\mathcal{L}_1(I_{pred}, I_{gt}) + \lambda \mathcal{L}_{SSIM}(I_{pred}, I_{gt}),
\end{equation}
\begin{equation}
    \mathcal{L}_{mes} = (1-\lambda)\mathcal{L}_1(M_{pred}, M_{gt}) + \lambda \mathcal{L}_{SSIM}(M_{pred}, M_{gt}),
\end{equation}

总损失为~$\mathcal{L}_{total} = \mathcal{L}_{rgb} + \beta \mathcal{L}_{mes}$，其中~$\beta$~平衡原始场景和消息场景的优化。训练过程中沿用~3DGS~的自适应密度控制策略（高斯分裂/克隆），以适应两个场景的共同需求。

GS-Hider~解决了隐写的基本框架问题，后续工作从不同应用场景进一步拓展了隐写技术。GaussianStego（2024.07）~\cite{2407_01301} 面向生成式~3DGS~提出首个优化框架，在保持渲染质量的同时从生成~3D~资产的渲染图像中准确提取隐藏信息。与锚点~GS-Hider~（在重建式~3DGS~中用耦合特征~$f_i$~替代~SH~系数~$h_i$，直接修改高斯参数同时编码两个场景）不同，GaussianStego~的核心区别在于操作域——它不修改~3DGS~参数本身，而是在~3DGS~生成模型的前馈过程中嵌入隐藏信息。具体地，GS-Hider~作用于已训练好的高斯点云的属性空间，而~GaussianStego~作用于~image-to-3D~生成网络~$\mathcal{F}_\theta$（基于~LGM~实现）的权重空间——通过微调生成网络的参数间接影响其输出高斯~$\mathcal{G} = \mathcal{F}_\theta(I)$。其核心挑战是：生成式~3DGS~的输出是随机初始化的高斯点云（而非从真实图像重建的），隐写嵌入需要在不破坏生成质量的前提下将信息编码到这些高斯的参数中。

\paragraph{隐藏信息嵌入}GaussianStego~的框架包含嵌入和解码两个阶段。嵌入阶段的关键是选择合适的视觉编码器：将~2D~隐藏图像投影到~3D~感知空间需要编码器具有语义连续的空间表示能力。作者选择~DINOv2~而非~CLIP~作为隐藏图像~$H$~的特征提取器——DINOv2~的密集预测特性使其空间表示在语义一致性和空间连续性方面均优于~CLIP~的对比学习特征，更适合将~2D~图像水印逆投影到~3D~表示空间。DINOv2~编码器提取隐藏图像的语义特征~$f_H$，随后通过交叉注意力将其注入~3D~生成模型的中间特征~$f_I$：

\begin{equation}
    f_I := \mathrm{softmax}\!\left(\frac{Q K^\top}{\sqrt{d}}\right) V,
\end{equation}

其中~$d$~为~Key/Query~特征的维度长度，$K = F_K(f_H)$，$V = F_V(f_H)$，$Q = F_Q(f_I)$。此处~$F_K, F_V$~为作用于隐藏图像特征的线性变换层，$F_Q$~为作用于~3D~生成中间特征的线性变换层——隐藏信息特征提供~Key~和~Value~决定"注入什么内容"，3D~特征提供~Query~决定"从哪里注入"，实现隐藏信息对~3D~特征的自适应加权调制。这种早期注入策略将隐藏图像的空间细节融入生成过程而非仅在输出端操作，使隐藏信息能通过网络传播影响~3D~表示的更新。

\paragraph{对比水印机制}解码阶段的关键创新是对比水印（Contrastive Watermarking）机制。给定检查视角~$P_C$~和用户自定义的无意义常量图像~$\emptyset$，对比损失定义为一对互补约束：

\begin{equation}
    \mathcal{L}_{dec^+} = |F_\omega(\mathcal{G}(P_C)) - H|, \quad \mathcal{L}_{dec^-} = |F_\omega(\mathcal{G}_0(P_C)) - \emptyset|,
\end{equation}

其中~$\mathcal{G}$~为含隐写信息的生成高斯，$\mathcal{G}_0$~为原始生成网络~$\mathcal{F}_{\theta_0}$~产生的不含隐写信息的高斯，$F_\omega$~为~U-Net~解码器。$\mathcal{L}_{dec^+}$~确保解码器从含隐写信息的渲染中恢复隐藏图像~$H$；$\mathcal{L}_{dec^-}$~确保不含隐写信息的原始渲染不会触发误检——解码器输出预定义的无意义常量图像~$\emptyset$~而非任何看似合理的图案，从而有效避免假阳性。

\paragraph{自适应梯度协调}隐写嵌入面临渲染质量与信息隐藏之间的固有冲突。GaussianStego~发现，简单惩罚含隐写高斯与原始高斯之间所有参数的偏差会过度限制生成网络的隐写能力。关键洞察是：生成网络的权重对渲染质量的贡献具有强稀疏性，且渲染和隐写两个目标对网络权重的梯度需求在不同权重组上方向可能一致也可能冲突。自适应梯度协调（AGH，Adaptive Gradient Harmonisation）策略仅对梯度方向一致的权重组执行隐写更新：

\begin{equation}
    M = \mathbb{I}\!\left(\cos\!\left(\nabla \mathcal{L}_{rgb},\, \nabla(\mathcal{L}_{dec^+} + \mathcal{L}_{dec^-})\right) > 0\right),
\end{equation}

其中~$\mathbb{I}$~为指示函数，$\cos(\cdot,\cdot)$~计算两个梯度分量的余弦相似度。梯度掩码~$M$~通过~Hadamard~积作用于总损失的梯度：$\nabla\mathcal{L} \odot M$。具体地，余弦相似度大于零的权重组（两个目标方向一致）正常更新，小于零的权重组（方向冲突）被屏蔽，最小化隐写学习对渲染视觉质量的影响。总体训练损失为~$\mathcal{L} = \lambda_1 \mathcal{L}_{dec^+} + \lambda_2 \mathcal{L}_{dec^-} + \lambda_3 \mathcal{L}_{rgb}$，其中~$\mathcal{L}_{rgb} = |\mathcal{G}(P) - \mathcal{G}_0(P)|$~约束含隐写高斯与原始高斯的渲染一致性。

\paragraph{实现与泛化}GaussianStego~基于~LGM~生成模型实现，为每个水印微调独立的~LoRA（Low-Rank Adaptation，低秩适配）~适配器（而非微调整个生成网络），在约~100~个~Objaverse~对象上训练~30~轮（约~20~分钟）即可泛化到未见过的对象。框架可扩展至多模态嵌入（文本、QR~码、音频、视频），通过为每种模态设计专用解码器分支实现。实验表明文本嵌入恢复精度（ACC）、QR~码恢复质量（SSIM）和音频恢复质量（PSNR）均达到实用水平。消融实验验证了~DINOv2~编码器、交叉注意力和梯度协调三个组件各自的必要性。

从生成式隐写~GaussianStego~过渡到后处理隐写~ConcealGS，研究路线从修改生成过程转向直接对已有~3DGS~资产进行操作。

GaussianStego~面向生成式~3DGS~的生成过程嵌入，而~ConcealGS~则从知识蒸馏的角度提供了一种互补的后处理隐写思路。ConcealGS（2025.01）~\cite{2501_03605} 的核心区别在于操作时机——GaussianStego~在~3DGS~生成模型的前馈过程中嵌入隐藏信息（通过修改生成网络权重间接影响输出高斯），而~ConcealGS~直接对已训练好的~3DGS~资产进行后处理嵌入，无需访问原始训练数据或生成管线。这样，ConcealGS~适用于更广泛的场景：任何已有的~3DGS~资产都可以直接添加隐式信息，而不依赖特定的生成架构。

\paragraph{知识蒸馏框架}ConcealGS~的总体框架包含教师模型、学生模型和通用解码器三个组件。教师模型为预训练的~3DGS~$G_o$，学生模型~$G_w$~为嵌入了隐式信息的~3DGS。直接在~3DGS~中嵌入隐式信息会显著降低渲染质量（类似~StegaNeRF~在~NeRF~中的退化问题），因此~ConcealGS~引入知识蒸馏正则化约束嵌入过程：

\begin{equation}
    L_{kd} = |G_w(P_m) - G_o(P_m)|,
\end{equation}

其中~$G_w(P_m)$~和~$G_o(P_m)$~分别为学生模型和教师模型在相机位置~$P_m$~处的渲染图像。该损失确保含隐式信息的学生模型在渲染结果上与教师模型保持一致，隐式信息的嵌入不会导致可感知的视觉降质。解码器~$F$（基于~U-Net~实现）负责从渲染图像中提取隐式信息，其训练目标由正负对比损失共同驱动。

\paragraph{对比式信息嵌入与一致性保持}ConcealGS~采用检查视角（check view）$P_c$~和正常视角（normal view）$P_n$~的双视角机制实现信息嵌入与渲染保真的平衡。检查视角~$P_c$~是所有者预先选定的特定视角，用于在侵权检测时提取隐式信息；正常视角~$P_n \in \{P_{n_j}\}_{j=1}^M$~是用户常规浏览时使用的视角集合。解码器~$F$~在检查视角上的正对比损失确保隐式信息可被恢复：

\begin{equation}
    L_d^+ = |F(G_w(P_c)) - I|,
\end{equation}

其中~$I$~为目标嵌入的隐式信息（如~$64 \times 64$~分辨率的隐藏图像），$L_d^+$~度量解码器预测与隐式信息真值之间的差异。在正常视角上的负对比损失确保解码器不干扰常规渲染：

\begin{equation}
    L_d^- = |F(G_o(P_n)) - G_o(P_n)|,
\end{equation}

$L_d^-$~训练解码器在正常视角输入上表现为恒等映射——输出与输入尽可能一致，仅引入极小的扰动或噪声。这种对比式设计使解码器在不同视角下呈现截然不同的行为：检查视角恢复隐藏信息，正常视角保持原始渲染不变。

\paragraph{梯度引导优化}隐写嵌入面临信息提取（$L_d^+$）与渲染一致性（$L_d^-$）之间的固有冲突——两个损失函数对解码器参数的梯度方向可能不一致甚至相反，直接联合优化会导致解码器在两个目标间振荡。ConcealGS~提出梯度引导优化策略，逐层动态调整解码器的更新强度。对于~$N$~层解码器~$F = \{f_i\}_{i=1}^N$，计算每层~$f_i$~对两个损失梯度的余弦相似度：

\begin{equation}
    s_i = \frac{\nabla_{f_i} L_d^+ \cdot \nabla_{f_i} L_d^-}{\|\nabla_{f_i} L_d^+\| \cdot \|\nabla_{f_i} L_d^-\|}, \quad w_i = \sigma(s_i),
\end{equation}

其中~$s_i$~衡量第~$i$~层在两个任务上的梯度方向一致性，$\sigma(\cdot)$~为~sigmoid~函数将其映射为~$(0, 1)$~区间的梯度权重~$w_i$。优化时每层的梯度乘以对应权重：$\nabla_{f_i} L \leftarrow w_i \cdot \nabla_{f_i} L$。该策略的物理含义清晰——梯度方向相似的层（$s_i \approx 1$，即两个任务在该层上一致推进）保持正常更新；梯度方向相反的层（$s_i \approx -1$，即两个任务在该层上相互冲突）被抑制更新，避免一个目标的优化损害另一个目标。这种动态协调机制在不同解码器层之间自适应地平衡信息嵌入与渲染保真。

\paragraph{实验验证}ConcealGS~在~NeRF-Synthetic（lego, drums, chair, ship）和~T\&T（playroom, drjohnson, truck, train）两个数据集上评估。与~LSB、DeepStega~和~StegaNeRF~三种基线方法相比，ConcealGS~在场景渲染质量上~PSNR~提升~0.89~dB，在隐式信息恢复质量上~PSNR~提升~5.20~dB。相比未嵌入任何隐式信息的标准~3DGS~模型，ConcealGS~仅引入~0.17~dB~的~PSNR~下降和~0.0013~的~SSIM~下降，几乎不影响渲染质量。鲁棒性分析表明，在高斯模糊和~JPEG~压缩等常见图像处理操作下，隐藏信息的~SSIM~仅下降约~0.02，验证了方法在真实环境中的实用性。消融实验分别验证了解码器、一致性策略和梯度引导优化三个组件各自的贡献。

\subsubsection{频域加密与双流隐写}

本路线介绍通过频域加密和双流高斯共存实现无需修改渲染管线的隐写方法。前面三种方法都通过修改高斯参数来嵌入隐藏信息——GS-Hider~用耦合特征替代~SH~系数，GaussianStego~在生成过程中注入隐藏特征，ConcealGS~对已训练模型做后处理蒸馏。这导致一个共同的实用性缺陷：修改后的~3DGS~资产无法直接兼容原始~3DGS~渲染引擎和下游任务管线。Splats in Splats（2024.12）~\cite{2412_03121} 提出了一种根本性不同的策略——完全不修改任何高斯属性，也不改变渲染管线，而是让两组独立的高斯点云共存在同一空间中，实现"高斯中藏高斯"的隐写范式。

从参数修改式的隐写过渡到不修改参数的"高斯中藏高斯"范式，隐写技术向更高的兼容性和安全性演进。

\paragraph{动机与核心洞察}Splats in Splats~的设计源自对~SH~系数物理含义的深入分析。SH~的阶数~$l$（band index）决定了频率特性：低阶~SH~承载低频成分（主要视觉特征），高阶~SH~承载高频成分（细节纹理）。在大多数场景中，高频反射仅占很小比例，高阶~SH~系数的数值贡献微乎其微，存在显著的信息冗余。这意味着可以在高阶~SH~的冗余空间中嵌入额外信息，而不会对视觉效果产生可感知的影响。基于这一洞察，Splats in Splats~将隐藏场景的信息编码到封面场景高阶~SH~系数的低位比特中——这是与~GS-Hider~的本质区别：GS-Hider~用高维耦合特征~$f_i \in \mathbb{R}^M$~完全替代原始~SH~系数~$h_i \in \mathbb{R}^{48}$，修改了每个高斯的属性结构和渲染方式；Splats in Splats~则保留原始~SH~系数不变，仅在高阶分量的低位比特中嵌入隐藏信息，封面场景的点云文件结构、属性格式和渲染管线完全不变。

\paragraph{重要性分级~SH~系数加密}设封面场景的~SH~系数为~$c_i \in \mathbb{R}^{n \times 3}$，隐藏场景的~SH~系数为~$c_i' \in \mathbb{R}^{n \times 3}$，其中~$c_{i,j}$~表示第~$i$~个高斯的第~$j$~个~SH~分量（$0 \leq j \leq n-1$），其阶数为~$\lfloor \sqrt{j} \rfloor$。阶数越低，该系数对视觉效果越重要。加密策略遵循"重要性反转"原则：将隐藏场景中最重要的低阶~$c_{i,j}'$~嵌入到封面场景中最不重要的高阶~$c_{i,j}$~中。具体地，首先将~SH~系数缩放至整数并转为二进制表示，然后对~$c_{i,j}$~按其阶数分配保留位数~$\Delta_j$~，截断低位：$\tilde{c}_{i,j} = c_{i,j}~\&~((1 \ll (k + \Delta_j)) - 1)$，其中~$k$~为对应~0~阶系数的移位长度，$\ll$~和~$\&$~分别为左移和按位与操作。随后将~$c_{i,j}'$~按重要性反转索引（$n-1-j$）移位后与~$\tilde{c}_{i,j}$~异或融合：

\begin{equation}
    c_{i,j}^W = \tilde{c}_{i,j} \oplus (c_{i,n-1-j}' \gg (\Delta - (k + \Delta_j))),
\end{equation}

其中~$\Delta$~为最大比特长度，$\gg$~为右移操作。索引~$n-1-j$~确保低阶（高重要性）的隐藏系数被嵌入到高阶（低重要性）的封面系数中。提取时，版权持有者通过逆运算恢复隐藏系数：$c_{i,j}' = c_{i,n-1-j}^W~\&~((1 \ll (k + \Delta_{n-1-j})) - 1)$。这种位级加密确保嵌入操作对封面场景的视觉质量几乎无影响，同时对噪声攻击表现出强鲁棒性——重要性分级策略保证即使高阶分量受到噪声干扰，低阶（最关键）的隐藏信息仍可正确恢复。

\paragraph{卷积自编码器不透明度映射}3D~场景由几何结构和不透明度共同决定。SH~加密处理了外观信息（颜色与纹理），但不透明度~$\alpha_i$~的隐藏需要另一条路径。直接在原始不透明度上叠加隐藏不透明度会破坏封面场景的几何结构，因此~Splats in Splats~采用卷积自编码器学习从封面不透明度到隐藏不透明度的非线性映射。首先设定阈值~$\tau$~筛选出隐藏场景中重要的不透明度索引集合~$I = \{i \mid \alpha_i' > \tau,~i \in \{1,\ldots,N\}\}$，仅保留对隐藏场景重建有显著贡献的高斯点（索引~$I$~之外的隐藏不透明度设为零），同时存储这些高斯的~3D~坐标~$x_I$~作为私密信息。基于封面不透明度~$\alpha_I$~与隐藏不透明度~$\alpha_I'$~在许多位置呈现互补关系的观察，自编码器以~$1 - \alpha_I$~为输入，经过编码器~$E$~和解码器~$D$~的编解码后输出隐藏不透明度估计：

\begin{equation}
    W_p^* = \arg\min_{E,D}~\mathcal{L}_{MSE}(D(E(1 - \alpha_I)),~\alpha_I').
\end{equation}

自编码器由若干轻量卷积层和反卷积层构成，确保实时推理。训练好的模型参数~$W_p^*$~作为资产持有者的私钥——只有拥有~$W_p^*$~和索引~$I$~的授权用户才能从封面不透明度中恢复隐藏场景的几何结构。

\paragraph{整体隐写与提取流程}整个框架分为嵌入和提取两个阶段。嵌入阶段：（1）分别训练封面场景和秘密场景的~3DGS，获取两组独立的高斯点云；（2）对隐藏场景的~SH~系数和不透明度执行上述加密与映射，将隐藏信息编码进封面场景的参数中；（3）输出的~3DGS~资产保持与原始~3DGS~完全相同的属性格式（$\{\mu_i, q_i, s_i, \alpha_i, c_i^W\}$）和渲染管线，可直接在任何标准~3DGS~渲染引擎中使用。提取阶段：版权持有者利用私钥（索引集合~$I$~和自编码器参数~$W_p^*$~）从公开资产中解密隐藏~SH~系数~$c_{i,j}'$~并估计隐藏不透明度~$\alpha_I'$，配合封面场景的几何属性（位置、旋转、尺度）恢复出完整的隐藏~3DGS~场景。这种设计确保了"可用性"——普通用户完全感知不到隐藏内容的存在，版权持有者则可在需要时提取完整~3D~场景作为权属凭证。

\paragraph{实验验证}在~Mip-NeRF 360~的~9~个场景上（隐藏场景与封面场景的对应关系与~GS-Hider~一致），Splats in Splats~在保真度和效率上均显著优于~GS-Hider~和~SecureGS~。封面场景的新视角合成~PSNR~提升~5.31\%，同时渲染速度达到~100 FPS——是使用神经网络的~GS-Hider~的~$3\times$，训练时间也仅为现有方法的一半。这一效率优势源于框架完全不依赖任何神经网络进行渲染：嵌入后的~3DGS~资产通过标准~3DGS~光栅化器直接渲染，无需额外的解码器推理。鲁棒性方面，在~25\%~序列化剪枝（删除低不透明度高斯）下~PSNR~仅下降~$0.002$~（GS-Hider~下降~$0.012$，SecureGS~退化~$5.171$），在~25\%~随机剪枝下~PSNR~仅下降~$0.09$~（GS-Hider~下降~$1.260$，SecureGS~下降~$1.435$），展现出卓越的抗攻击能力。

其升级版~Splats in Splats++（2026.04）~\cite{2604_15862} 从鲁棒性和泛化性两个维度显著增强原版。与原版相比，核心改进体现在三个方面：保留重要性分级~SH~系数加密的核心框架但重新设计了双流优化策略、用多分辨率哈希网格编码器替代卷积自编码器实现不透明度映射、以及引入梯度门控不透明度一致性损失强化场景融合。

\paragraph{双流隐写架构}与原版将两个独立训练的~3DGS~场景通过后处理融合不同，Splats in Splats++~采用双流联合优化架构。每个高斯基元~$i$~关联两组独立的属性：场景属性~$\{\alpha_i^S, c_i^S\}$~和消息属性~$\{\alpha_i^M, c_i^M\}$，而几何属性（位置~$x_i$、旋转~$r_i$、尺度~$s_i$）在两个表示之间共享。训练目标同时优化场景重建损失~$\mathcal{L}_{scene}$~和消息重建损失~$\mathcal{L}_{mes}$：

\begin{equation}
    \mathcal{L}_{recon} = \mathcal{L}_{scene} + \lambda^M \mathcal{L}_{mes},
\end{equation}

其中每个损失均为~$\mathcal{L}_1$~和~$\mathcal{L}_{SSIM}$~的加权组合，$\lambda^M$~平衡场景保真度与消息恢复精度（实验中设为~1.0）。采用两组独立不透明度而非单一共享不透明度的设计，有助于缓解几何歧义和潜在信息泄漏——如果两个场景共享同一不透明度场，封面场景渲染时的不透明度梯度可能间接泄露隐藏场景的几何结构。

\paragraph{重要性分级~SH~系数加密}原版~Splats in Splats~已提出重要性反转加密策略，Splats in Splats++~沿用了这一核心设计并进一步阐明了其理论基础。加密策略将隐藏场景的低阶~SH~系数嵌入到原始场景的高阶~SH~系数中，利用~SH~的物理洞察：低阶~SH~表示低频成分（承载主要视觉特征），高阶~SH~表示高频成分（信息冗余且对人眼不敏感）。将~SH~系数缩放至~$\Delta$~位整数表示~$\{0, 1\}^\Delta$~后，首先按阶数分配保留位数~$\Delta_j$~清空封面系数的低位嵌入空间：$\tilde{c}_{i,j} = c_{i,j}~\&~((1 \ll (k + \Delta_j)) - 1)$，其中~$k$~为~0~阶系数的基础移位长度（实验中设为~17）。随后执行重要性反转位融合：

\begin{equation}
    c_{i,j}^W = \tilde{c}_{i,j} \oplus (c_{i,n-1-j}' \gg (\Delta - (k + \Delta_j))),
\end{equation}

索引~$n-1-j$~确保隐藏场景中最重要的低阶系数被嵌入到封面场景中最不重要的高阶系数中，最大化嵌入容量同时保持视觉不可感知。提取时，版权持有者通过位平面隔离恢复隐藏系数：$\hat{c}_{i,j}^M = c_{i,n-1-j}^W~\&~((1 \ll (k + \Delta_{n-1-j})) - 1)$。消融实验表明，重要性分级加密在高斯噪声（$0.0005$--$0.01$）下相比均匀加密~(AVG)~具有显著优势——随着噪声水平增大，均匀加密导致恢复质量急剧退化，而分级加密则保持稳定的高质量恢复，因为最关键的隐藏信息被锚定在视觉冗余度最高的分量中。

\paragraph{哈希网格引导不透明度映射}原版采用卷积自编码器学习封面不透明度到隐藏不透明度的映射，但自编码器缺乏空间感知能力——它将不透明度视为独立标量处理，忽略了高斯基元之间的空间相关性。Splats in Splats++~用多分辨率哈希网格编码器替代卷积自编码器，将离散的属性映射投影到连续的潜在流形上，使恢复的隐藏属性继承原始场景的几何一致性。具体地，对于归一化坐标~$x_i \in [0, 1]^3$~的高斯基元，建立~$L$~层分辨率的层次结构以捕获多尺度几何特征。每层分辨率~$R_l$~由几何级数确定：

\begin{equation}
    R_l = R_{min} \cdot b^l, \quad b = \exp\!\Big(\frac{\ln R_{max} - \ln R_{min}}{L-1}\Big),
\end{equation}

实验中~$L=16$、$R_{min}=16$、$R_{max}=1024$。在每层中，坐标~$x_i$~缩放至~$R_l$~以定位其所在体素单元的~8~个整数顶点~$V_l$，通过固定空间哈希函数~$\mathcal{H}$~从大小为~$T = 2^{21}$~的哈希表~$T_l$~中检索可学习嵌入向量：

\begin{equation}
    h(v, T) = \Big(\bigoplus_{d=1}^{3} v_d \pi_d\Big) \mod T, \quad v \in V_l,
\end{equation}

其中~$\pi_d$~为不同的大素数。获取顶点嵌入后，基于~$x_i$~在单元内的相对位置执行三线性插值，得到层特征~$\phi_{i,l} \in \mathbb{R}^F$（$F=4$）。多层特征拼接形成多尺度空间描述符~$h_i = [\phi_{i,0}, \ldots, \phi_{i,L-1}] \in \mathbb{R}^{LF}$。为桥接空间几何与视觉外观，将~$h_i$~与原始场景不透明度~$\alpha_i^S$~和主~SH~系数~$c_i^S$~拼接后输入映射网络~$\mathcal{M}$（由哈希表嵌入~$T$~和~MLP~权重~$W$~参数化），推断隐藏不透明度：

\begin{equation}
    \mathcal{L}_{op}(T, W) = \sum_{i=1}^{N} \|\mathcal{M}(h_i \oplus \alpha_i^S \oplus c_i^S;\, W) - \alpha_i^M\|^2,
\end{equation}

其中~$\oplus$~为拼接操作。映射模块联合优化~4000~轮（学习率~$5 \times 10^{-3}$），使~MLP~通过神经插值推理隐藏不透明度值。消融实验对比了三种方案：（a）直接~MLP~预测、（b）不透明度掩码+CNN（原版方案）和~(c)~哈希网格引导映射——前两种方案即使在添加一致性损失后仍无法建立两个场景间的精确几何对应关系，而哈希网格方案始终产生高质量、几何一致的映射。

\paragraph{梯度门控不透明度一致性损失}采用两组独立不透明度虽然缓解了几何歧义，但若不加约束，优化过程可能驱使两个不透明度场过度分离，导致场景与消息在几何上解耦，增加被对抗攻击剥离的风险。为此引入可见性与梯度感知的不透明度一致性正则化，选择性地仅在稳定且信息丰富的区域强制两个场景的几何一致性。将不透明度建模为~Bernoulli~分布的参数，用对称~KL~散度度量场景不透明度~$\alpha_i^S$~和消息不透明度~$\alpha_i^M$~之间的差异，并通过两个因子调制约束强度：（1）可见性权重~$v_i$~由体渲染过程推导，反映高斯对最终图像的贡献程度；（2）梯度门控：

\begin{equation}
    g_i = \exp\!\big(-\|\nabla_{\alpha_i^S} \mathcal{L}_{recon}\|\big),
\end{equation}

该门控机制在活跃优化区域（$\nabla_{\alpha_i^S} \mathcal{L}_{recon}$~较大）放松约束，允许外观优化自由进行；在稳定区域（梯度接近零）强制对齐，确保两个场景共享一致的几何支撑。最终的一致性损失为~$\mathcal{L}_{cons} = \mathbb{E}_i[v_i \cdot g_i \cdot D_{KL}^{sym}(\alpha_i^S \| \alpha_i^M)]$，总训练目标为：

\begin{equation}
    \mathcal{L}_{total} = \mathcal{L}_{recon} + \lambda \mathcal{L}_{cons},
\end{equation}

其中~$\lambda = 0.02$。这种密度门控机制有效耦合了场景与消息的几何结构，使隐藏信息难以被结构攻击移除。消融实验表明，关闭一致性损失后，方法在多种对抗攻击下性能显著退化；而启用后，在所有攻击评估中始终达到最优。

\paragraph{实验验证与泛化能力}在~NeRF~合成数据集、Tanks and Temples~和~Mip-NeRF 360~三个数据集上评估，Splats in Splats++~在保真度和效率上均达到最优。与~GS-Hider~和~SecureGS~相比，消息保真度提升~6.28 dB（在所有数据集上的平均~Message PSNR），渲染速度保持~100 FPS（是~GS-Hider~的~$3\times$），训练时间减半。鲁棒性方面引入了鲁棒性评分~$S_R$——定义为场景相对退化率与消息相对退化率之差，值越接近零或正值表示隐藏信息比场景本身更抗攻击。在~Opacity Pruning~攻击下，消息~PSNR~为~27.252 dB（$S_R = -1.10$），远优于次优方法（$S_R = -5.49$）；在更具攻击性的~GSPure（专门针对~3DGS~隐写的剪枝方法）下，消息~PSNR~保持~22.185 dB，而次优方法的鲁棒性评分骤降至~$-24.51$。此外，方法天然兼容所有~3DGS~渲染引擎和下游任务——原始与隐藏场景可直接用于网格提取（SuGaR）和~3D~分割（LangSplat）。泛化性方面：（1）2D~图像嵌入——将图像嵌入视为仅含单一视角的~3D~场景嵌入的退化情形，嵌入图像可高质量恢复且对封面场景几乎无影响；（2）4D~动态场景隐写——将框架扩展至~4DGS~变形场（先训练静态场景再微调变形场），在~D-NeRF~和~HyperNeRF~数据集上保持最高~NVS~质量，渲染结果时间连贯且视觉保真，而~SecureGS~在动态区域出现明显伪影，GS-Hider~出现纹理和色彩失真。

\subsubsection{密钥安全与混合解耦隐写}

本路线介绍引入密钥控制机制和端到端联合优化以提升隐写安全性的方法。"高斯中藏高斯"的范式实现了不修改属性的秘密嵌入，但安全性依赖于秘密场景的隐蔽性——一旦攻击者知道隐写的存在，仍有可能检测到异常。更关键的是，无论是~Splats in Splats~还是~Splats in Splats++，其隐写框架本身不具备访问控制机制：任何持有~3DGS~资产的用户都可以尝试提取隐藏内容。此外，两种方法均采用~SH~加密与不透明度映射分离的非端到端管线，且单一封面场景只能嵌入一个秘密场景，限制了多秘密隐藏的灵活性。为此，引入密钥机制成为提升隐写安全性的自然延伸。

从频域加密隐写过渡到密钥安全隐写，安全性逐步从"隐蔽性"升级到"访问控制"。

\paragraph{端到端密钥安全框架}KeySS（2025.03）~\cite{2503_07191} 提出密钥安全~3D~隐写框架（Key-Secured 3D Steganography），从架构层面重新设计隐写范式。与~Splats in Splats/Splats in Splats++~的分离式管线（分别处理~SH~加密和不透明度映射）不同，KeySS~采用端到端学习框架，从零开始联合优化封面~3DGS~模型~$G_{cover}$~和解码器~$D$~，通过解码器直接学习从封面高斯到秘密高斯的完整变换，同时保持标准~3DGS~特征格式和渲染管线的完全兼容。给定~$S+1$~组真值~2D~图像（一组封面场景~$\{I_{gt\_cover}^i\}_{i=1}^{M_{cover}}$~和~$S$~组秘密场景~$\{I_{gt\_secret^{s}}^i\}_{i=1}^{M_{secret}}$），框架同时优化封面重建和秘密恢复：$L_{cover} = (1-\lambda_{cover})L_1(I_{pred\_cover}, I_{gt\_cover}) + \lambda_{cover}\mathcal{L}_{SSIM}(I_{pred\_cover}, I_{gt\_cover})$，秘密损失~$L_{secret}^s$~同理。核心区别在于密钥控制的解码架构——解码器~$D$~以用户密钥~$k^s$~为条件：$G_{secret}^s = D(G_{cover}, k^s)$，使授权用户通过正确密钥恢复秘密场景，而未授权用户即使获取了点云文件也无法提取隐藏内容。

\paragraph{密钥编码与安全训练}密钥通过预训练~CLIP~文本编码器（2021）~\cite{clip_radford2021learning}（ViT-L/14）编码为语义嵌入：用户密钥~$k$~首先被分词，然后经过~transformer~编码器~$E$~和平均池化操作得到最终密钥嵌入~$k = \mathrm{AvgPool}(E(k))$。密钥空间限定为~16~字符的字母数字字符串（大小写字母），在安全强度和计算可行性之间取得平衡。CLIP~编码器的优势在于能将多样化的文本输入映射到语义丰富且结构化的嵌入空间，使密钥空间的语义覆盖足够宽广。密钥嵌入与归一化的~3DGS~特征拼接后输入解码器。训练中同时优化两种密钥场景：（1）正确密钥场景——使用对应秘密场景的密钥~$k^s$~，解码器恢复秘密场景（损失~$L_{secret}^s$）；（2）随机错误密钥场景——使用训练中未见的随机密钥，强制解码器重建原始封面场景：

\begin{equation}
    L_{incorrect} = (1-\lambda_{cover}')L_1(I_{pred\_incorrect}, I_{gt\_cover}) + \lambda_{cover}'\mathcal{L}_{SSIM}(I_{pred\_incorrect}, I_{gt\_cover}),
\end{equation}

其中~$I_{pred\_incorrect}$~为错误密钥下的渲染结果。这一约束确保错误密钥下解码器输出封面场景本身，不泄露任何隐藏信息。多秘密扩展仅需为每个额外秘密增加一组损失项~$L_{secret}^s$，无需引入额外架构修改——解码器根据不同密钥~$k^{s_1}, k^{s_2}, \ldots$~自然切换到对应的秘密场景，体现了密钥机制的可扩展性。

\paragraph{共享公共分支与特征特定分支解码器}解码器由共享公共分支~$D_{common}$~和五个特征特定分支组成，系统化地量化各高斯属性对隐写的独立贡献。所有~3DGS~特征与密钥嵌入归一化后拼接为输入向量~$f = \mathrm{concat}(\alpha, r, s, \mu, c, k)$，公共分支提取综合特征表示~$h = D_{common}(f)$。随后，公共特征~$h$~与封面高斯的原始属性分别通过五个特征特定分支，得到更新后的属性：

\begin{equation}
    [\alpha^*, r^*, s^*, \mu^*, c^*] = [\alpha, r, s, \mu, c] + \delta^5 \odot [\Delta_{op}(h), \Delta_{ro}(h), \Delta_{sc}(h), \Delta_{po}(h), \Delta_{sh}(h)],
\end{equation}

其中~$\Delta_{op}$、$\Delta_{ro}$、$\Delta_{sc}$、$\Delta_{po}$、$\Delta_{sh}$~分别对应不透明度、旋转、缩放、位置和~SH~特征的特征特定分支，$\delta^5$~为二值选择向量控制哪些属性参与更新，$\odot$~为~Hadamard~积。这一设计的核心洞察来自系统的特征贡献分析：高斯特征对秘密隐藏的贡献不均等。不透明度对秘密隐藏贡献最大——这符合直觉，因为不透明度直接控制每个高斯的可见性；但仅依赖不透明度存在安全隐患：攻击者可通过检测和恢复零值不透明度属性暴露隐藏信息。相比之下，SH~特征因其高维性和复杂性反而最难学习，对隐写过程的辅助有限。消融实验表明，四特征组合（不透明度+旋转+缩放+位置）在重建保真度和隐写不可感知性之间取得最优平衡——该组合的低不透明度区域分布更均匀，3D-Sinkhorn~距离更低，统计不可检测性显著优于仅更新不透明度的方案。

\paragraph{3D-Sinkhorn~距离评估}传统~2D~隐写分析（如~StegExpose）基于渲染图像的像素级统计检测，无法有效捕获~3DGS~参数空间中的分布异常——不同的特征更新策略在~StegExpose~下产生几乎相同的~AUC~值，无法区分其安全性差异。为此，KeySS~首次提出~3D-Sinkhorn~距离评估，在~3DGS~参数空间中直接量化隐写的不可检测性。具体地，将真值封面场景和隐写封面场景的~3D~高斯归一化后投影为特征特定直方图：$g_i = \mathrm{hist}(f_i)$，$g_i^{gt} = \mathrm{hist}(f_i^{gt})$，然后用~Sinkhorn~距离（基于熵正则化最优传输的分布度量）分别计算各属性维度上的分布差异：$d = \sum_i \mathrm{Sinkhorn}(g_i, g_i^{gt})$。Sinkhorn~距离相比~Wasserstein~距离通过熵正则化实现了计算效率与几何对应保持的理想平衡，对~3D~分析至关重要。为系统化评估，引入综合评分平衡重建质量和统计不可感知性：$\mathcal{S} = (\mathcal{M}_{cover} + \mathcal{M}_{secret}) \cdot (1 - d)$，其中~$\mathcal{M}$~为保真度指标。该评分同时奖励高质量重建和低分布差异，激励方法在两个目标之间寻找最优点。

\paragraph{组合训练策略}KeySS~采用三项组合策略充分利用双场景信息：（1）合并~SfM~点云初始化——将封面场景和秘密场景的~SfM~点云合并作为~3DGS~初始化点，保留双场景的完整空间信息。消融实验证明这一策略至关重要：仅使用封面或秘密场景的~SfM~点会导致性能偏向对应场景；（2）合并相机姿态——对于某一场景独有的相机视角，使用预训练模型生成对应视角的真值图像（用封面~3DGS~渲染秘密视角、用秘密~3DGS~渲染封面视角），丰富训练样本多样性，实现双场景间的平衡性能；（3）合并致密化——在反向传播中利用封面和秘密场景的视图空间位置梯度联合指导高斯的分裂和克隆，细化~$G_{cover}$~的表示精度，在保持封面完整性的同时提升隐藏信息的准确度。整体损失为：$L = \lambda_0 L_{cover} + \sum_{s=1}^{S} \lambda^s L_{secret}^s + \lambda_{inc} L_{incorrect}$，其中~$\lambda_0 = \lambda^s = 0.5$，$\lambda_{inc} = 0.01$。

\paragraph{实验验证}在~Mip-NeRF 360~的~9~个场景和~Deep Blending~的~1~个场景上（按~GS-Hider~的协议配对为~9~个封面-秘密对），KeySS~在保真度上显著优于~GS-Hider~：尽管~KeySS~的真值~3DGS~封面上限较低（训练设置差异），封面重建仅下降~0.511 dB（GS-Hider~下降~1.388 dB），秘密恢复~PSNR~优于基线~4.9\%。多秘密隐藏时（一个封面嵌入两个秘密），仅需增加一组损失项，封面保真度和秘密恢复质量均与单秘密场景可比，密钥可控解码机制根据不同密钥输入准确切换到对应秘密场景。安全性方面，在~StegExpose~隐写分析下，KeySS~的~ROC~曲线接近对角线（随机猜测水平），显著优于现有方法，且多秘密嵌入后安全性不降。错误密钥攻击评估中，随机生成训练中未见的错误密钥，解码器稳定输出封面场景（高~PSNR~对封面、低~PSNR~对秘密），验证了~$L_{incorrect}$~损失的有效性。

KeySS~通过密钥机制增强访问控制，而~SecureGS~则从另一个角度利用了锚点~3DGS~的结构化特性来提升隐写的安全性和效率。这一改进的动机来自~GS-Hider~的一个关键安全缺陷——耦合特征虽然从属性格式上看不出异常，但将~GS-Hider~的点云可视化后发现，隐藏场景的几何结构暴露在公开点云中（例如将"麦克风"隐藏在"盆景"场景中时，麦克风的形状在可视化点云中清晰可见）。其根本原因是~GS-Hider~用同一组高斯点同时表示原始场景和隐藏场景，不可避免地导致几何与空间位置的重叠，且缺乏对隐藏高斯点生长的优化控制。SecureGS（2025.03）~\cite{2503_06118} 受~Scaffold-GS~锚点设计启发，通过混合解耦高斯加密表示（HDGER）将隐藏高斯与原始高斯在表示层面完全分离，从根本上解决了这一问题。

\paragraph{混合解耦高斯加密表示}与~GS-Hider~用耦合特征~$f_i$~替代~SH~系数的"耦合-解耦"范式不同，SecureGS~采用~Scaffold-GS~的锚点-神经高斯架构作为基础，设计了一种显式-隐式混合表示。场景被体素化为锚点集合~$V \in \mathbb{R}^{N \times 3}$，每个锚点~$v \in V$~配备局部上下文特征~$f_v \in \mathbb{R}^{32}$、缩放因子~$l_v \in \mathbb{R}^3$~以及两组可学习偏移~$\{O_{v,i}^{ori}\}_{i=1}^k$~和~$\{O_{v,j}^{hid}\}_{j=1}^k$，分别生成渲染原始场景和隐藏对象的高斯点。关键设计在于偏移的存储方式：原始场景的高斯偏移~$\{O_{v,i}^{ori}\}_{i=1}^k$~显式存储在点云文件中（与~Scaffold-GS~格式完全一致），位置由~$\mu_{v,i}^{ori} = x_v + O_{v,i}^{ori} \cdot l_v$~计算；而隐藏对象的高斯偏移则通过隐式神经解码器~$F_o^*$~从锚点特征中动态推断——$\mu_{v,j}^{hid} = x_v + F_o^*(\hat{f}_v, \hat{d}_{vc}, \hat{n}_{vc})$，其中~$\hat{f}_v$~为多分辨率视角依赖混合特征，$\hat{d}_{vc} = \|x_v - x_c\|_2$~为相对距离，$\hat{n}_{vc}$~为视角方向。类似地，高斯属性也由公私两套~MLP~分别预测：公开~MLP~$\{F_c, F_\alpha, F_q, F_s\}$~生成原始场景的不透明度~$\alpha_{v,i}^{ori}$、旋转~$q_{v,i}^{ori}$、缩放~$s_{v,i}^{ori}$~和颜色~$c_{v,i}^{ori}$；私有~MLP~$\{F_c^*, F_\alpha^*, F_q^*, F_s^*\}$~生成隐藏对象的对应属性。偏移预测器~$F_o^*$~和私有~MLP~仅授权用户持有——从点云文件格式看，SecureGS~与标准~Scaffold-GS~完全相同，不包含任何可疑的额外属性。

\paragraph{密度区域感知优化}混合解耦表示虽然解决了属性格式的安全性，但锚点点云较为稀疏，隐藏对象的锚点仍可能在可视化时暴露几何线索。简单粗暴地增加所有区域的高斯密度虽能遮掩隐藏结构，但会显著增大存储和降低渲染速度。为此，SecureGS~提出区域感知密度优化（RDO）策略，自适应控制锚点的分裂与修剪，仅在隐藏对象所在区域大量生长锚点。该策略包含三个步骤：（1）异步梯度累积——对原始场景和隐藏对象分别以不同频率累积神经高斯点的反向传播梯度~$\tau_{ori}$~和~$\tau_{hid}$，通过减缓~$\tau_{hid}$~的累积速度抑制隐藏锚点的分裂，增强隐蔽性；（2）高密度区域定位——基于~$\tau_{hid} > \tau_{fix}$~（$\tau_{fix}$~为预设分裂阈值）的触发条件确定隐藏锚点集~$V_{hid}$，然后对其执行~DBSCAN~点云聚类，获取隐藏锚点密集分布的包围盒~$S_{bbx}$；（3）自适应生长与修剪——对原始场景锚点的生长采用自适应阈值替代固定阈值：$\tau_{ada} = \tau_{fix} / r_{down}$（当锚点位于~$S_{bbx}$~内时）或~$\tau_{fix}$（否则），其中~$r_{down}$~为梯度下采样比率。这使得原始场景的锚点仅在隐藏对象附近大量分裂以遮掩几何线索，其余区域保持稀疏以控制存储开销。

\paragraph{训练与损失函数}训练使用像素级~$\mathcal{L}_1$~损失、SSIM~项和体积正则化~$\mathcal{L}_{vol}$~联合优化锚点的可学习属性和所有~MLP~的权重。给定原始训练视角~$I_{gt}^{ori}$~和对应的隐藏视角~$I_{gt}^{hid}$，两个场景的损失分别为：

\begin{equation}
    \mathcal{L}_{ori} = (1-\lambda)\mathcal{L}_1(I_{ori}, I_{gt}^{ori}) + \lambda\mathcal{L}_{SSIM}(I_{ori}, I_{gt}^{ori}) + \mathcal{L}_{vol}(s^{ori}),
\end{equation}
\begin{equation}
    \mathcal{L}_{hid} = (1-\lambda)\mathcal{L}_1(I_{hid}, I_{gt}^{hid}) + \lambda\mathcal{L}_{SSIM}(I_{hid}, I_{gt}^{hid}) + \mathcal{L}_{vol}(s^{hid}),
\end{equation}

其中体积正则化~$\mathcal{L}_{vol}$~通过约束每个高斯缩放的乘积促使神经高斯紧凑、减少重叠。总损失为~$\mathcal{L}_{total} = \mathcal{L}_{ori} + \gamma \mathcal{L}_{hid}$，$\gamma$~平衡原始场景与隐藏对象的渲染质量。

\paragraph{安全性与扩展能力}SecureGS~从三个维度确保安全性：（1）点云文件格式安全——公开存储的内容与标准~Scaffold-GS~完全一致（锚点特征~$f_v$、原始偏移~$O_{v,i}^{ori}$~和缩放因子~$l_v$）；（2）点云结构安全——通过~RDO~策略，可视化锚点点云中几乎检测不到隐藏对象的几何痕迹；（3）渲染图像安全——渲染的原始场景不包含隐藏对象的任何伪影或边缘。除隐藏~3D~对象外，SecureGS~还支持位序列隐藏（通过额外的私有~MLP~$F_b^*$~将锚点特征解码为比特序列，所有体素解码结果取平均后交叉验证，实现~100\%~解码精度）和单张图像隐藏（在特定视角下约束渲染结果接近版权图像）。与~GS-Hider~和~3DGS+StegaNeRF~在~Mip-NeRF 360~的~9~个场景上对比，SecureGS~在渲染保真度、存储效率和速度方面均显著领先：原始场景~PSNR~提升~1.16~dB，隐藏对象~PSNR~提升~1.66~dB（甚至略超~Scaffold-GS~基线，证明解耦设计不干扰原始渲染），存储空间减少~201.24~MB，渲染速度提升近~$3\times$。鲁棒性方面，随机剪枝~5\%~锚点几乎不影响渲染质量，25\%~剪枝下原始场景和隐藏对象仍分别维持~24.96~dB~和~33.90~dB~的~PSNR。

\begin{takeaway}
隐写与信息隐藏利用 3DGS 的显式参数空间隐藏秘密 3D 场景或信息：解耦表示将原始与隐藏内容的渲染完全分离，隐式神经场解码实现容量和安全性兼顾，Scaffold-GS 锚点结构为大规模隐写提供了高效载体。
\end{takeaway}

%=====================================================================


\section{其他安全方法}
\label{sec:21_4}

除水印、对抗和隐写之外，3DGS~的安全研究还涉及篡改定位、图像保护和操控检测等新兴方向。

从水印、对抗和隐写三大核心方向过渡到其他安全方法，研究者的关注点扩展到更广泛的安全场景。

与锚点~GuardMarkGS~（在统一框架中联合优化水印和对抗以主动保护~3DGS~场景）不同，GS-Checker（2025.11）~\cite{2511_20354} 从防御视角转向取证视角——不再预防编辑行为，而是在事后检测~3DGS~场景是否已被恶意编辑并精确定位篡改区域。这一问题的紧迫性源于~3DGS~编辑工具（如~GaussianEditor、Gaussian Grouping）的快速发展：恶意用户可以利用这些工具篡改他人创建的~3DGS~模型的特定属性，不仅侵犯原始所有者的权益，还可能生成具有负面社会影响的篡改内容。传统~2D~篡改定位方法（如~SAFIRE、IML-ViT）虽然已相当成熟，但应用于~3DGS~时面临根本性局限：篡改操作发生在~3DGS~模型参数层面，~2D~方法只能从渲染图像中部分识别篡改痕迹，且不同视角的~2D~检测结果常常不一致。因为~2D~方法忽略了~3D~空间中跨视角的相关性，可能出现将真实区域误判为篡改区域的情况。更关键的是，传统~2D~方法中"篡改像素偏离正常统计分布"的现象同样延伸到~3D~领域：篡改后的~3D~高斯在参数分布上呈现异常，这些异常反映了篡改操作与原始重建过程之间的差异，为~3D~级别的取证分析提供了关键线索。GS-Checker~充分利用~3D~高斯参数本身的内在属性来识别篡改痕迹，通过三个协同组件实现~3DGS~场景的精确篡改定位。

\subsection{其他安全方法}
\label{ss:21_4}

本节讨论三条技术路线：（1）篡改检测与验证；（2）几何防护与版权保护；（3）3D伪造检测。

\subsubsection{篡改检测与验证}

本路线介绍通过 2D-3D 联合迭代精确定位 3DGS 场景中被篡改区域的方法。

\paragraph{篡改属性集成} GS-Checker~为每个~3D~高斯引入一个标量篡改属性~$r_i \in \mathbb{R}^1$，将其集成到~3DGS~模型参数中——对于第~$i$~个高斯，除原有的均值位置~$\mu_i$、不透明度~$\alpha_i$、缩放因子~$s_i$、旋转因子~$q_i$~和~SH~系数~$h_i$~外，新增的~$r_i$~指示该高斯是否被篡改。篡改属性的初始化采用~2D-3D~投影与多视角投票策略：首先将~3DGS~模型的渲染图像输入预训练的~2D~篡改定位网络（如~SAFIRE）获得每个视角的~2D~掩码~$M_j = F(R_j)$，其中~$R_j$~为第~$j$~个视角的渲染图像；然后将~2D~掩码投影回~3D~空间，对每个高斯的中心点进行多视角投票——若其投影落入篡改区域则投票为~$+1$，落入真实区域则为~$0$，落在掩码外则弃权~$-1$。最终通过累计投票数~$T_i = \sum_{j=0}^{N-1} V_{i,j}$~确定每个高斯的初始篡改属性值，其中获得多数篡改票的高斯（即多数视角认为其属于篡改区域）被标记为疑似篡改。这一投票机制的核心优势在于利用~3D~结构在不同视角间的一致性来消除~2D~模型的错误预测——由于~2D~篡改定位模型不考虑~3D~空间结构，其在不同视角可能给出不一致的预测结果，而多视角投票通过跨视角共识有效过滤这些不一致。

\paragraph{3D~对比机制} 篡改高斯与真实高斯在参数分布上存在系统性差异，GS-Checker~利用~3D~对比机制从~3D~空间中直接挖掘这些异常线索。具体地，根据篡改属性值将所有高斯分为三组：高值组~$G_h$（疑似篡改）、低值组~$G_l$（疑似真实）和中间组~$G_m$（待判定）。对比机制聚焦于~$G_m$~中高斯的归属判定——通过比较其属性与~$G_h$、$G_l$~的相似性来决定更新方向。相似性度量基于高斯的原始参数（位置~$\mu_i$、不透明度~$\alpha_i$、缩放~$s_i$、旋转~$q_i$~和~SH~系数~$h_i$），计算方式为~$\delta_i^h = \|F_i^m - F^h\|_2^2$~和~$\delta_i^l = \|F_i^m - F^l\|_2^2$，其中~$F^h$、$F^l$~分别为~$G_h$、$G_l$~中高斯属性的平均值。更新方向由差异的相对大小决定：$u_i = \mathrm{sgn}(\delta_i^l - \delta_i^h)$——若~$G_m$~中的高斯在属性上更接近~$G_h$（即~$\delta_i^h < \delta_i^l$），则增大其篡改属性值；反之则减小。直觉上，这一机制利用了篡改操作在参数分布中留下的统计足迹：篡改区域的编辑操作（如修改颜色、形状或插入/删除对象）会在参数空间中形成与原始重建过程不同的分布模式，对比机制正是通过比较这种分布差异来精确定位篡改区域。

\paragraph{循环优化策略} 循环优化策略通过~2D-3D~联合迭代实现篡改属性的逐步精化。由于篡改属性已被集成到~3DGS~模型参数中，可以利用可微光栅化将其渲染到~2D~：$M^R(r) = \sum_{i=1}^{G} r_i\,\alpha_i\,\prod_{j=1}^{i-1}(1-\alpha_j)$，其中~$G$~为光线上重叠的高斯数量。渲染结果与~2D~篡改定位模型生成的掩码之间的差异构成~2D~渲染损失：$\mathcal{L}_{2d} = -\alpha_1 \sum_{r \in \mathcal{R}(I_j)} M_j(r)\,\log M^R(r) + \alpha_2 \sum_{r \in \mathcal{R}(I_j)} (1-M_j(r))\,\log M^R(r)$，其中~$\alpha_1 = 1.0$、$\alpha_2 = 10.0$~为平衡超参数——对真实区域的惩罚权重（$\alpha_2$）远高于篡改区域（$\alpha_1$），确保真实区域的篡改属性被强烈压低。优化过程采用两阶段策略：第一阶段仅使用~2D~渲染损失优化篡改属性，经过若干迭代后~$G_h$~与~$G_l$~之间的篡改属性值差异拉大，便于后续对比机制区分；第二阶段联合~2D~渲染损失和~3D~对比损失进行优化，总损失为~$\mathcal{L}_{\mathrm{cyc}} = \beta_1\,\mathcal{L}_{2d} + \beta_2\,\mathcal{L}_{3d}$（$\beta_1 = 1.0$、$\beta_2 = 10.0$），其中~3D~对比损失~$\mathcal{L}_{3d} = -\frac{1}{K}\sum_{i}^{K} u_i\,r_i$~（$K$~为~$G_m$~中高斯数量）驱动中间组高斯的篡改属性向正确方向更新。整个优化过程中仅更新篡改属性，不影响~3DGS~模型的其他参数。使用~Adam~优化器，学习率为~1.0，归一化后的篡改属性阈值为~0.1，所有实验在单块~NVIDIA Tesla V100 GPU~上完成。

\paragraph{实验评估} GS-Checker~在自建的~3DGS~篡改数据集上进行评估，该数据集包含~11~个被篡改的~3DGS~模型，覆盖三种篡改类型：对象插入（在~3D~场景中添加新对象）、对象修改（编辑~3D~场景中对象的颜色、形状等属性）和对象删除（移除~3D~场景中的对象并填补空隙）。篡改操作使用~GaussianEditor~和~Gaussian Grouping~两种编辑方法，基于~Mip-NeRF 360~和~Instruct-NeRF2NeRF~数据集训练的~3DGS~模型。由于此前没有专门针对~3DGS~篡改定位的方法，基线比较采用三种策略组合：（1）SAFIRE + SA3D（2D~篡改定位 + 3DGS~分割）；（2）SAFIRE + SAGD（2D~篡改定位 + 另一种~3DGS~分割）；（3）Inverse with SAFIRE（将~2D~掩码直接投影回~3D~高斯）。实验表明，GS-Checker~在三种篡改类型上均取得最佳定位精度（以~$F_1$~分数和~IoU~衡量），证明了~3D~对比机制和循环优化在~3DGS~篡改定位中的有效性。鲁棒性测试表明，GS-Checker~在~2D~高斯噪声、高斯模糊、3D~缩放噪声和不透明度噪声等失真条件下仍能维持较为准确的定位性能。消融实验验证了每个组件的贡献：单独移除~3D~对比机制或循环优化策略均导致定位性能下降，两者的组合才能达到最佳效果。此外，在真实场景（未被篡改）上的测试表明，~2D~篡改定位模型会错误地将部分区域预测为篡改，而~GS-Checker~的~3D~投票机制能够有效滤除这些误报——这进一步证实了从~2D~向~3D~提升的必要性。泛化性测试还验证了~GS-Checker~对其他~3DGS~编辑方法（如~DGE）和不同~2D~篡改定位骨干网络（如~IML-ViT）的适应性。

\subsubsection{几何防护与版权保护}

本路线介绍在图像发布前嵌入几何防护以阻止未经授权~3D~重建的预防性方法。篡改定位解决的是事后取证问题，而~Geometry Cloak~则从预防的角度出发——在图像发布之前就嵌入防护措施。如果说~GS-Checker~回答的是"3DGS~场景是否被篡改"，那么~Geometry Cloak~回答的是"如何在图像层面阻止未经授权的~3D~重建"。两者构成事后取证与事前预防的完整安全闭环。

\paragraph{问题定义与动机} Geometry Cloak（2024.10）~\cite{2410_22705} 针对~TGS~单视角~3D~重建管线提出图像保护方法。TGS~采用显式点云与隐式~triplane~的混合~3D~表示：点云~$P \in \mathbb{R}^{N \times 3}$~提供显式几何，triplane~$T = (T_{xy}, T_{xz}, T_{yz}) \in \mathbb{R}^{3 \times C \times H \times W}$~编码隐式特征场；对于点云中的位置~$x \in \mathbb{R}^3$，通过对三个正交平面的三线性插值与拼接获得~triplane~特征~$f_t$，再由~MLP~解码为~3D~高斯属性（不透明度、协方差、SH~系数）。Geometry Cloak~的核心发现是：在~TGS~的重建管线中，图像特征对扰动具有显著鲁棒性——即使施加明显的对抗扰动，重建的~3D~模型也几乎不受影响；而点云解码器则存在固有脆弱性——仅需极小的不可见扰动即可大幅改变重建的~3D~模型。这一脆弱性的根源在于：点云是场景结构的蒸馏和直接表示，对其施加扰动在~3D~输出中极为显著；同时，点云作为后续~triplane~特征采样的基础几何信息，其扰动通过~$f_t = (T_{xy}, P_{xy}) \oplus (T_{xz}, P_{xz}) \oplus (T_{yz}, P_{yz})$~的插值过程级联传播至整个重建结果。因此，Geometry Cloak~选择绕过图像级特征的固有鲁棒性，直接攻击点云解码器的显式几何输出。

\paragraph{视角特异性~PGD} Geometry Cloak~通过视角特异性~PGD~（Projected Gradient Descent）优化扰动~$\delta$，使~TGS~从受保护图像~$I + \delta$~重建的点云在指定视角~$\theta$~下的投影逼近预设的~2D~点云水印图案。总体优化目标为~$\delta^* = \arg\min_{\|\delta\|_\infty \leq \epsilon}\, L_{CD}\!\bigl(E(I+\delta),\, P_w\bigr)$，其中~$E$~表示~TGS~中将图像映射为点云的网络，$\epsilon$~为扰动预算，$L_{CD}$~为~Chamfer~距离，$P_w$~为目标水印点云。具体地，视角特异性~PGD~的迭代更新公式为~$I^{(i+1)} = I^{(i)} + \alpha\,\mathrm{sign}\!\bigl(\nabla_I L_{CD}(P_\theta^{(i)}, P_w)\bigr)$，其中~$P_\theta^{(i)} = \mathrm{Proj}_\theta\!\bigl(E(I^{(i)} + \delta^{(i)})\bigr)$~表示从受保护图像重建的点云在视角~$\theta$~下的~2D~投影，$\alpha$~为步长。实验中采用~mask~版本~PGD——仅对图像中的前景物体区域施加扰动，因为~TGS~仅利用前景进行~3D~重建。优化过程迭代~$N=100$~步，步长~$\alpha = 0.001$，更新后的图像裁剪至有效范围~$[0,1]$。对于自定义~3D~图案（而非~2D~预定义图案），则直接计算重建点云与目标~3D~点云之间的~Chamfer~距离，无需视角投影。

\paragraph{可验证图案} Geometry Cloak~支持两种可验证图案。（1）预定义图案：将字母数字字符（A--Z、1--9）的图像进行分割，记录分割区域的坐标，从分割坐标中采样点形成~2D~点云作为水印图案。在~GSO~数据集（1030~个物体，每个随机选取一个~2D~水印字符）和~OmniObject3D~数据集（190~类、950~张图像）上的实验表明，嵌入预定义图案后，TGS~重建的~3D~模型在视觉和几何质量上均被严重破坏（PSNR/SSIM~大幅下降，CD~距离显著增大），且从特定视角可重新浮现嵌入的水印信息——这既保证了~3D~模型的不可用性，又保留了可验证的版权标记。（2）自定义图案：用户可选择性保护图像的特定区域。首先采用与~TGS~相同的方法从图像中提取点云，然后通过~InstructP2P~等文本引导编辑技术或~MeshLab~等开源工具对点云进行精细裁剪，选择需要保护的区域（如物体的尾部、身体、头部）生成定制~3D~点云作为目标图案。实验表明，用户可通过自定义图案有效地影响重建结果的特定部分，进一步证实了~TGS~点云对扰动的敏感性。

\paragraph{泛化性与鲁棒性} Geometry Cloak~的防护效果不仅限于~TGS~本身。由于所有基于~3DGS~的单视角~3D~重建方法均依赖显式几何特征来表示~3D~场景，Geometry Cloak~可自然扩展到其他~GS-based~方法。论文在~LGM~（另一种~GS-based~单视角~3D~重建方法）上验证了这一泛化性：在~LGM~的三个默认场景中，Geometry Cloak~施加的扰动有效地降低了重建~3D~模型的质量，且效果远优于同等强度的~Gaussian~噪声。此外，论文还测试了~triplane~扰动与点云扰动的组合，结果表明组合扰动并未提升攻击性能——这进一步证实了~triplane~特征是~TGS~中鲁棒的组成部分，难以被有效扰动。鲁棒性实验还表明，Geometry Cloak~对~JPEG~压缩等常见图像操作具有一定抵抗力，且在~$\epsilon = 0$~到~$4$~的扰动强度范围内，重建质量呈明显下降趋势。

\paragraph{安全意义} Geometry Cloak~将~3DGS~安全从"事后检测"范式拓展到"事前预防"范式。这使得版权持有者可在发布~2D~图像前嵌入几何隐身衣，对任何从受保护图像生成的~3D~重建主张所有权——无论谁使用何种~TGS~模型进行重建，输出都将携带预设的版权标记或遭受不可逆的质量破坏。

从图像保护到~3D~伪造检测，安全研究从"保护"视角扩展到"鉴别真伪"的取证视角。

\subsubsection{3D伪造检测}

本路线介绍检测~3DGS~场景表示是否被编辑操控以及防御不安全内容生成的安全方法。篡改定位和图像保护是两种互补的安全策略，而随着~3DGS~操控技术日趋成熟，一个此前被忽视的安全盲区逐渐显现：当编辑操作直接发生在~3D~场景表示层面时，所产生的渲染图像在~2D~层面几乎不留任何可辨识的伪影，使得传统~2D~假图像检测方法面临根本性失效。Fake3DGS（2026.04）~\cite{2604_27590} 正式将~3D~假检测定义为一项独立任务——判定~3DGS~场景表示是否在其几何、空间布局或外观参数上被修改，而不论其渲染图像看起来多么逼真。与~GS-Checker~（检测~3DGS~参数篡改并精确定位篡改区域）不同，Fake3DGS~的核心贡献在于构建首个大规模~3DGS~场景操控检测基准数据集，并系统证明当前最先进的~2D~深度假检测器在~3D~操控场景中严重失效。

\paragraph{3D~操控的独特威胁} 传统~2D~假图像检测基于一个核心假设：伪造内容在像素或纹理层面会留下可辨识的伪影（如频谱异常、几何不一致或光照不合理）。然而，~3DGS~操控从根本上瓦解了这一假设——编辑操作直接作用于~3D~高斯参数（位置、颜色、几何形状等），所产生的~3D~场景在渲染时天然具备多视角几何一致性，生成的假图像从任何角度看都"完美无瑕"。这意味着传统~2D~检测器赖以区分真伪的统计足迹（如~GAN（Generative Adversarial Network，生成对抗网络）~生成器特有的频谱特征、扩散模型特有的重建残差）在~3D~操控图像中几乎不存在——攻击者不再需要掩饰~2D~伪影，因为~3D~一致性本身就是"伪造"的一部分。更严峻的是，~3DGS~编辑能力正被加速部署到增强现实/虚拟现实、数字孪生和自动驾驶仿真等对真实性有严格要求的场景中，使得~3D~假内容的检测不仅是科学问题，更是安全问题。

\paragraph{基准数据集构建} Fake3DGS~数据集包含超过~41,000~个重建的~3DGS~场景，其中真实样本与操控样本各占一半，是迄今为止规模最大的~3DGS~操控检测基准。真实场景基于~UCO3D~数据集——涵盖超过~1,000~个物体类别、多样化的室内外环境，并提供预训练的~3DGS~重建结果。为避免类别偏差，数据集在每个类别上均衡采样。操控场景采用两种~3DGS~编辑方法生成：（1）GaussCtrl——通过基于~ControlNet~的深度条件编辑和基于注意力的潜码对齐实现~3D~场景编辑，保持几何一致性；（2）Instruct-GS2GS——利用~InstructPix2Pix~迭代编辑输入图像并同时优化底层~3D~场景。编辑指令由~Meta-Llama-3-8B-Instruct~自动生成，覆盖三类操控族：背景/表面变化（如草地变为雪地）、物体外观变化（如颜色修改）和物体类型/材质替换（如羊变为长颈鹿），确保每类编辑在数据集中均匀分布。所有~3DGS~模型导出为~nerfstudio~格式，并采用~Self-Organizing Gaussian Grids (SOGS)~压缩策略进行后处理——通过~8-bit~均匀标量量化对高斯参数独立量化，再将量化张量重排为局部连贯的~2D~网格并以~PNG~无损压缩编码，在忠实保留重建质量的同时将数据集占用从约~7TB~压缩至约~200GB，使大规模训练和评估在标准研究硬件上可行。数据集提供两种划分策略：混合划分（80\%--20\%训练-测试分割，约~33k~训练场景和~8k~测试场景）评估已知编辑器下的整体检测性能；交叉编辑划分（用一种编辑方法训练、另一种测试）评估对未知编辑器的泛化能力。

\paragraph{2D~检测器的系统性失效} 实验评测涵盖六类~2D~深度假检测基线：基于~ResNet-50~的~CNN~检测器（Wang et al.）、基于扩散重建残差的~DIRE、基于对比学习的~CoDE~嵌入、基于~CLIP~的线性分类器（ViT-B/16~和~ViT-L/14~两种变体）以及基于~DINOv2~自监督骨干的线性分类头，所有基线均在训练集上微调以确保公平比较。在混合设置下，2D~检测器表现尚可——最强基线~CoDE~达到~92.2\%~的总体准确率，多数方法能维持较高水平。然而，在交叉编辑泛化设置下，所有~2D~方法出现显著性能崩溃：CLIP~和~DINO~变体虽然能保持~93--98\%~的真实样本准确率，但假样本准确率骤降至~41.5--45.9\%——检测器倾向于将大量操控场景误判为真实。即使是交叉编辑设置下最强的~2D~基线（在~GaussCtrl~上训练、在~Instruct-GS2GS~上测试），其~83.8\%~的总体准确率仍由真实样本~95.8\%~的高准确率驱动，假样本准确率仅为~74.8\%。这些结果清晰揭示了一个关键发现：2D~检测器依赖的是编辑器特有的伪影线索，这些线索无法跨编辑方法迁移；而~3DGS~操控生成的渲染图像本身在~2D~层面几乎不留通用的可检测痕迹。

\paragraph{3D~感知检测方法} 针对上述问题，Fake3DGS~提出直接在~3D~高斯表示上操作的检测器~Fake3DD。该方法以~PointTransformerV3~为骨干——一种基于~Transformer~的点云架构，通过局部自注意力机制为每个点聚合空间邻域信息。Fake3DD~将~3DGS~场景自然地视为~3D~空间中分布的基元集合，对原始点云输入进行扩展以容纳高斯特有属性：每个高斯以其均值点的~3D~坐标作为位置输入，以不透明度、缩放因子、四元数旋转、球谐函数~$s_h$~及其零阶项~$s_0$（视角无关颜色分量）作为特征输入。编码器-解码器阶段提取上下文感知的高斯级特征后，通过全局均值池化将可变数量的高斯特征聚合为固定维度的场景级嵌入，再经分类头输出二元预测（假/真实）。这一设计充分利用~3DGS~的表示特性：真伪线索可能取决于高斯之间的空间上下文关系，而非单个高斯的孤立属性。

\paragraph{实验结果与消融分析} 在交叉编辑泛化设置下，Fake3DD~在~GaussCtrl~训练、Instruct-GS2GS~测试的协议上达到~98.7\%~的总体准确率，且类别间高度均衡（假样本~99.1\%、真实样本~98.4\%），相比最强~2D~基线提升~14.9~个百分点——其中假样本准确率提升~24.3~个百分点（$74.8\% \to 99.1\%$）。这一差距直接证明：直接操作~3DGS~表示能获取编辑器无关的真伪证据，而~2D~渲染图像中缺乏这种可迁移的判别信号。消融实验逐一移除高斯属性以评估各特征的贡献：移除不透明度导致准确率下降最大（$92.5\%$），表明透明度信息对检测最为关键——编辑操作会改变高斯的不透明度分布模式；缩放因子和零阶球谐系数~$s_0$~提供中等贡献，分别捕获受编辑影响的空间分布属性和视角无关颜色信息；四元数旋转和高阶球谐系数的影响较小，说明高斯方向和视角相关外观在真实与操控场景间保持相对一致，判别信息有限。这些发现为~3DGS~安全研究提供了明确的启示：操控痕迹最显著地编码在高斯的透明度和几何属性中，而非外观细节中。

如果说~Fake3DGS~关注的是操控检测（事后取证），那么~3DEditSafe（2026.05）~\cite{2605_15398} 则从预防的角度出发——首次系统研究并防御文本驱动~3D~编辑管线中的不安全内容生成问题。3DEditSafe~的出发点是一个此前被忽视的安全盲区：在以~EditSplat~为代表的~3D~编辑管线中，不安全语义不仅存在于单张~2D~编辑图像，还会通过多视角优化被传播并固化到最终的~3D~高斯场景中，使得不安全内容在新视角下持续可见。这意味着，即使~2D~扩散生成阶段已对不安全概念进行了一定抑制，残留的不安全信号仍可在~3D~优化过程中被跨视角对齐和强化。3DEditSafe~首先从攻击面分析入手，构建了包含~30~个提示-场景对的评测基准（基于~Unsafe Diffusion~提示集和~10~个~EditSplat~场景，涵盖暴力、血腥和恐怖等语义类别），验证了原始~EditSplat~管线在不安全提示下的脆弱性——平均视角级攻击成功率（ASR，Attack Success Rate）为~$0.347$，场景级攻击成功率高达~$0.722$。

在方法论层面，3DEditSafe~采用三阶段安全防线。第一阶段是提示风险门控（PRG，Prompt Risk Gate）：计算编辑提示与预定义不安全概念集~$C_\text{unsafe}=\{c_1, c_2, \ldots, c_K\}$~的语义相似度和关键词匹配得分的加权和，输出风险分数~$s_\text{risk}(p)$。仅当~$\mathbb{1}_\text{risk}(p)=1$~时才激活后续安全机制，良性提示直接走标准编辑路径——这与提示过滤或提示清洗方法（在生成前拒绝或改写用户指令）不同，3DEditSafe~始终保持原始编辑指令不变，安全约束仅作用于~3D~优化过程。第二阶段是语义安全投影（SSP，Semantic Safety Projection）：给定编辑提示的~CLIP~文本嵌入~$E_\text{CLIP}(p)$，找到与其最对齐的不安全概念方向~$\mathbf{u}$，然后通过~$E_\text{safe}(p) = \text{normalize}\big(E_\text{CLIP}(p) - \beta\, \langle E_\text{CLIP}(p), \mathbf{u}\rangle\,\mathbf{u}\big)$~投影到最近的安全语义方向。这种"软投影"策略避免了将所有不安全提示映射到同一安全模板的退化问题，为~3D~优化器提供了一个接近原始编辑意图的正向语义目标。第三阶段是渲染视角~3D~安全正则化（RV3SR，Rendered-View 3D Safety Regularization），也是方法的核心贡献——直接在~3D~优化过程中对渲染视角施加语义安全约束，而非仅在~2D~去噪阶段防御。具体包含三个正则项：不安全概念排斥损失~$\mathcal{L}_\text{unsafe}^\text{3D} = \sum_i \max\big(0,\, \max_{c_j \in C_\text{unsafe}} \cos\big(F_\text{img}(R(G, v_i)),\, E_\text{CLIP}(c_j)\big) - m\big)$~惩罚渲染视角与不安全概念的语义对齐（$m$~为安全边界）；安全语义引导损失~$\mathcal{L}_\text{safe}^\text{3D} = 1 - \cos\big(F_\text{img}(R(G, v_i)),\, E_\text{safe}(p)\big)$~引导渲染视角朝向投影后的安全语义方向；以及弱保持损失~$\mathcal{L}_\text{preserve} = \|M_i \odot (R(G, v_i) - I_i)\|_1$~稳定编辑区域内的视觉保真度（$M_i$~为~LangSAM~预测的目标掩码）。最终优化目标为~$\mathcal{L}_\text{total} = \mathcal{L}_\text{edit} + \lambda_o \mathcal{L}_\text{outside} + \mathbb{1}_\text{risk}(p)\big(\lambda_u \mathcal{L}_\text{unsafe}^\text{3D} + \lambda_s \mathcal{L}_\text{safe}^\text{3D} + \lambda_p \mathcal{L}_\text{preserve}\big)$，其中风险门控~$\mathbb{1}_\text{risk}(p)$~确保安全正则项仅在检测到风险提示时激活。

实验结果揭示了几个关键发现。首先，~2D~安全引导（仅作用于扩散生成阶段的~Safe Latent Diffusion~式安全引导）虽在某些场景有帮助，但效果不一致——例如在~Bear~场景上将视角级~ASR~从~$0.236$~降至~$0.122$，但场景级~ASR~仍为~$1.0$；在~Dinosaur~场景上~ASR~反而从~$0.586$~上升至~$0.649$。这直接证实了~2D~层面安全约束的局限性。相比之下，3DEditSafe~将平均视角级~ASR~从~$0.347$~降至~$0.190$，场景级~ASR~从~$0.722$~降至~$0.389$，相比~2D~安全引导基线分别进一步降低~$46\%$~和~$49\%$。安全性与质量的权衡同样清晰：在不安全提示上~3DEditSafe~的~CLIP~目标相似度从~EditSplat~的~$0.266$~降至~$0.233$，但考虑到目标提示本身包含不安全语义，这一下降正是安全抑制的预期效果。保持偏差保持在低水平（$0.024$）。这些结果表明，在~3D~编辑管线中，安全性控制必须在~3D~优化阶段本身实施，仅依赖~2D~扩散阶段的防御是不充分的。3DEditSafe~与~Fake3DGS~（操控检测）和~AdLift~（扩散编辑防御）形成互补——前者侧重事后取证，后者聚焦~2D~扩散过程，而~3DEditSafe~直接约束~3D~场景的语义演化。

\begin{takeaway}
其他安全方法从篡改检测扩展到操控防御：高斯属性的不透明度和几何信息是操控检测的最关键特征，而 3D 优化阶段本身的安全性控制是不可或缺的——仅依赖 2D 扩散阶段的防御在跨视角传播中会失效。
\end{takeaway}

%% ============================================================

\section{近期研究（2026 年 8--11 月）}
\paragraph{NGS-Marker：从局部高斯子集解码资产水印}
2026.08~\cite{2608_17447}。图像水印往往依赖特定渲染结果，资产只被取走一部分时未必仍可解码。NGS-Marker 学习高斯空间的消息注入与提取，并以预训练解码器监督随机局部片段的渐进嵌入，使被截取的高斯也携带可恢复的身份信息。渲染误差和消息恢复共同约束训练。

作者用 24 个训练场景和四个留出场景，在局部子集被提取并插入另一场景的设置下比较位准确率与原语级识别。方法需要访问高斯原语，超大场景高效分区与遍历仍是作者指出的扩展问题。

\paragraph{TRACE：用射线汇总量做跨私有子地图选视}
2026.10~\cite{2610_00822}。多机器人在各自持有高斯子地图时，某条候选射线的遮挡仍受其他机器人地图影响。TRACE 按深度区间交换光学厚度、辐射与位姿导数的汇总量，用来近似合并地图上的选视信息增益，无需发送整个原语集合。

在 Habitat-Sim 的 Gibson 室内场景，作者固定建图和规划管线，只比较选视策略。不同机器人命中落在同一深度区间时会产生近似误差，稀疏通信实现才体现显著带宽优势。该工作限制了直接共享的内容，仿真也未给出统计意义上的抗信息泄露保证。

\paragraph{FlyMark：在静态高斯的基本颜色内写入模型域水印}
2026.10~\cite{2610_11364}。FlyMark 将隐式所有权消息写入已经训练好的静态高斯资产，仅在色差受限的条件下调整零阶颜色。读出依赖合法密钥和对模型参数的访问；无需改变位置、尺度或训练额外解码器。水印所作用的是模型本身，并非单张渲染图。

原文按场景和密钥组合分别考察正常读出、错密钥和模型变化后的结果，并用固定相机进行画质对照。适用条件包括可读取静态参数数组以及作者定义的模型域变换。所测结果不能解释为对任意再训练、蒸馏或模型重建的普遍鲁棒性。

\paragraph{PrinGS：在高斯表示层处理人脸和文字隐私}
2026.10~\cite{2610_36168}。PrinGS 根据敏感区域的图像掩码和光栅化时的可见性定位相关高斯，再对其局部外观或位置进行匿名化。它既可在现有场景上操作，也可把匿名化约束加入训练；现有场景模式可用渲染视图上的掩码，因此可以在原始训练图像缺失时使用。

作者在脸部和文字实验中分别检查检测率、再识别相似度或 OCR 字符错误率。效果受掩码覆盖、可见性估计以及是否采到了关键视角影响，留出视图上的隐私指标不能代替证明所有视角都无法恢复信息。


\section*{本章小结}
\addcontentsline{toc}{section}{本章小结}

\begin{futurebox}{本章总结与展望}{ch21_summary}
本章从安全、水印与隐私角度综述了~3D Gaussian Splatting~的保护技术，涵盖以下核心内容：
\begin{itemize}
\item \textbf{数字水印}（\ref{sec:21_1}节）：从基础~3DGS~水印框架（3D-GSW、WATER-GS、GuardSplat）到高容量水印（BitC-3DGS）、可泛化水印（GS-Marker）、生成模型水印（GaussianSeal）和压缩鲁棒水印（CompMarkGS），水印技术正在从概念验证走向实用化。安全分析揭示了水印框架的脆弱性——GSPure~等攻击方法可有效移除水印，推动了对更强鲁棒性的需求。
\item \textbf{对抗攻击与防御}（\ref{sec:21_2}节）：攻击方法从重建质量攻击（Poison-splat、PatchPoison）到下游检测器攻击（CLOAK、3DGAA、AdvSplat），揭示了~3DGS~在多个层面的脆弱性。防御方法（DefenseSplat、Spectral Defense、RemedyGS）从频率分析和检测-净化两个方向提升系统鲁棒性。
\item \textbf{隐写与信息隐藏}（\ref{sec:21_3}节）：从~GS-Hider~到~Splats in Splats++，隐写框架实现了在~3DGS~资产中嵌入完整~3D~场景（而非仅版权标识），重要性分级~SH~加密和哈希网格引导不透明度映射是核心技术，支持~2D~图像、~3D~场景和~4D~动态场景的多模态隐写。
\item \textbf{其他安全方法}（\ref{sec:21_4}节）：篡改定位（GS-Checker）和图像保护（Geometry Cloak）拓展了~3DGS~安全的研究边界。
\end{itemize}

当前~3DGS~安全研究面临的主要挑战包括：
\begin{itemize}
\item \textbf{攻防军备竞赛}：水印攻击（GSPure、GMEA）和水印防御（RDSplat、GuardMarkGS）的对抗升级仍在持续，缺乏理论上可证明安全的水印方案。
\item \textbf{物理可实现性}：大多数对抗攻击在数字域验证，物理世界的部署面临打印失真、视角变化和天气干扰等挑战。R-PGA~和~PGA~等方法的物理实验提供了初步验证，但距离可靠的物理攻击评估仍有差距。
\item \textbf{标准化评估}：缺乏统一的~3DGS~安全评估基准和威胁模型定义，Deng et al.（2026.02）~\cite{2602_02602} 的场景驱动框架是一个重要起点但尚未被广泛采用。
\item \textbf{4DGS~安全}：动态~3DGS~场景的水印和隐写需要处理时序一致性（4D-GSW~的初步探索），对抗~4DGS~的攻击和防御研究几乎是空白。
\end{itemize}
\end{futurebox}
